% Leçon 35 — Grammaire : 不管是…….或者是……. ; 是。。。。。。的 ; 比字句 — Chinois Tle A — Cours signé OnBuch+
% Programme officiel MINESEC. Compiler avec : tectonic ZH35-grammaire.tex
\newcommand{\DOCMATIERE}{Chinois}
\newcommand{\DOCNIVEAU}{Tle A}
\newcommand{\DOCMODULE}{Module 5 : Mass médias et télécommunication}
\newcommand{\DOCLECON}{ZH35}
\newcommand{\DOCTITRE}{Leçon 35 — Grammaire : 不管是…….或者是……. ; 是。。。。。。的 ; 比字句}
% OnBuch+ — préambule « cours signés » ENRICHI (compilé avec tectonic / XeLaTeX)
% Système visuel « Pop » d'OnBuch+ : fond crème (jamais blanc), bordures encre
% épaisses, ombres « dures » décalées, titres Archivo Black en capitales.
% Version MATHÉMATIQUES : graphes (pgfplots/tikz), boîtes pédagogiques
% (définition, propriété, méthode, attention, le savais-tu, exercice résolu,
% exercice, TP, bilan), page de garde et signature OnBuch+ ancrée en bas.
%
% Macros attendues dans main.tex AVANT \input{preamble} :
%   \DOCMATIERE  \DOCNIVEAU  \DOCMODULE  \DOCLECON (numéro)  \DOCTITRE
\documentclass[11pt]{article}

\usepackage{fontspec}
\usepackage[french]{babel} % français = langue principale ; le chinois via \zh{..} / textebox (xeCJK)
\usepackage{xeCJK}
\usepackage{pifont} % \ding{51} \ding{55} utilisés par les modèles
\usepackage[a4paper,margin=2cm,top=3.4cm,bottom=2.8cm,headheight=1.4cm,footskip=1.3cm]{geometry}
\usepackage{amsmath,amssymb,amsfonts}
\usepackage{mathtools} % \xmapsto, \coloneqq... souvent utilisés par les modèles
\usepackage{cancel} % simplifications barrées (bilans de forces)
% \abs{z} (module, valeur absolue) et \norm{u} (norme), utilisés par les modèles
\providecommand{\abs}{}\let\abs\relax\DeclarePairedDelimiter{\abs}{\lvert}{\rvert}
\providecommand{\norm}{}\let\norm\relax\DeclarePairedDelimiter{\norm}{\lVert}{\rVert}
\usepackage[table]{xcolor} % [table] : fournit aussi \rowcolors (alternance de lignes)
\usepackage{pagecolor}
\usepackage{tikz}
\usetikzlibrary{arrows.meta,positioning,calc,shapes.geometric,shapes.misc,shapes.arrows,decorations.pathmorphing,decorations.pathreplacing,decorations.markings,patterns,shadows,mindmap,trees,fit,backgrounds,matrix,intersections,angles,quotes}
\usepackage{pgfplots}
\usepackage[version=4]{mhchem} % équations nucléaires \ce{^{235}_{92}U}
\usepackage[european,siunitx]{circuitikz} % schémas électriques
\pgfplotsset{compat=1.18}
\usepgfplotslibrary{fillbetween,groupplots} % aires entre courbes (calcul intégral)
\usepackage{enumitem}
\usepackage{fancyhdr}
% [most] charge toutes les bibliothèques tcolorbox (listings, minted, raster,
% documentation...) et gonfle inutilement la mémoire TeX sur un document long
% (~300 boîtes) : on ne charge que ce qui est réellement utilisé.
\usepackage[skins,breakable]{tcolorbox}
\usepackage{graphicx}
\usepackage{colortbl}
\usepackage{booktabs}
\usepackage{tabularx}
\usepackage{diagbox} % case d'en-tête diagonale des tableaux à double entrée
% Colonnes extensibles centrée / gauche / droite (C, L, R), souvent utilisées par les modèles
\newcolumntype{C}{>{\centering\arraybackslash}X}
\newcolumntype{L}{>{\raggedright\arraybackslash}X}
\newcolumntype{R}{>{\raggedleft\arraybackslash}X}
\usepackage{array}
\usepackage{multicol}
\usepackage{multirow}
\usepackage{siunitx}
% parse-numbers=false : accepte aussi \SI{2,5 \times 10^{-3}}{...}
\sisetup{locale=FR,output-decimal-marker={,},group-minimum-digits=4,parse-numbers=false}
\DeclareSIUnit\ohm{\ensuremath{\symup{\Omega}}} % U+2126 absent de la police
\DeclareSIUnit\division{div} % réglages d'oscilloscope (ms/div, V/div)
\DeclareSIUnit\tr{tr} % tours (vitesse de rotation, ex. \SI{1500}{\tr\per\minute})
\DeclareSIUnit\molar{M} % concentration molaire (ex. \SI{3}{\molar}, \SI{1}{\micro\molar})
\DeclareSIUnit\an{an} % année (le modèle confond parfois avec \year, un compteur TeX) — ex. \SI{5}{\kilo\meter\per\an}
\DeclareSIUnit\FCFA{FCFA} % franc CFA (ex. \SI{500}{\FCFA\per\kilo\gram})
\DeclareSIUnit\degreeBrix{\textdegree Brix} % teneur en sucre d'un sirop/jus (réfractométrie)
\DeclareSIUnit\month{mois} % le modèle utilise parfois \month, un compteur TeX, comme unité de durée
\usepackage{qrcode}
% \degree : commande fréquemment utilisée par les modèles pour un angle ou une
% température en dehors de \SI{}{} (non fournie par les paquets chargés).
\providecommand{\degree}{^{\circ}}
% \qed (fin de démonstration), utilisé par les modèles sans amsthm
\providecommand{\qed}{\hfill$\square$}
% \barre : double barre des valeurs interdites dans un tableau de variations (tkz-tab)
\providecommand{\barre}{\ensuremath{\|}}
% environnement proof (amsthm non chargé) utilisé par les modèles
\ifdefined\proof\else\newenvironment{proof}[1][Démonstration]{\par\noindent\textit{#1.}\ }{\qed\par}\fi
\usepackage{lastpage}
\usepackage{hyperref}
\hypersetup{hidelinks}

% ============================================================
% Palette « Pop » OnBuch+ — mêmes hex que la classe P (plus_ui.dart)
% ============================================================
\definecolor{poporange}{HTML}{F59321}
\definecolor{popdark}{HTML}{E07A0C}
\definecolor{popcream}{HTML}{FFF6E8}
\definecolor{popink}{HTML}{1C1714}
\definecolor{poppurple}{HTML}{7A5AE0}
\definecolor{popgreen}{HTML}{2E9E6A}
\definecolor{poppink}{HTML}{E0457B}
\definecolor{popblue}{HTML}{2F7DE1}
\definecolor{popgold}{HTML}{C98A12}
\definecolor{popmuted}{HTML}{8A7F76}
\definecolor{popline}{HTML}{EADFD2}
\definecolor{popgray}{HTML}{8A7F76}
\colorlet{popgrey}{popgray}
% Alias tolérants (le modèle nomme parfois les couleurs différemment)
\colorlet{popwhite}{white}
\colorlet{popblack}{popink}
\colorlet{popred}{poppink}
\colorlet{popyellow}{popgold}
% Teintes claires (fonds de tableaux/encadrés) — le modèle nomme parfois
% les couleurs avec un suffixe L (ex. popblueL) sans les avoir définies.
\colorlet{poporangeL}{poporange!20}
\colorlet{popdarkL}{popdark!20}
\colorlet{poppurpleL}{poppurple!20}
\colorlet{popgreenL}{popgreen!20}
\colorlet{poppinkL}{poppink!20}
\colorlet{popblueL}{popblue!20}
\colorlet{popgoldL}{popgold!20}
\colorlet{popredL}{popred!20}
\colorlet{popgrayL}{popgray!20}
% Teintes claires pour fonds de tableaux / zones
\colorlet{popblueL}{popblue!10!white}
\colorlet{poporangeL}{poporange!14!white}
\colorlet{popgreenL}{popgreen!12!white}
\colorlet{poppurpleL}{poppurple!12!white}
\colorlet{poppinkL}{poppink!12!white}
\colorlet{popgoldL}{popgold!14!white}

\pagecolor{popcream}

% ============================================================
% Typographie
% ============================================================
\defaultfontfeatures{Ligatures=TeX}
\setmainfont{PlusJakartaSans-Medium.ttf}[Path=../../../fonts/,BoldFont=PlusJakartaSans-Bold.ttf]
% Symboles, API et emojis absents de la police principale : repli sur DejaVu Sans (caractères actifs)
\newfontface\symfont{DejaVuSans.ttf}[Path=../../../fonts/]
\makeatletter
\newcommand\symactive[1]{\catcode#1=13 \begingroup\lccode`\~=#1 \lowercase{\endgroup\def~}{{\symfont\char#1}}}
\newcommand\symblank[1]{\catcode#1=13 \begingroup\lccode`\~=#1 \lowercase{\endgroup\def~}{}}
\newcommand\symhyph[1]{\catcode#1=13 \begingroup\lccode`\~=#1 \lowercase{\endgroup\def~}{\mbox{-}}}
\newcount\symc
\newcommand\symrange[2]{\symc=#1 \@whilenum\symc<#2 \do{\expandafter\symactive\expandafter{\the\symc}\advance\symc by 1 }}
\newcommand\symblankrange[2]{\symc=#1 \@whilenum\symc<#2 \do{\expandafter\symblank\expandafter{\the\symc}\advance\symc by 1 }}
\makeatother
\symrange{"0250}{"02FF}\symactive{"2713}\symactive{"2714}\symactive{"2717}\symactive{"2718}\symactive{"2610}\symactive{"2611}\symactive{"2612}
\symrange{"2600}{"27BF}
\catcode"2705=13 \begingroup\lccode`\~="2705 \lowercase{\endgroup\def~}{{\symfont\char"2713}}\symblank{"26BD}\symblank{"26A1}
\symrange{"2460}{"257F}\symactive{"223C}\symactive{"2248}\symactive{"2260}
\symactive{"01F8}\symactive{"01F9}\symactive{"1E3E}\symactive{"1E3F}
\symhyph{"2011}\symhyph{"2E1A}\symblankrange{"1F300}{"1FAFF}\symblank{"FE0F}
% Caractères chinois : WenQuanYi Zen Hei (sous-ensemble GB2312 + pinyin), gras/italique simulés
\setCJKmainfont{WenQuanYiZenHei-subset.ttf}[Path=../../../fonts/,AutoFakeBold=3,AutoFakeSlant=0.2]
\xeCJKsetup{CJKspace=false}
% Pinyin : voyelles à caron (ǎ ǐ ǒ ǔ ǖ ǘ ǚ ǜ) absentes de la police principale -> caractères actifs
\newfontface\pyfont{WenQuanYiZenHei-subset.ttf}[Path=../../../fonts/]
\newcommand\pyactive[1]{\catcode#1=13 \begingroup\lccode`\~=#1 \lowercase{\endgroup\def~}{{\pyfont\char#1}}}
\pyactive{"01CD}\pyactive{"01CE}\pyactive{"01CF}\pyactive{"01D0}\pyactive{"01D1}\pyactive{"01D2}\pyactive{"01D3}\pyactive{"01D4}\pyactive{"01D5}\pyactive{"01D6}\pyactive{"01D7}\pyactive{"01D8}\pyactive{"01D9}\pyactive{"01DA}\pyactive{"01DB}\pyactive{"01DC}
% Maths en sans-serif (Fira Math) avec chiffres et lettres droites de Plus
% Jakarta, pour que les formules se fondent dans le texte.
\usepackage[math-style=ISO,bold-style=ISO]{unicode-math}
\setmathfont{FiraMath-Regular.otf}[Path=../../../fonts/]
\setmathfont{PlusJakartaSans-Medium.ttf}[Path=../../../fonts/,range={up/num,up/{Latin,latin},\mathup}]
\usepackage{icomma} % virgule décimale sans espace en mode math
% Caractères mathématiques Unicode absents des polices de texte (Plus Jakarta,
% Archivo Black) : rendus par leur équivalent mathématique, en texte comme en
% formule — sinon ils s'affichent en carré vide (ex. « Arithmétique dans ℤ »).
\usepackage{newunicodechar}
\newunicodechar{ℝ}{\ensuremath{\mathbb{R}}}
\newunicodechar{ℤ}{\ensuremath{\mathbb{Z}}}
\newunicodechar{ℂ}{\ensuremath{\mathbb{C}}}
\newunicodechar{ℕ}{\ensuremath{\mathbb{N}}}
\newunicodechar{ℚ}{\ensuremath{\mathbb{Q}}}
\newunicodechar{𝔻}{\ensuremath{\mathbb{D}}}
\newunicodechar{α}{\ensuremath{\alpha}}
\newunicodechar{β}{\ensuremath{\beta}}
\newunicodechar{γ}{\ensuremath{\gamma}}
\newunicodechar{δ}{\ensuremath{\delta}}
\newunicodechar{ε}{\ensuremath{\varepsilon}}
\newunicodechar{θ}{\ensuremath{\theta}}
\newunicodechar{λ}{\ensuremath{\lambda}}
\newunicodechar{μ}{\ensuremath{\mu}}
\newunicodechar{σ}{\ensuremath{\sigma}}
\newunicodechar{τ}{\ensuremath{\tau}}
\newunicodechar{φ}{\ensuremath{\varphi}}
\newunicodechar{ψ}{\ensuremath{\psi}}
\newunicodechar{ω}{\ensuremath{\omega}}
\newunicodechar{Σ}{\ensuremath{\Sigma}}
% majuscules produites par \MakeUppercase dans les titres (α-aminés -> Α)
\newunicodechar{Α}{\ensuremath{\alpha}}
\newunicodechar{Β}{\ensuremath{\beta}}
\newunicodechar{Γ}{\ensuremath{\Gamma}}
\newunicodechar{Δ}{\ensuremath{\Delta}}
\newunicodechar{Ε}{\ensuremath{\varepsilon}}
\newunicodechar{Θ}{\ensuremath{\Theta}}
\newunicodechar{Λ}{\ensuremath{\Lambda}}
\newunicodechar{Μ}{\ensuremath{\mu}}
\newunicodechar{Τ}{\ensuremath{\tau}}
\newunicodechar{Φ}{\ensuremath{\Phi}}
\newunicodechar{Ψ}{\ensuremath{\Psi}}
\newunicodechar{Ω}{\ensuremath{\Omega}}
\newunicodechar{↦}{\ensuremath{\mapsto}}
\newunicodechar{∈}{\ensuremath{\in}}
\newunicodechar{∉}{\ensuremath{\notin}}
\newunicodechar{∧}{\ensuremath{\wedge}}
\newunicodechar{⇔}{\ensuremath{\Leftrightarrow}}
\newunicodechar{⟺}{\ensuremath{\Longleftrightarrow}}
\newunicodechar{⇒}{\ensuremath{\Rightarrow}}
\newunicodechar{⟹}{\ensuremath{\Longrightarrow}}
\newunicodechar{∘}{\ensuremath{\circ}}
\newunicodechar{∪}{\ensuremath{\cup}}
\newunicodechar{∩}{\ensuremath{\cap}}
\newunicodechar{≡}{\ensuremath{\equiv}}
\newunicodechar{⊂}{\ensuremath{\subset}}
\newunicodechar{⊥}{\ensuremath{\perp}}
\newunicodechar{∃}{\ensuremath{\exists}}
\newunicodechar{∀}{\ensuremath{\forall}}
\newunicodechar{∛}{\ensuremath{\sqrt[3]{\,}}}
\newunicodechar{∜}{\ensuremath{\sqrt[4]{\,}}}
\newunicodechar{⃗}{\ensuremath{{}^{\to}}}
\newunicodechar{ⁿ}{\ifmmode^{n}\else\textsuperscript{\ensuremath{n}}\fi}
\newunicodechar{ˣ}{\ifmmode^{x}\else\textsuperscript{\ensuremath{x}}\fi}
\newunicodechar{ᵇ}{\ifmmode^{b}\else\textsuperscript{\ensuremath{b}}\fi}
\newunicodechar{ʳ}{\ifmmode^{r}\else\textsuperscript{\ensuremath{r}}\fi}
\newunicodechar{ᵘ}{\ifmmode^{u}\else\textsuperscript{\ensuremath{u}}\fi}
\newunicodechar{ʸ}{\ifmmode^{y}\else\textsuperscript{\ensuremath{y}}\fi}
\newunicodechar{ᵃ}{\ifmmode^{a}\else\textsuperscript{\ensuremath{a}}\fi}
\newunicodechar{ᵉ}{\ifmmode^{e}\else\textsuperscript{\ensuremath{e}}\fi}
\newunicodechar{ᵏ}{\ifmmode^{k}\else\textsuperscript{\ensuremath{k}}\fi}
\newunicodechar{ᵖ}{\ifmmode^{p}\else\textsuperscript{\ensuremath{p}}\fi}
\newunicodechar{ᴺ}{\ifmmode^{N}\else\textsuperscript{\ensuremath{N}}\fi}
\newunicodechar{ᶜ}{\ifmmode^{c}\else\textsuperscript{\ensuremath{c}}\fi}
\newunicodechar{ʰ}{\ifmmode^{h}\else\textsuperscript{\ensuremath{h}}\fi}
\newunicodechar{ᵠ}{\ifmmode^{q}\else\textsuperscript{\ensuremath{q}}\fi}
\newunicodechar{ₙ}{\ifmmode_{n}\else\textsubscript{\ensuremath{n}}\fi}
\newunicodechar{ₐ}{\ifmmode_{a}\else\textsubscript{\ensuremath{a}}\fi}
\newunicodechar{ᵢ}{\ifmmode_{i}\else\textsubscript{\ensuremath{i}}\fi}
\newunicodechar{ᵣ}{\ifmmode_{r}\else\textsubscript{\ensuremath{r}}\fi}
\newunicodechar{ₖ}{\ifmmode_{k}\else\textsubscript{\ensuremath{k}}\fi}
\newunicodechar{ᵦ}{\ifmmode_{\beta}\else\textsubscript{\ensuremath{\beta}}\fi}
\newunicodechar{ₓ}{\ifmmode_{x}\else\textsubscript{\ensuremath{x}}\fi}
\newfontfamily\popdisplay{ArchivoBlack-Regular.ttf}[Path=../../../fonts/]
\newfontfamily\poplogo{SpaceGrotesk-Bold.ttf}[Path=../../../fonts/]
\newfontfamily\popsemi{PlusJakartaSans-SemiBold.ttf}[Path=../../../fonts/]
\newfontfamily\popxbold{PlusJakartaSans-ExtraBold.ttf}[Path=../../../fonts/]
\newfontfamily\popxboldls{PlusJakartaSans-ExtraBold.ttf}[Path=../../../fonts/,LetterSpace=3.0]

\setlist[enumerate]{leftmargin=1.7em,itemsep=5pt,topsep=4pt}
\setlist[itemize]{leftmargin=1.7em,itemsep=4pt,topsep=4pt,label={\color{poporange}\textbullet}}
\setlength{\parindent}{0pt}
\setlength{\parskip}{5pt}
\renewcommand{\arraystretch}{1.25}

% pgfplots : style Pop par défaut
\pgfplotsset{
  every axis/.append style={axis line style={popink,line width=1pt},tick style={popink},
    grid style={popline,line width=0.5pt},label style={font=\small},tick label style={font=\footnotesize}},
  popcurve/.style={poporange,line width=1.8pt,no markers,smooth},
}
% Même style disponible dans un \draw TikZ simple (hors pgfplots)
\tikzset{popcurve/.style={poporange,line width=1.8pt}}
\tikzset{popgrid/.style={popline,very thin}}
\colorlet{popcurve}{poporange} % le nom du style est parfois utilisé comme couleur

% ============================================================
% Pill / squiggle
% ============================================================
\newcommand{\poppill}[3]{%
  \tikz[baseline=-0.55ex]\node[draw=popink,line width=1.2pt,rounded corners=2.6pt,%
    fill=#2,inner xsep=6.5pt,inner ysep=3pt]{%
    \popxboldls\fontsize{7}{7}\selectfont\color{#3}\MakeUppercase{#1}};%
}
\newcommand{\popsquiggle}[1]{%
  \tikz[baseline=-0.4ex]{
    \draw[#1,line width=2.6pt,line cap=round]
      plot [smooth] coordinates {(0,0.09) (0.34,0.01) (0.68,0.21) (1.02,0.01) (1.36,0.10)};
  }%
}
% Mot-clé surligné (marqueur orange sous le texte)
\newcommand{\cle}[1]{{\popxbold\color{popdark}#1}}

% ============================================================
% En-tête / pied
% ============================================================
\pagestyle{fancy}
\fancyhf{}
\renewcommand{\headrulewidth}{0pt}
\renewcommand{\footrulewidth}{0pt}
\lhead{\raisebox{-0.15ex}{\poplogo\fontsize{13}{13}\selectfont\color{popink}OnBuch\color{poporange}+}}
\rhead{\poppill{\DOCMATIERE\ \textbullet\ \DOCNIVEAU\ \textbullet\ Leçon \DOCLECON}{white}{popink}}
\lfoot{\popxbold\fontsize{7.6}{7.6}\selectfont\color{popmuted}\MakeUppercase{Cours signé OnBuch+}\ \textbullet\ {\popsemi\DOCTITRE}}
\rfoot{\tikz[baseline=-0.6ex]\node[rounded corners=8.5pt,fill=popink,minimum height=17pt,minimum width=17pt,inner xsep=5pt,inner sep=0]{\popxbold\fontsize{8}{8}\selectfont\color{popcream}\thepage\,/\,\pageref*{LastPage}};}

% ============================================================
% Titres
% ============================================================
\newcommand{\chaptitre}[2]{%
  \begin{tcolorbox}[enhanced,colback=poporange,colframe=popink,boxrule=2.6pt,arc=11pt,
      drop shadow=popink,left=15pt,right=15pt,top=12pt,bottom=11pt,before skip=3pt,after skip=18pt]
    \poppill{#2}{popink}{popcream}\par\vspace{9pt}
    {\raggedright\hyphenpenalty=10000\exhyphenpenalty=10000\popdisplay\fontsize{22}{24}\selectfont\color{popcream}\MakeUppercase{#1}\par}
  \end{tcolorbox}
}
% Section numérotée : pastille orange avec le numéro + titre Archivo
\newcounter{popsec}
\newcommand{\coursec}[1]{%
  \stepcounter{popsec}\setcounter{popsub}{0}%
  \par\vspace{16pt}\noindent
  \begin{minipage}{\linewidth}
  \tikz[baseline=-0.7ex]\node[draw=popink,line width=1.6pt,fill=poporange,rounded corners=4pt,minimum size=22pt,inner sep=0]{\popdisplay\fontsize{12}{12}\selectfont\color{popcream}\thepopsec};%
  \hspace{8pt}{\popdisplay\fontsize{15}{17}\selectfont\color{popink}\MakeUppercase{#1}}\par
  \vspace{4pt}\noindent{\color{poporange}\rule{36pt}{3.6pt}}
  \end{minipage}\par\nopagebreak\vspace{8pt}
}
\newcounter{popsub}
\newcommand{\courssub}[1]{%
  \stepcounter{popsub}%
  \par\vspace{9pt}\noindent{\popxbold\fontsize{11.5}{13}\selectfont\color{popink}{\color{poporange}\thepopsec.\thepopsub}\hspace{6pt}#1}\par\nopagebreak\vspace{4pt}
}

% ============================================================
% Boîtes pédagogiques — même recette : fond blanc, bordure épaisse colorée,
% ombre dure encre, badge plein. Titre optionnel [#1].
% ============================================================
\tcbset{popbase/.style={breakable,enhanced,colback=white,boxrule=2.3pt,arc=9pt,
  shadow={3pt}{-3pt}{0pt}{popink},left=14pt,right=14pt,top=11pt,bottom=11pt,before skip=11pt,after skip=15pt,
  fontupper=\popsemi\fontsize{10.6}{14.5}\selectfont\color{popink},
  fontlower=\popsemi\fontsize{10.6}{14.5}\selectfont\color{popink}}}
% Coche dessinée (le glyphe ✓ n'existe pas dans les polices du document)
\newcommand{\popcheck}[1]{\tikz[baseline=-0.5ex]\draw[#1,line width=1.6pt,line cap=round,line join=round](-0.13,0.0)--(-0.04,-0.09)--(0.13,0.1);}
\providecommand{\coche}{\popcheck{popgreen}}
\providecommand{\resultat}[1]{\par\smallskip\noindent\fcolorbox{popgreen}{popgreenL}{\parbox{\dimexpr\linewidth-2\fboxsep-2\fboxrule\relax}{\textbf{Résultat.} #1}}\par\smallskip}
\newcommand{\popboxhead}[3]{\poppill{#1}{#2}{white}\ifx\relax#3\relax\else\hspace{6pt}{\popxbold\fontsize{10.6}{12}\selectfont\color{popink}#3}\fi\par\vspace{7pt}}

\newtcolorbox{aretenir}[1][]{popbase,colframe=popgreen,before upper={\popboxhead{À retenir}{popgreen}{#1}}}
% Texte en chinois (document de lecture, dialogue, extrait) : cadre
% d'étude. Un titre optionnel (source, auteur) s'affiche dans l'en-tête.
\newtcolorbox{textebox}[1][]{popbase,colframe=popblue,before upper={\popboxhead{文本 · Texte}{popblue}{#1}},after upper={}}
% Marques ✓ / ✗ (forme correcte / fautive) souvent utilisées par les modèles
\providecommand{\cross}{\textcolor{poppink}{\ensuremath{\times}}}
\providecommand{\xmark}{\cross}\providecommand{\crossmark}{\cross}\providecommand{\cmark}{\ensuremath{\checkmark}}
% Ligne de réponse / justification à compléter (inventée par les modèles)
\providecommand{\justification}[1]{\par\noindent\textit{Justification~:} \dotfill\par}
% \textipa (package tipa non chargé) : la police ne couvre pas l'API -> simple passage
\providecommand{\textipa}[1]{#1}
% Soulignement (ulem non chargé) : \uline, \ul
\providecommand{\uline}[1]{\underline{#1}}\providecommand{\ul}[1]{\underline{#1}}
% \glossaire{\item ...} inventée par les modèles : liste de glossaire
\providecommand{\glossaire}[1]{\begin{itemize}#1\end{itemize}}
% \alert (beamer) inventée par les modèles : texte mis en évidence
\providecommand{\alert}[1]{\textbf{\textcolor{poppink}{#1}}}
% variantes ASCII d'ordinaux écrites en macro
\providecommand{\iere}{\textsuperscript{re}}\providecommand{\ieme}{\textsuperscript{e}}\providecommand{\ere}{\textsuperscript{re}}\providecommand{\eme}{\textsuperscript{e}}\providecommand{\er}{\textsuperscript{er}}\providecommand{\ie}{\textsuperscript{e}}
% Ordinaux français écrits en macro par les modèles : 1\ière, 3\ième
\newcommand{\ière}{\textsuperscript{re}}\newcommand{\ième}{\textsuperscript{e}}
% \exemple (inventée par les modèles) : étiquette « Exemple : »
\providecommand{\exemple}{\textbf{Exemple~:}\ }
% \square utilisable aussi hors mode math (cases à cocher)
\AtBeginDocument{\let\squareMath\square\renewcommand{\square}{\ensuremath{\squareMath}}}
% Mot ou phrase en chinois à l'intérieur d'un paragraphe français
\newcommand{\zh}[1]{\ifmmode\text{#1}\else{#1}\fi}
\let\eng\zh \let\esp\zh \let\all\zh % alias tolérés
% flèches d'intonation inventées par les modèles
\providecommand{\textsearrow}{}\renewcommand{\textsearrow}{\ensuremath{\searrow}}\providecommand{\textnearrow}{}\renewcommand{\textnearrow}{\ensuremath{\nearrow}}
% \fr{..} : mot français (inventé par les modèles)
\providecommand{\fr}[1]{\textit{#1}}
% zhbox : encadré de texte chinois inventé par les modèles
\newenvironment{zhbox}{\begin{quote}}{\end{quote}}
% \courssubsub : titre de niveau 3 inventé par les modèles
\providecommand{\courssubsub}[1]{\par\medskip\noindent\textbf{#1}\par\smallskip}
% \zhbf : texte chinois en gras inventé par les modèles
\providecommand{\zhbf}[1]{\textbf{#1}}
% \vs : « versus » inventé par les modèles
\providecommand{\vs}{\textit{vs}\ }
% \blank : case à compléter inventée par les modèles
\providecommand{\blank}[1][]{\underline{\hspace{1.6em}}}
\newtcolorbox{exemplebox}[1][]{popbase,colframe=poporange,before upper={\popboxhead{Exemple}{poporange}{#1}}}
\newtcolorbox{definition}[1][]{popbase,colframe=popblue,before upper={\popboxhead{Définition}{popblue}{#1}}}
\newtcolorbox{propriete}[1][]{popbase,colframe=poppurple,before upper={\popboxhead{Propriété}{poppurple}{#1}}}
\newtcolorbox{methode}[1][]{popbase,colframe=popgold,before upper={\popboxhead{Méthode}{popgold}{#1}}}
\newtcolorbox{attention}[1][]{popbase,colframe=poppink,before upper={\popboxhead{Attention}{poppink}{#1}}}
\newtcolorbox{experience}[1][]{popbase,colframe=popblue,colback=popblueL,before upper={\popboxhead{Expérience}{popblue}{#1}}}
% « Le savais-tu ? » : carte violette pleine (contexte, culture, Cameroun)
\newtcolorbox{savaistu}[1][]{popbase,colback=poppurple,colframe=popink,
  fontupper=\popsemi\fontsize{10.4}{14.2}\selectfont\color{white},
  before upper={\poppill{Le savais-tu\,?}{white}{poppurple}\ifx\relax#1\relax\else\hspace{6pt}{\popxbold\color{white}#1}\fi\par\vspace{7pt}}}
% Exercice résolu : énoncé en haut, \tcblower puis la solution sur fond crème
\newcounter{popexo}\newcounter{popexr}
\newtcolorbox{exoresolu}[1][]{popbase,colframe=poporange,colbacklower=popcream,
  segmentation style={popink,line width=1.2pt,dashed},
  before upper={\stepcounter{popexr}\popboxhead{Exercice résolu \thepopexr}{poporange}{#1}},
  before lower={\poppill{Solution}{popink}{popcream}\par\vspace{6pt}}}
% Exercice (non résolu en place) — la correction est dans la partie Corrigés
\newtcolorbox{exercice}[1][]{popbase,colframe=popink,
  before upper={\stepcounter{popexo}\popboxhead{Exercice \thepopexo}{popink}{#1}}}
% Corrigé
\newtcolorbox{corrige}[1][]{popbase,colframe=popgreen,colback=popgreenL,
  before upper={\popboxhead{Corrigé}{popgreen}{#1}}}
% Objectifs / prérequis / bilan
\newtcolorbox{objectifs}{popbase,colframe=popink,colback=poporangeL,
  before upper={\popboxhead{Objectifs de la leçon}{popink}{}}}
\newtcolorbox{prerequis}{popbase,colframe=popink,colback=popblueL,
  before upper={\popboxhead{Prérequis}{popblue}{}}}
\newtcolorbox{bilan}{popbase,colframe=popink,colback=white,boxrule=2.8pt,
  before upper={\popboxhead{Fiche bilan}{poporange}{L'essentiel à connaître pour l'examen}}}
% Figure encadrée avec légende
\newtcolorbox{popfigure}[1][]{enhanced,breakable=false,colback=white,colframe=popline,boxrule=1.4pt,arc=8pt,
  left=10pt,right=10pt,top=10pt,bottom=8pt,before skip=10pt,after skip=12pt,halign=center,
  lower separated=false,halign lower=center,
  fontlower=\popsemi\fontsize{9}{11.5}\selectfont\color{popmuted},
  #1}
\newcommand{\legende}[1]{\par\vspace{6pt}{\popsemi\fontsize{9}{11.5}\selectfont\color{popmuted}#1}}

% ============================================================
% Signature OnBuch+
% ============================================================
\newcommand{\poplogoinline}[1]{{\poplogo\fontsize{#1}{#1}\selectfont\color{popink}OnBuch\color{poporange}+}}
\newcommand{\signatureonbuch}{%
  \begin{tcolorbox}[enhanced,colback=popink,colframe=popink,boxrule=2.2pt,arc=10pt,
      drop shadow=poporange,left=14pt,right=14pt,top=10pt,bottom=10pt,before skip=0pt,after skip=0pt]
    \tikz[baseline=-0.7ex]\node[circle,fill=popgreen,draw=popcream,line width=1.3pt,minimum size=22pt,inner sep=0]{\popcheck{white}};%
    \hspace{9pt}\begin{minipage}[c]{0.62\linewidth}
      {\popxbold\fontsize{12}{13}\selectfont\color{popcream}Signé\ \textbullet\ {\poplogo OnBuch\color{poporange}+}}\par\vspace{2pt}
      {\raggedright\popsemi\fontsize{8}{10.5}\selectfont\color{popline}\MakeUppercase{Équipe pédagogique OnBuch+}\\\MakeUppercase{Conforme au programme officiel MINESEC}\par}
    \end{minipage}\hfill
    {\poplogo\fontsize{22}{22}\selectfont\color{popcream}OnBuch\color{poporange}+}
  \end{tcolorbox}%
}

% ============================================================
% Page de garde
% #1 titre · #2 matière · #3 niveau · #4 module · #5 n° leçon · #6 durée ·
% #7 description / objectifs (texte ou itemize)
% ============================================================
\newcommand{\pagegarde}[7]{%
  \thispagestyle{empty}
  \newgeometry{a4paper,margin=1.7cm,top=1.6cm,bottom=1.4cm}%
  \begin{tikzpicture}[remember picture,overlay]
    \fill[poporange!18] ([xshift=-3.2cm,yshift=-3.6cm]current page.north east) circle (4.6cm);
    \fill[poppurple!14] ([xshift=2.4cm,yshift=7.6cm]current page.south west) circle (3.8cm);
  \end{tikzpicture}%
  {\poplogo\fontsize{30}{30}\selectfont\color{popink}OnBuch\color{poporange}+}\hfill
  \raisebox{0.6ex}{\poppill{Cours signé \textbullet\ Premium}{popink}{popcream}}\par
  \vspace{1pt}\popsquiggle{poppurple}\par
  \vspace{8pt}
  \begin{tcolorbox}[enhanced,colback=poporange,colframe=popink,boxrule=2.8pt,arc=13pt,
      drop shadow=popink,left=18pt,right=18pt,top=12pt,bottom=13pt,after skip=10pt]
    \poppill{#2 \textbullet\ #3}{popink}{popcream}\hspace{5pt}\poppill{#4}{white}{popink}\par\vspace{8pt}
    {\popdisplay\fontsize{14}{14}\selectfont\color{popink}LEÇON #5}\par\vspace{4pt}
    {\raggedright\hyphenpenalty=10000\exhyphenpenalty=10000\popdisplay\fontsize{25}{27}\selectfont\color{popcream}\MakeUppercase{#1}\par}
    \vspace{8pt}
    {\popxbold\fontsize{9.5}{9.5}\selectfont\color{popink}\MakeUppercase{Durée conseillée : #6}}
  \end{tcolorbox}
  \begin{tcolorbox}[enhanced,colback=white,colframe=popink,boxrule=2.2pt,arc=10pt,
      drop shadow=popink,left=16pt,right=16pt,top=10pt,bottom=10pt,after skip=8pt]
    \poppill{Dans cette leçon}{poppurple}{white}\par\vspace{5pt}
    \popsemi\fontsize{9.3}{12.6}\selectfont\color{popink}#7
  \end{tcolorbox}
  \noindent\begin{minipage}[t]{0.487\linewidth}
    \begin{tcolorbox}[enhanced,colback=popgreen,colframe=popink,boxrule=2.2pt,arc=10pt,
        drop shadow=popink,left=11pt,right=11pt,top=10pt,bottom=10pt,height=3.1cm,valign=center]
      \tcbox[colback=white,colframe=popink,boxrule=1.3pt,arc=5pt,boxsep=0pt,nobeforeafter,
        left=3pt,right=3pt,top=3pt,bottom=3pt,valign=center]{\qrcode[height=1.9cm]{https://play.google.com/store/apps/details?id=com.luvvix.onbuch&hl=fr}}%
      \hspace{8pt}\begin{minipage}[c]{0.5\linewidth}
        {\popdisplay\fontsize{11}{12.5}\selectfont\color{white}\MakeUppercase{Télécharge OnBuch}}\par\vspace{4pt}
        {\raggedright\popsemi\fontsize{8.4}{11}\selectfont\color{white}Gratuit \textbullet\ Google Play\\Tuteur IA, annales, concours, résultats\par}
      \end{minipage}
    \end{tcolorbox}
  \end{minipage}\hfill
  \begin{minipage}[t]{0.487\linewidth}
    \begin{tcolorbox}[enhanced,colback=poppurple,colframe=popink,boxrule=2.2pt,arc=10pt,
        drop shadow=popink,left=11pt,right=11pt,top=10pt,bottom=10pt,height=3.1cm,valign=center]
      \tikz[baseline=-0.7ex]\node[circle,fill=white,draw=popink,line width=1.3pt,minimum size=1.6cm,inner sep=0]{\popdisplay\fontsize{24}{24}\selectfont\color{poppurple}+};%
      \hspace{8pt}\begin{minipage}[c]{0.6\linewidth}
        {\popdisplay\fontsize{11}{12.5}\selectfont\color{white}\MakeUppercase{Passe à OnBuch+}}\par\vspace{4pt}
        {\raggedright\popsemi\fontsize{8.4}{11}\selectfont\color{white}Tous les cours signés, Tuteur IA illimité, fiches PDF — onglet OnBuch+ de l'app\par}
      \end{minipage}
    \end{tcolorbox}
  \end{minipage}
  \vspace{8pt}
  \popcontenu
  \vfill
  \signatureonbuch
  \restoregeometry
  \newpage
}

% Rangée « Ce cours contient » (page de garde)
\newcommand{\poptile}[3]{% #1 couleur #2 gros texte #3 légende
  \begin{tcolorbox}[enhanced,colback=#1,colframe=popink,boxrule=2pt,arc=9pt,shadow={2.5pt}{-2.5pt}{0pt}{popink},
      left=6pt,right=6pt,top=9pt,bottom=9pt,halign=center,valign=center,height=1.75cm,width=\linewidth,nobeforeafter]
    {\popdisplay\fontsize{12.5}{13}\selectfont\color{white}\MakeUppercase{#2}}\par\vspace{3pt}
    {\centering\popsemi\fontsize{7.8}{9.5}\selectfont\color{white}#3\par}
  \end{tcolorbox}}
\newcommand{\popcontenu}{%
  \par\noindent{\popxboldls\fontsize{7.5}{7.5}\selectfont\color{popmuted}CE COURS SIGNÉ CONTIENT}\par\vspace{7pt}
  \noindent\begin{minipage}[t]{0.235\linewidth}\poptile{popblue}{Cours}{complet, illustré, pas à pas}\end{minipage}\hfill
  \begin{minipage}[t]{0.235\linewidth}\poptile{poporange}{Méthodes}{et exercices résolus}\end{minipage}\hfill
  \begin{minipage}[t]{0.235\linewidth}\poptile{poppink}{Exercices}{type examen + corrigés}\end{minipage}\hfill
  \begin{minipage}[t]{0.235\linewidth}\poptile{popgreen}{Fiche bilan}{l'essentiel à retenir}\end{minipage}\par}

% Sommaire
\newtcolorbox{sommaire}{enhanced,colback=white,colframe=popink,boxrule=2.2pt,arc=10pt,
  drop shadow=popink,left=18pt,right=18pt,top=13pt,bottom=13pt,before skip=6pt,after skip=20pt,
  before upper={\poppill{Sommaire}{poppurple}{white}\par\vspace{9pt}\popsemi\fontsize{10.6}{14.5}\selectfont\color{popink}}}

% Signature de fin de cours — poussée tout en bas de la dernière page
\newcommand{\findecours}{%
  \par\vfill\nopagebreak
  \begin{center}\popsquiggle{poporange}\end{center}\vspace{4pt}
  \signatureonbuch
}
\AtBeginDocument{\def\llbracket{\mathopen{[\mkern-3.5mu[}}\def\rrbracket{\mathclose{]\mkern-3.5mu]}}} % intervalles d'entiers (glyphe ⟦ absent de la police)
% Glyphes absents de Fira Math : redéfinis à partir de glyphes disponibles
\AtBeginDocument{%
  \def\setminus{\mathbin{\backslash}}\let\smallsetminus\setminus
  \def\vdots{\mathord{\vbox{\baselineskip3.2pt\lineskiplimit0pt\kern4pt\hbox{.}\hbox{.}\hbox{.}}}}%
  \def\ddots{\mathinner{\mkern1mu\raise6.5pt\hbox{.}\mkern2mu\raise3.6pt\hbox{.}\mkern2mu\raise0.7pt\hbox{.}\mkern1mu}}%
  \def\triangle{\mathord{\scalebox{0.9}{$\symup{\Delta}$}}}%
  \def\nabla{\mathord{\raisebox{\depth}{\scalebox{1}[-1]{$\symup{\Delta}$}}}}%
  \def\checkmark{\popcheck{popdark}}%
  \def\star{\mathbin{\ast}}\def\bigstar{\mathord{\ast}}%
  \def\diamond{\mathbin{\scalebox{0.6}{$\Diamond$}}}%
  \def\bigcup{\mathop{\mathchoice{\raisebox{-0.3em}{\scalebox{1.7}{$\cup$}}}{\scalebox{1.25}{$\cup$}}{\cup}{\cup}}\displaylimits}%
  \def\bigcap{\mathop{\mathchoice{\raisebox{-0.3em}{\scalebox{1.7}{$\cap$}}}{\scalebox{1.25}{$\cap$}}{\cap}{\cap}}\displaylimits}%
  \def\lesseqqgtr{\mathrel{\lessgtr}}\def\bigoplus{\mathop{\oplus}\displaylimits}%
}
\providecommand{\ssi}{si et seulement si} % abréviation fréquente
\AtBeginDocument{\providecommand{\wideparen}{\overparen}} % arc de cercle

\begin{document}

\pagegarde{Leçon 35 — Grammaire : 不管是…….或者是……. ; 是。。。。。。的 ; 比字句}{Chinois}{Tle A}{Module 5 : Mass médias et télécommunication}{ZH35}{2 h}{Cette leçon de grammaire présente trois structures essentielles du chinois moderne : la construction concessive 不管是……或者是……, la tournure d'insistance 是……的 et la phrase comparative en 比. À travers des exemples tirés de l'univers des médias et des télécommunications, les élèves apprennent à les construire, à les employer correctement et à éviter les pièges typiques des francophones.
\vspace{4pt}
\begin{multicols}{2}\small
\begin{itemize}
\item À la fin de cette leçon, je sais reconnaître et schématiser la structure 不管是……或者是……
\item À la fin de cette leçon, je sais construire une phrase concessive avec 都/也 en prédicat
\item À la fin de cette leçon, je sais employer la structure 是……的 pour insister sur le temps, le lieu, la manière ou l'agent d'une action passée
\item À la fin de cette leçon, je sais former la négation correcte de 是……的 (不是……的)
\item À la fin de cette leçon, je sais construire une phrase comparative avec 比 (A 比 B + adjectif)
\item À la fin de cette leçon, je sais exprimer le degré de la différence avec 得多、一点儿、多了
\end{itemize}\end{multicols}}

\begin{sommaire}
\begin{multicols}{2}
\begin{enumerate}
\item Observation guidée d'exemples authentiques
\item Structure 1 : 不管是……或者是……
\item Structure 2 : 是……的 (insistance sur une action passée)
\item Structure 3 : la phrase comparative en 比 (比字句)
\item Comparaison avec le français et pièges des francophones
\item Exercices d'application et de transformation
\item Activité d'intégration
\item Méthodes et exercices résolus
\item Exercices
\item Corrigés des exercices
\item Fiche bilan
\end{enumerate}
\end{multicols}
\end{sommaire}

\begin{objectifs}
\begin{itemize}
\item À la fin de cette leçon, je sais reconnaître et schématiser la structure 不管是……或者是……
\item À la fin de cette leçon, je sais construire une phrase concessive avec 都/也 en prédicat
\item À la fin de cette leçon, je sais employer la structure 是……的 pour insister sur le temps, le lieu, la manière ou l'agent d'une action passée
\item À la fin de cette leçon, je sais former la négation correcte de 是……的 (不是……的)
\item À la fin de cette leçon, je sais construire une phrase comparative avec 比 (A 比 B + adjectif)
\item À la fin de cette leçon, je sais exprimer le degré de la différence avec 得多、一点儿、多了
\item À la fin de cette leçon, je sais comparer avec le français et éviter les erreurs d'ordre des mots
\item À la fin de cette leçon, je sais utiliser ces trois structures dans des phrases et un court texte sur les médias
\end{itemize}
\end{objectifs}

\begin{prerequis}
\begin{itemize}
\item Les pronoms personnels, le verbe 是 et la négation 不/没 (Première)
\item Les adjectifs qualificatifs et leur emploi prédicatif sans 是 (很 + adjectif)
\item Les compléments de temps et de lieu placés avant le verbe
\item Le vocabulaire des médias du Module 5 : 电视、报纸、网络、手机、广播、信息
\item La structure 虽然……但是…… (concession simple, déjà vue)
\end{itemize}
\end{prerequis}

\newpage

\coursec{Observation guidée d'exemples authentiques}

Pendant les vacances, Manka'a, élève de Terminale A à Yaoundé, hésite devant son écran : « Que je regarde la télévision ou que j'écoute la radio, je veux absolument suivre les informations sur la Coupe du monde ! » Son correspondant chinois, Li Ming, lui répond aussitôt par message : \zh{不管是看电视，或者是听广播，都可以。} (\emph{Bùguǎn shì kàn diànshì, huòzhě shì tīng guǎngbō, dōu kěyǐ.} — « Que ce soit en regardant la télévision ou en écoutant la radio, c'est possible. »). Intriguée, Manka'a remarque aussi que Li Ming dit \zh{我是昨天收到你的信的。} (\emph{Wǒ shì zuótiān shōudào nǐ de xìn de.} — « C'est hier que j'ai reçu ta lettre. ») pour insister sur le moment, et \zh{网络比报纸快。} (\emph{Wǎngluò bǐ bàozhǐ kuài.} — « Internet est plus rapide que le journal. ») pour comparer les médias. Trois structures, trois outils puissants pour parler des mass médias comme un vrai sinophone !

\textbf{Problématique :} comment le chinois exprime-t-il (1) l'alternative indifférente « que... ou que... », (2) l'insistance sur les circonstances d'une action passée, et (3) la comparaison ? Observons d'abord, expliquons ensuite.

\courssub{Lecture à voix haute : trois exemples contextualisés « médias »}

Lisons ensemble, à voix haute, en respectant les tons. Chaque exemple est suivi de son pinyin et de sa traduction.

\begin{exemplebox}[Exemple 1 — l'alternative indifférente]
\zh{不管是看报纸，或者是上网，我们都能得到消息。}\\
\emph{Bùguǎn shì kàn bàozhǐ, huòzhě shì shàngwǎng, wǒmen dōu néng dédào xiāoxi.}\\
« Que l'on lise le journal ou que l'on aille sur Internet, nous pouvons obtenir l'information. »
\end{exemplebox}

\begin{exemplebox}[Exemple 2 — l'insistance sur une action passée]
\zh{这条新闻是我在网上看到的。}\\
\emph{Zhè tiáo xīnwén shì wǒ zài wǎngshàng kàndào de.}\\
« C'est sur Internet que j'ai vu cette nouvelle. »
\end{exemplebox}

\begin{exemplebox}[Exemple 3 — la comparaison]
\zh{手机比电脑便宜。}\\
\emph{Shǒujī bǐ diànnǎo piányi.}\\
« Le téléphone portable est moins cher que l'ordinateur. »
\end{exemplebox}

\begin{textebox}[Dialogue de Manka'a et Li Ming]
\zh{曼卡：李明，你平时怎么了解新闻？}\\
\emph{Mànkǎ : Lǐ Míng, nǐ píngshí zěnme liǎojiě xīnwén ?}\\
« Manka'a : Li Ming, comment t'informes-tu d'habitude ? »\\
\zh{李明：不管是看电视，或者是听广播，都可以了解新闻。}\\
\emph{Lǐ Míng : Bùguǎn shì kàn diànshì, huòzhě shì tīng guǎngbō, dōu kěyǐ liǎojiě xīnwén.}\\
« Li Ming : Que ce soit en regardant la télévision ou en écoutant la radio, on peut s'informer. »\\
\zh{曼卡：这条消息你是在哪儿看到的？}\\
\emph{Mànkǎ : Zhè tiáo xiāoxi nǐ shì zài nǎr kàndào de ?}\\
« Manka'a : C'est où que tu as vu cette information ? »\\
\zh{李明：我是在网上看到的。网络比报纸快，手机比电视方便。}\\
\emph{Lǐ Míng : Wǒ shì zài wǎngshàng kàndào de. Wǎngluò bǐ bàozhǐ kuài, shǒujī bǐ diànshì fāngbiàn.}\\
« Li Ming : C'est sur Internet que je l'ai vue. Internet est plus rapide que le journal, et le téléphone portable est plus pratique que la télévision. »
\end{textebox}

\courssub{Vocabulaire des exemples}

\begin{center}
\begin{tabularx}{\linewidth}{|X|X|X|X|}
\hline
\rowcolor{popblueL} Caractères & Pinyin & Nature & Traduction \\
\hline
报纸 & bàozhǐ & nom & journal (papier) \\
\hline
上网 & shàngwǎng & verbe & aller sur Internet \\
\hline
消息 & xiāoxi & nom & information, nouvelle \\
\hline
新闻 & xīnwén & nom & nouvelle, actualité \\
\hline
广播 & guǎngbō & nom & radio (diffusion) \\
\hline
网络 & wǎngluò & nom & réseau, Internet \\
\hline
手机 & shǒujī & nom & téléphone portable \\
\hline
电脑 & diànnǎo & nom & ordinateur \\
\hline
便宜 & piányi & adjectif & bon marché \\
\hline
方便 & fāngbiàn & adjectif & pratique, commode \\
\hline
收到 & shōudào & verbe & recevoir \\
\hline
条 & tiáo & classificateur & (pour nouvelles, rues, poissons) \\
\hline
\end{tabularx}
\end{center}

\courssub{Questions d'observation guidée}

Avant toute règle, répondons collectivement à ces questions en regardant les trois exemples :

\begin{enumerate}
\item Dans l'exemple 1, où se place \zh{不管} ? (Réponse attendue : en tête de phrase, suivi de \zh{是}.)
\item Quel mot introduit la deuxième possibilité ? (Réponse : \zh{或者是}.)
\item Quel petit mot accompagne \cle{toujours} la deuxième partie de la structure 1, juste avant le verbe principal ? (Réponse : \zh{都} « tous, dans tous les cas ».)
\item Dans l'exemple 2, où sont \zh{是} et \zh{的} ? Qu'y a-t-il entre les deux ? (Réponse : \zh{是} avant l'élément mis en relief, \zh{的} en fin de phrase ; entre eux, les circonstances : qui, où, quand, comment.)
\item Dans l'exemple 3, où se place \zh{比} ? Que trouve-t-on après lui ? (Réponse : \zh{比} se place entre les deux termes comparés ; après lui, le terme de référence, puis l'adjectif.)
\end{enumerate}

\begin{methode}[Observer avant d'expliquer : la méthode des couleurs]
Pour découvrir une structure nouvelle, procédez en trois étapes :
\begin{enumerate}
\item \textbf{Soulignez en rouge les mots-outils} : \zh{不管}, \zh{或者是}, \zh{都}, \zh{是}, \zh{的}, \zh{比}. Ce sont les « piliers » de la structure.
\item \textbf{Soulignez en bleu les verbes} : \zh{看}, \zh{听}, \zh{收到}, \zh{看到}. Ce sont les « briques » de sens.
\item \textbf{Dessinez le schéma} : remplacez les briques par des cases vides (A, B, verbe...) et gardez les piliers. Vous obtenez le schéma général, valable pour toutes les phrases du même type.
\end{enumerate}
C'est seulement après cette observation que la règle est énoncée : vous la comprendrez alors comme une évidence.
\end{methode}

\courssub{Mise en évidence collective : les trois schémas au tableau}

\begin{popfigure}
\begin{tikzpicture}[mindmap, grow cyclic, every node/.style={concept, align=center},
root concept/.append style={concept color=popblue, font=\bfseries},
level 1/.append style={level distance=4.6cm, sibling angle=120},
level 1 concept/.append style={concept color=poporange},
edge from parent/.style={concept color=popmuted}]
\node[root concept, align=center]{3 structures\\ médias}
child { node[align=center]{不管是…\\ 或者是…\\ …都…}
  child[concept color=popgreen] { node[concept color=popgreen, font=\small, align=center]{不管是看报纸，\\ 或者是上网，\\ 我们都能得到消息。} }
}
child { node[align=center]{是……的\\ (insistance)}
  child[concept color=popgreen] { node[concept color=popgreen, font=\small, align=center]{这条新闻是\\ 我在网上看到的。} }
}
child { node[align=center]{比字句\\ (comparaison)}
  child[concept color=popgreen] { node[concept color=popgreen, font=\small] {手机比电脑便宜。} }
};
\end{tikzpicture}
\legende{Carte mentale des trois structures de la leçon, avec un mini-exemple par branche.}
\end{popfigure}

\begin{aretenir}[Les trois schémas découverts]
\begin{itemize}
\item \textbf{Structure 1 :} 不管(是) + A，或者是 + B，…… \cle{都} / 也 + verbe. — « Que... ou que..., dans tous les cas... »
\item \textbf{Structure 2 :} (Sujet) + \cle{是} + circonstances (temps, lieu, manière, agent) + verbe + \cle{的}. — « C'est... que... » (action passée).
\item \textbf{Structure 3 :} A + \cle{比} + B + adjectif. — « A est plus... que B. »
\end{itemize}
\end{aretenir}

\begin{center}
\begin{tabularx}{\linewidth}{|X|X|X|}
\hline
\rowcolor{popblueL} Structure & Exemple (caractères + pinyin) & Traduction \\
\hline
不管是…或者是…都… & 不管是看报纸，或者是上网，我们都能得到消息。\newline \emph{Bùguǎn shì kàn bàozhǐ, huòzhě shì shàngwǎng, wǒmen dōu néng dédào xiāoxi.} & Que l'on lise le journal ou qu'on aille sur Internet, on obtient l'information. \\
\hline
是……的 & 我是昨天收到你的信的。\newline \emph{Wǒ shì zuótiān shōudào nǐ de xìn de.} & C'est hier que j'ai reçu ta lettre. \\
\hline
A 比 B + adj. & 网络比报纸快。\newline \emph{Wǎngluò bǐ bàozhǐ kuài.} & Internet est plus rapide que le journal. \\
\hline
\end{tabularx}
\end{center}

\begin{exoresolu}[Repérage de structures]
Énoncé : dans les phrases suivantes, soulignez les mots-outils et identifiez la structure (1, 2 ou 3).\\
a) 不管是打电话，或者是发短信，都很方便。\\
b) 这条新闻是我在网上看到的。\\
c) 手机比电脑便宜。
\tcblower
Solution détaillée :\\
a) Mots-outils en rouge : \textbf{不管是}…\textbf{或者是}…\textbf{都}. Verbes en bleu : 打电话, 发短信. Structure 1 (alternative indifférente). Traduction : « Que ce soit en téléphonant ou en envoyant un SMS, c'est très pratique. »\\
b) Mots-outils : \textbf{是}…\textbf{的} ; entre eux : 我在网上 (qui + où) ; verbe : 看到. Structure 2 (insistance sur le lieu d'une action passée). Traduction : « C'est sur Internet que j'ai vu cette nouvelle. »\\
c) Mot-outil : \textbf{比} ; A = 手机, B = 电脑, adjectif = 便宜. Structure 3 (comparaison). Traduction : « Le portable est moins cher que l'ordinateur. »
\end{exoresolu}

\coursec{Structure 1 : 不管是……或者是…… (Alternative indifférente)}

\begin{propriete}[Schéma et emplois]
\textbf{Schéma de base :} \zh{不管是} + \cle{Choix 1} + \zh{，或者是} + \cle{Choix 2} + \zh{，} + \cle{Sujet} + \zh{都 / 也} + \cle{Prédicat}.\\
\textbf{Variante formelle :} \zh{无论} (wúlùn) remplace \zh{不管} ; \zh{还是} (háishì) remplace \zh{或者是} à l'écrit soutenu.\\
\textbf{Emploi :} Exprime que le résultat est le même quelle que soit l'alternative choisie. Les deux choix sont souvent des verbes ou des phrases verbales (VP). Le \zh{是} après \zh{不管} est facultatif si le choix est un verbe simple (\zh{不管看报纸...}), mais recommandé pour la clarté en Terminale.
\end{propriete}

\begin{exemplebox}[Variantes de position du sujet]
\begin{itemize}
\item \zh{我们不管是看报纸，或者是上网，都能得到消息。} (Sujet avant \zh{不管} = Thème)
\item \zh{不管是看报纸，或者是上网，我们都能得到消息。} (Sujet après la virgule = Rhème, plus courant à l'oral)
\end{itemize}
\end{exemplebox}

\begin{propriete}[Négation]
La négation porte sur le prédicat final, \cle{jamais} sur \zh{不管} ou \zh{或者是}.
\zh{不管是看报纸，或者是上网，我们都不能 / 也不能 得到消息。}
\end{propriete}

\begin{attention}[Pièges francophones — Structure 1]
\begin{itemize}
\item \textbf{Oublier \zh{都} / \zh{也}} : *\zh{不管是A，或者是B，我们可以。} $\rightarrow$ Incorrect. Le corréla \zh{都} (ou \zh{也}) est \cle{obligatoire} pour fermer la structure.
\item \textbf{Confondre \zh{或者是} et \zh{还是}} : \zh{还是} s'utilise dans les questions alternatives (\zh{你喝茶还是咖啡？}). Dans l'affirmation « que... ou... », on utilise \zh{或者是} (ou \zh{或者} seul).
\item \textbf{Ordre des mots} : Ne mettez pas le sujet \cle{entre} \zh{不管} et \zh{或者是}. *\zh{不管我们是看报纸，或者是上网...} $\rightarrow$ Lourd/incorrect. Sujet avant tout ou après la virgule.
\end{itemize}
\end{attention}

\begin{exercice}[Entraînement Structure 1]
Complétez avec \zh{不管是}，\zh{或者是}，\zh{都} / \zh{也} :\\
1. \_\_\_\_ 听广播，\_\_\_\_ 看电视，我们 \_\_\_\_ 能听到新闻。\\
2. \_\_\_\_ 去北京，\_\_\_\_ 去上海，他 \_\_\_\_ 要坐飞机。\\
3. Traduisez : « Que ce soit par SMS ou par WeChat, c'est rapide. » (indice : 发短信 fā duǎnxìn, 微信 Wēixìn).
\end{exercice}

\coursec{Structure 2 : 是……的 (Insistance sur une action passée)}

\begin{propriete}[Schéma et principe fondamental]
\textbf{Schéma :} \cle{Sujet / Thème} + \zh{是} + \cle{Élément mis en relief (Temps / Lieu / Manière / Agent / But)} + \cle{Verbe} + \zh{的}.\\
\textbf{Principe :} Cette structure \cle{présuppose} que l'action est \textbf{accomplie} (passé). Elle sert à \textbf{mettre en relief} (focaliser) une circonstance, non l'action elle-même. C'est l'équivalent du \textbf{« C'est... que... »} français (phrase clivée).
\end{propriete}

\begin{exemplebox}[Les quatre types de focalisation courants]
\begin{itemize}
\item \textbf{Temps :} \zh{我是昨天买的手机。} (\emph{Wǒ shì zuótiān mǎi de shǒujī.} — C'est hier que j'ai acheté le portable.)
\item \textbf{Lieu :} \zh{这条新闻是在网上看到的。} (\emph{Zhè tiáo xīnwén shì zài wǎngshàng kàndào de.} — C'est sur Internet qu'on a vu cette nouvelle.)
\item \textbf{Manière / Moyen :} \zh{他是坐飞机去的。} (\emph{Tā shì zuò fēijī qù de.} — C'est en avion qu'il y est allé.)
\item \textbf{Agent (sujet logique) :} \zh{这封信是王老师写的。} (\emph{Zhè fēng xìn shì Wáng lǎoshī xiě de.} — C'est M. Wang qui a écrit cette lettre.)
\end{itemize}
\end{exemplebox}

\begin{propriete}[Règles strictes — « Trois interdits »]
\begin{enumerate}
\item \textbf{Pas de \zh{了} après le verbe.} La structure \zh{是...的} marque déjà l'accompli. *\zh{我是昨天买了手机的。} $\rightarrow$ \textbf{Faux}.
\item \textbf{Pas de complément de durée / fréquence entre \zh{是} et \zh{的}} (ex: \zh{两小时}, \zh{三次}). On dit : \zh{我学了两小时中文} (phrase neutre), pas *\zh{我是两小时学中文的}.
\item \textbf{Négation} : \zh{不是...的}. \zh{我不是昨天买的手机，是前天买的。}
\end{enumerate}
\end{propriete}

\begin{attention}[Pièges francophones — Structure 2]
\begin{itemize}
\item \textbf{Traduire « C'est... que... » par \zh{是...是...}} : *\zh{是昨天是我买的手机。} $\rightarrow$ Un seul \zh{是} (le premier), le second est \zh{的}.
\item \textbf{Placer \zh{的} après le verbe seulement} : *\zh{我是昨天买的手机。} (Correct ici car verbe transitif + objet). Mais si verbe intransitif : \zh{我是昨天来的} (pas de \zh{的} après \zh{来} ? Si, \zh{的} est en fin de phrase : \zh{我是昨天来的}. Toujours \zh{的} final).
\item \textbf{Oublier le \zh{的} final} : *\zh{我是昨天买手机。} $\rightarrow$ Phrase incomplète, sonne comme « Je suis hier acheter téléphone » (non-sens).
\end{itemize}
\end{attention}

\begin{exercice}[Entraînement Structure 2]
Transformez la phrase neutre en phrase avec \zh{是...的} en focalisant l'élément entre parenthèses :\\
1. 我 (在图书馆) 看的书。 $\rightarrow$ \_\_\_\_\_\_\_\_\_\_\_\_\\
2. 他 (坐地铁) 去学校。 $\rightarrow$ \_\_\_\_\_\_\_\_\_\_\_\_\\
3. 他们 (上个星期) 到了北京。 $\rightarrow$ \_\_\_\_\_\_\_\_\_\_\_\_\\
4. Négation : « Ce n'est pas moi qui ai cassé le verre. » (indice : 打破 dǎpò, 杯子 bēizi).
\end{exercice}

\coursec{Structure 3 : La phrase comparative en 比 (比字句)}

\begin{propriete}[Schéma de base et degré]
\textbf{Schéma de base :} A + \zh{比} + B + \cle{Adjectif (verbe d'état)}.\\
\textbf{Traduction :} « A est plus [Adj] que B. » \textbf{Jamais de \zh{很} (hěn) avant l'adjectif} dans la phrase affirmative simple.
\begin{center}
\begin{tabularx}{\linewidth}{|X|X|X|}
\hline
\rowcolor{popblueL} Degré & Structure & Exemple \\
\hline
Supériorité simple & A 比 B Adj & \zh{手机比电脑便宜。} \\
\hline
Supériorité marquée (un peu) & A 比 B Adj \zh{一点儿 / 一些} & \zh{手机比电脑便宜\textbf{一点儿}。} \\
\hline
Supériorité nette (beaucoup) & A 比 B Adj \zh{得多 / 多了} & \zh{网络比报纸快\textbf{得多}。} \\
\hline
\end{tabularx}
\end{center}
\end{propriete}

\begin{propriete}[Négation et égalité]
\begin{itemize}
\item \textbf{Négation (A moins Adj que B) :} A + \zh{没有 / 不比} + B + Adj. \textit{Nuance :} \zh{没有} = « pas aussi... que » (différence objective) ; \zh{不比} = « pas plus... que » (réfutation).
\item \textbf{Égalité :} A + \zh{跟 / 和} + B + \zh{一样} + Adj. (\zh{手机跟电脑一样贵。})
\item \textbf{Superlatif :} A + \zh{最} + Adj. (\zh{手机最方便。}) — hors \zh{比} mais lié.
\end{itemize}
\end{propriete}

\begin{exemplebox}[Comparaison avec objets différents (ellipse de B)]
\zh{这部手机比那部贵。} (\emph{Zhè bù shǒujī bǐ nà bù guì.} — Celui-ci est plus cher que celui-là.)\\
\zh{今天比昨天冷。} (\emph{Jīntiān bǐ zuótiān lěng.} — Aujourd'hui il fait plus froid qu'hier.)
\end{exemplebox}

\begin{attention}[Pièges francophones — Structure 3]
\begin{itemize}
\item \textbf{Ajouter \zh{很} :} *\zh{手机比电脑很便宜。} $\rightarrow$ \textbf{Interdit}. L'adjectif est nu dans la comparatif simple.
\item \textbf{Inverser A et B :} *\zh{电脑比手机便宜} (si on veut dire « le portable est moins cher »). \zh{比} introduit \cle{toujours} le terme de référence (le second terme de comparaison).
\item \textbf{Confondre \zh{没有} et \zh{不比}} : \zh{我没有他高} (Je ne suis pas aussi grand que lui = il est plus grand). \zh{我不比他高} (Je ne suis pas plus grand que lui = je suis aussi grand ou plus petit).
\item \textbf{Comparer des noms sans adjectif :} *\zh{我比你苹果。} $\rightarrow$ Il faut un adjectif/verbe d'état : \zh{我比你多两个苹果。} / \zh{我的苹果比你的多。}
\end{itemize}
\end{attention}

\begin{exercice}[Entraînement Structure 3]
1. Mettez au comparatif (supériorité simple) :\\
   a) 这本书，那本书，贵 (这本书 \_\_\_\_ 那本书 \_\_\_\_) \\
   b) 法语，中文，难 (法语 \_\_\_\_ 中文 \_\_\_\_) \\
2. Ajoutez le degré indiqué :\\
   a) 网络比报纸快 (beaucoup) $\rightarrow$ \_\_\_\_\_\_\_\_\_\_\_\_\\
   b) 手机比电脑便宜 (un peu) $\rightarrow$ \_\_\_\_\_\_\_\_\_\_\_\_\\
3. Négation : « Internet n'est pas aussi cher que la télévision. » (indice : 贵 guì).
\end{exercice}

\coursec{Comparaison synthétique avec le français et pièges transversaux}

\begin{center}
\begin{tabularx}{\linewidth}{|X|X|X|}
\hline
\rowcolor{popblueL} Notion & Français (Structure) & Chinois (Structure + Mots-outils) \\
\hline
Alternative indifférente & Subjonctif : « Que je \textbf{regarde} la TV \textbf{ou que j'écoute} la radio... » & Indicatif + Corréla : \zh{不管是看电视，或者是听广播，\textbf{都}可以。} \\
\hline
Insistance (Passé) & Clivé : « \textbf{C'est hier que} j'ai reçu... » & \zh{是} + Circonstance + V + \zh{的} : \zh{我\textbf{是}昨天收到\textbf{的}。} \\
\hline
Comparaison & « A \textbf{est plus} Adj \textbf{que} B » (2 mots : plus... que) & A \zh{比} B Adj (1 mot : \zh{比} avant B, Adj nu) \\
\hline
\end{tabularx}
\end{center}

\begin{aretenir}[Check-list « Zéro faute » pour l'examen]
\begin{enumerate}
\item \zh{不管是...或者是...} $\rightarrow$ Vérifier la présence de \zh{都} / \zh{也} avant le verbe principal.
\item \zh{是...的} $\rightarrow$ Vérifier : Action passée ? \zh{的} final présent ? \zh{了} absent ? Élément focalisé bien placé après \zh{是} ?
\item \zh{比} $\rightarrow$ Vérifier : Ordre A \zh{比} B Adj ? \zh{很} absent ? Degré (\zh{一点儿}, \zh{得多}) après l'adjectif ?
\end{enumerate}
\end{aretenir}

\begin{savaistu}[Le téléphone mobile, roi des médias en Chine]
En Chine, plus d'un milliard de personnes sont connectées à Internet, et l'immense majorité s'informe d'abord via le téléphone mobile (\zh{手机}) plutôt que par la télévision ou le journal papier. Au Cameroun aussi, le smartphone est devenu le premier « poste de radio et de télévision » des jeunes, de Yaoundé à Buea. C'est pourquoi le comparatif en \zh{比} est partout dans les discussions sur les médias : \zh{手机比电视方便。} (\emph{Shǒujī bǐ diànshì fāngbiàn.} — « Le portable est plus pratique que la télévision. »)
\end{savaistu}

\coursec{Exercices d'application et de transformation (Niveau Examen)}

\begin{exercice}[Transformation de phrases]
Transformez les phrases suivantes selon le modèle indiqué :\\
\textbf{Modèle A (Alternative) :} 看电视可以，听广播也可以。 $\rightarrow$ \zh{不管是看电视，或者是听广播，都可以。}\\
1. 打电话方便，发微信也方便。 $\rightarrow$ \_\_\_\_\_\_\_\_\_\_\_\_\\
2. 去年去北京，今年去上海，都坐飞机。 $\rightarrow$ \_\_\_\_\_\_\_\_\_\_\_\_\\

\textbf{Modèle B (Insistance 是...的) :} 我昨天在网上看到了这条新闻。 (Focaliser sur le lieu) $\rightarrow$ \zh{这条新闻是我在网上看到的。}\\
3. 李明上个星期在图书馆借了这本书。 (Focaliser sur le temps) $\rightarrow$ \_\_\_\_\_\_\_\_\_\_\_\_\\
4. 妈妈用手机拍了这张照片。 (Focaliser sur l'agent) $\rightarrow$ \_\_\_\_\_\_\_\_\_\_\_\_\\

\textbf{Modèle C (Comparatif) :} 手机便宜，电脑贵。 $\rightarrow$ \zh{手机比电脑便宜。}\\
5. 网络快，报纸慢。 $\rightarrow$ \_\_\_\_\_\_\_\_\_\_\_\_\\
6. 这部电影有意思，那部电影没意思。 $\rightarrow$ \_\_\_\_\_\_\_\_\_\_\_\_\\
\end{exercice}

\begin{exercice}[Traduction thème (Français $\rightarrow$ Chinois)]
Traduisez en chinois (caractères + pinyin) :\\
1. Que ce soit pour regarder les informations ou pour discuter avec des amis, le portable est très pratique.\\
2. C'est par WeChat que j'ai reçu ton message hier soir.\\
3. La radio est moins rapide qu'Internet, mais plus pratique que le journal.\\
4. Ce n'est pas moi qui ai allumé la télévision, c'est mon petit frère.\\
5. Regarder la télé n'est pas aussi fatigant qu'étudier le chinois. (indice : 累 lèi, 学习 xuéxí).
\end{exercice}

\begin{exercice}[Expression écrite guidée : « Mon média préféré »]
Rédigez un court paragraphe (60-80 caractères) pour présenter votre média préféré. Utilisez \textbf{au moins une fois chacune des trois structures} étudiées.\\
\textbf{Plan suggéré :}
\begin{enumerate}
\item Alternative : « Que ce soit... ou..., je peux m'informer. » (\zh{不管是...或者是...都...})
\item Insistance : « C'est [quand/où/comment] que j'ai découvert cette info. » (\zh{是...的})
\item Comparaison : « Ce média est plus... que l'autre. » (\zh{比})
\end{enumerate}
\textbf{Mots-clés fournis :} 获取信息 (huòqǔ xīnxí - obtenir l'info), 实时 (shíshí - en temps réel), 深度 (shēndù - en profondeur).
\end{exercice}

\begin{corrige}[Exercices d'application et de transformation]
\textbf{Transformation :}\\
1. \zh{不管是打电话，或者是发微信，都很方便。} (\emph{Bùguǎn shì dǎ diànhuà, huòzhě shì fā Wēixìn, dōu hěn fāngbiàn.})\\
2. \zh{不管是去年去北京，或者是今年去上海，都坐飞机。} (\emph{Bùguǎn shì qùnián qù Běijīng, huòzhě shì jīnnián qù Shànghǎi, dōu zuò fēijī.})\\
3. \zh{这本书是李明上个星期在图书馆借的。} (\emph{Zhè běn shū shì Lǐ Míng shàng gè xīngqī zài túshūguǎn jiè de.})\\
4. \zh{这张照片是妈妈用手机拍的。} (\emph{Zhè zhāng zhàopiàn shì māma yòng shǒujī pāi de.})\\
5. \zh{网络比报纸快。} (\emph{Wǎngluò bǐ bàozhǐ kuài.})\\
6. \zh{这部电影比那部电影有意思。} (\emph{Zhè bù diànyǐng bǐ nà bù diànyǐng yǒuyìsi.})\\

\textbf{Traduction :}\\
1. \zh{不管是看新闻，或者是跟朋友聊天，手机都很方便。} (\emph{Bùguǎn shì kàn xīnwén, huòzhě shì gēn péngyou liáotiān, shǒujī dōu hěn fāngbiàn.})\\
2. \zh{这条信息是我昨天晚上用微信收到的。} (\emph{Zhè tiáo xìnxī shì wǒ zuótiān wǎnshàng yòng Wēixìn shōudào de.})\\
3. \zh{广播没有网络快，但是比报纸方便。} (\emph{Guǎngbō méiyǒu wǎngluò kuài, dànshì bǐ bàozhǐ fāngbiàn.})\\
4. \zh{不是我开的电视，是我弟弟开的。} (\emph{Bú shì wǒ kāi de diànshì, shì wǒ dìdi kāi de.})\\
5. \zh{看电视没有学习中文累。} (\emph{Kàn diànshì méiyǒu xuéxí Zhōngwén lèi.})\\

\textbf{Expression écrite (Exemple de production attendue) :}\\
\zh{我喜欢用手机上网。不管是看新闻，或者是看视频，都能获取信息。这条新闻是我在微信上看到的。手机比电脑方便，比报纸快，所以我最喜欢手机。}\\
(\emph{Wǒ xǐhuan yòng shǒujī shàngwǎng. Bùguǎn shì kàn xīnwén, huòzhě shì kàn shìpín, dōu néng huòqǔ xīnxi. Zhè tiáo xīnwén shì wǒ zài Wēixìn shàng kàndào de. Shǒujī bǐ diànnǎo fāngbiàn, bǐ bàozhǐ kuài, suǒyǐ wǒ zuì xǐhuan shǒujī.})
\end{corrige}

\begin{corrige}[Entraînement Structure 1]
1. \zh{不管是听广播，或者是看电视，我们都能听到新闻。}\\
2. \zh{不管是去北京，或者是去上海，他都要坐飞机。}\\
3. \zh{不管是发短信，或者是用微信，都很快。} (\emph{Bùguǎn shì fā duǎnxìn, huòzhě shì yòng Wēixìn, dōu hěn kuài.})
\end{corrige}

\begin{corrige}[Entraînement Structure 2]
1. \zh{这本书是我在图书馆看的。}\\
2. \zh{他是坐地铁去学校的。}\\
3. \zh{他们是上个星期到北京的。}\\
4. \zh{不是我打破的杯子，是我弟弟打破的。} (\emph{Bú shì wǒ dǎpò de bēizi, shì wǒ dìdi dǎpò de.})
\end{corrige}

\begin{corrige}[Entraînement Structure 3]
1. a) \zh{这本书比那本书贵。} b) \zh{法语比中文难。} (ou \zh{中文比法语难} selon l'avis de l'élève, les deux sont grammaticalement corrects).\\
2. a) \zh{网络比报纸快得多 / 多了。} b) \zh{手机比电脑便宜一点儿 / 一些。}\\
3. \zh{网络没有电视贵。} (ou \zh{网络不比电视贵。})
\end{corrige}

\coursec{Structure 1 : 不管是……或者是……}

Dans la section précédente, nous avons observé des phrases authentiques tirées d'un dialogue sur les médias. L'une d'elles a particulièrement retenu notre attention : \zh{不管是报纸，或者是电视，喀麦隆人都喜欢关注足球新闻。} \emph{(Bùguǎn shì bàozhǐ, huòzhě shì diànshì, Kāmàilóng rén dōu xǐhuan guānzhù zúqiú xīnwén.)} « Que ce soit le journal ou la télévision, les Camerounais aiment suivre les nouvelles du football. » Cette phrase exprime une idée très fréquente : peu importe le choix entre deux possibilités, la conclusion reste la même. C'est exactement le rôle de la structure \cle{不管是……或者是……}, première structure grammaticale de cette leçon. Observons-la d'abord dans un court texte, puis dégageons-en la règle.

\courssub{Observation : un texte authentique}

\begin{textebox}[Les habitudes médiatiques à Yaoundé]
\zh{不管是年轻人，或者是老年人，大家都用手机看新闻。}\emph{Bùguǎn shì niánqīngrén, huòzhě shì lǎoniánrén, dàjiā dōu yòng shǒujī kàn xīnwén.}\\
\zh{不管是早上，或者是晚上，我爸爸都听广播。}\emph{Bùguǎn shì zǎoshang, huòzhě shì wǎnshang, wǒ bàba dōu tīng guǎngbō.}\\
\zh{不管是法语节目，或者是英语节目，她都看得懂。}\emph{Bùguǎn shì Fǎyǔ jiémù, huòzhě shì Yīngyǔ jiémù, tā dōu kàndedǒng.}\\
\zh{不管是晴天还是雨天，他都骑自行车去买报纸。}\emph{Bùguǎn shì qíngtiān háishi yǔtiān, tā dōu qí zìxíngchē qù mǎi bàozhǐ.}
\end{textebox}

\textbf{Traduction.} « Que ce soient les jeunes ou les personnes âgées, tout le monde regarde les informations sur son téléphone. Que ce soit le matin ou le soir, mon père écoute la radio. Que ce soient les émissions en français ou en anglais, elle les comprend toutes. Qu'il fasse beau ou qu'il pleuve, il va acheter le journal à vélo. »

\textbf{Questions d'observation.}
\begin{enumerate}
\item Quels mots se répètent au début de chaque phrase ? (\zh{不管是……或者是……})
\item Quel mot apparaît dans la seconde partie de chaque phrase, juste avant le verbe ? (\zh{都})
\item Dans la dernière phrase, quel mot remplace \zh{或者是} ? (\zh{还是})
\end{enumerate}

\courssub{Définition et sens}

\begin{definition}[不管 bùguǎn]
\zh{不管} \emph{(bùguǎn)} est une conjonction de concession qui signifie « peu importe que… », « quelle que soit… ». Dans la structure \zh{不管是A，或者是B}, elle correspond au français « que ce soit A ou B », « aussi bien A que B ». Elle introduit deux (ou plusieurs) possibilités dont le choix est \textbf{indifférent} : la conclusion exprimée dans la seconde proposition reste vraie \textbf{dans tous les cas}.
\end{definition}

Le sens profond de la structure est donc double :
\begin{itemize}
\item \textbf{Concession avec choix multiples} : on énumère des alternatives (A, B, parfois C) ;
\item \textbf{Conclusion invariable} : le résultat ne dépend pas du choix. C'est ce que marque l'adverbe \zh{都} \emph{(dōu)} « tous, dans tous les cas » ou \zh{也} \emph{(yě)} « également ».
\end{itemize}

Comparaison avec le français : le français dit « que ce soit le matin \emph{ou} le soir » avec un simple « ou » ; le chinois, lui, ne peut pas se contenter de \zh{或者} : il faut encadrer les alternatives avec \zh{不管是} et conclure avec \zh{都/也}. La logique chinoise exige que l'on dise explicitement « peu importe (A ou B) → dans tous les cas (conclusion) ».

\courssub{Schéma général de la structure}

\begin{propriete}[Schéma de la structure 不管是……或者是……]
\begin{center}
不管是 + A，或者是 + B，(主语) + 都/也 + 谓语\\
\emph{Bùguǎn shì + A, huòzhě shì + B, (sujet) + dōu/yě + prédicat.}
\end{center}
\begin{itemize}
\item \zh{不管是} introduit la première possibilité ; \zh{是} peut être omis : \zh{不管白天晚上} est possible à l'oral.
\item \zh{或者是} introduit la seconde possibilité ; on peut ajouter une troisième : \zh{或者是C}.
\item Le prédicat de la proposition principale contient \textbf{obligatoirement} \zh{都} ou \zh{也}. Sans eux, la phrase est \textbf{agrammaticale}.
\end{itemize}
\end{propriete}

\begin{popfigure}
\begin{tikzpicture}[node distance=0.55cm, every node/.style={font=\small}]
\node[draw=popblue, fill=popblueL, rounded corners, align=center, minimum width=3.4cm, minimum height=1.1cm] (a) {不管是 + A\\ \emph{bùguǎn shì}\\ « que ce soit A »};
\node[draw=poporange, fill=poporangeL, rounded corners, align=center, minimum width=3.4cm, minimum height=1.1cm, right=of a] (b) {或者是 + B\\ \emph{huòzhě shì}\\ « ou que ce soit B »};
\node[draw=popgreen, fill=popgreenL, rounded corners, align=center, minimum width=3.6cm, minimum height=1.1cm, right=of b] (c) {都/也 + prédicat\\ \emph{dōu/yě}\\ « dans tous les cas »};
\draw[-{Stealth}, thick, popdark] (a) -- (b);
\draw[-{Stealth}, thick, popdark] (b) -- (c) node[midway, above, font=\footnotesize\itshape, text=poppurple]{conclusion identique};
\end{tikzpicture}
\legende{Chaîne logique de la structure 不管是……或者是…… : deux alternatives, une seule conclusion.}
\end{popfigure}

\courssub{Exemples gradués et commentés}

\begin{exemplebox}[Exemples de base]
\begin{enumerate}
\item \zh{不管是老师，或者是学生，都喜欢这个节目。}\\
\emph{Bùguǎn shì lǎoshī, huòzhě shì xuésheng, dōu xǐhuan zhège jiémù.}\\
« Que ce soient les professeurs ou les élèves, tous aiment cette émission. » (Le sujet est inclus dans les alternatives ; \zh{都} reprend l'ensemble.)
\item \zh{不管是白天，或者是晚上，他都在网上。}\\
\emph{Bùguǎn shì báitiān, huòzhě shì wǎnshang, tā dōu zài wǎngshang.}\\
« Que ce soit le jour ou la nuit, il est sur Internet. » (Ici le sujet \zh{他} est distinct et placé avant \zh{都}.)
\item \zh{不管是法语，或者是英语，他都会说。}\\
\emph{Bùguǎn shì Fǎyǔ, huòzhě shì Yīngyǔ, tā dōu huì shuō.}\\
« Qu'il s'agisse du français ou de l'anglais, il sait les parler. »
\end{enumerate}
\end{exemplebox}

\begin{exemplebox}[Exemples dans le contexte des médias (Module 5)]
\begin{enumerate}
\item \zh{不管是报纸，或者是网络，我们都应该核实消息。}\\
\emph{Bùguǎn shì bàozhǐ, huòzhě shì wǎngluò, wǒmen dōu yīnggāi héshí xiāoxi.}\\
« Que ce soit le journal ou Internet, nous devons vérifier les informations. »
\item \zh{不管是杜阿拉，或者是雅温得，人们都喜欢看足球比赛。}\\
\emph{Bùguǎn shì Dù'ālā, huòzhě shì Yǎwēndé, rénmen dōu xǐhuan kàn zúqiú bǐsài.}\\
« Que ce soit à Douala ou à Yaoundé, les gens aiment regarder les matchs de football. »
\item \zh{不管是晴天还是雨天，记者们都去现场采访。}\\
\emph{Bùguǎn shì qíngtiān háishi yǔtiān, jìzhěmen dōu qù xiànchǎng cǎifǎng.}\\
« Qu'il fasse beau ou qu'il pleuve, les journalistes vont interviewer sur le terrain. »
\end{enumerate}
\end{exemplebox}

\courssub{Variante avec 还是 : les alternatives de type question}

Quand les deux alternatives forment une opposition naturelle (beau temps / pluie, oui / non, aller / ne pas aller), le chinois préfère \zh{还是} \emph{(háishi)}, qui est le « ou » des questions, au lieu de \zh{或者}. Le schéma devient :

\begin{center}
不管是 + A，还是 + B，都/也 + 谓语
\end{center}

\begin{exemplebox}[Variante avec 还是]
\begin{enumerate}
\item \zh{不管是晴天还是雨天，他都骑自行车去买报纸。}\\
\emph{Bùguǎn shì qíngtiān háishi yǔtiān, tā dōu qí zìxíngchē qù mǎi bàozhǐ.}\\
« Qu'il fasse beau ou qu'il pleuve, il va acheter le journal à vélo. »
\item \zh{不管是你去还是我去，都一样。}\\
\emph{Bùguǎn shì nǐ qù háishi wǒ qù, dōu yíyàng.}\\
« Que ce soit toi ou moi qui y allions, c'est pareil. »
\end{enumerate}
\end{exemplebox}

\textbf{Nuance.} \zh{或者} est le « ou » des phrases affirmatives ; \zh{还是} est le « ou » des questions. Comme la structure \zh{不管} présente les alternatives comme des hypothèses ouvertes (proches d'une question : « est-ce A ? est-ce B ? »), \zh{还是} y est très naturel, surtout avec des verbes ou des adjectifs. Avec des noms simples (français/anglais, matin/soir), les deux sont acceptés.

\courssub{Emplois : habitudes et généralités}

La structure \zh{不管是……或者是……} sert principalement à exprimer :
\begin{itemize}
\item une \textbf{habitude} indépendante des circonstances : \zh{不管是工作日还是周末，他都六点起床。} « Que ce soit un jour de semaine ou le week-end, il se lève à six heures. » ;
\item une \textbf{généralité} : \zh{不管是大人还是小孩，都喜欢这个节目。} « Adultes ou enfants, tout le monde aime cette émission. » ;
\item une \textbf{condition indifférente} avant une consigne : \zh{不管是打电话还是发短信，都要先告诉老师。} « Qu'on téléphone ou qu'on envoie un SMS, il faut d'abord prévenir le professeur. »
\end{itemize}

\courssub{Complément utile pour l'examen : 不管 + question indirecte}

\begin{aretenir}[Extension : 不管 + mot interrogatif]
\zh{不管} peut aussi introduire une \textbf{question indirecte} avec un mot interrogatif (\zh{谁} qui, \zh{什么} quoi, \zh{哪儿} où, \zh{怎么} comment, \zh{多} combien). Le sens est « peu importe qui/quoi/où… ». La règle d'or reste identique : \zh{都} ou \zh{也} est obligatoire dans la conclusion.
\begin{itemize}
\item \zh{不管谁来，我们都欢迎。} \emph{Bùguǎn shéi lái, wǒmen dōu huānyíng.} « Peu importe qui vient, nous l'accueillons. »
\item \zh{不管多难，他也学汉语。} \emph{Bùguǎn duō nán, tā yě xué Hànyǔ.} « Si difficile que ce soit, il apprend le chinois. »
\end{itemize}
Cette extension est fréquente à l'examen : sachez la reconnaître.
\end{aretenir}

\courssub{Tableau récapitulatif}

\begin{center}
\begin{tabularx}{\linewidth}{|X|X|X|}
\hline
\rowcolor{popblueL} \textbf{Structure} & \textbf{Exemple (caractères + pinyin)} & \textbf{Traduction} \\
\hline
不管是 A，或者是 B，都 + V & \zh{不管是白天，或者是晚上，他都在网上。}\emph{Bùguǎn shì báitiān, huòzhě shì wǎnshang, tā dōu zài wǎngshang.} & Que ce soit le jour ou la nuit, il est en ligne. \\
\hline
不管是 A，还是 B，也 + V & \zh{不管是晴天还是雨天，他也去采访。}\emph{Bùguǎn shì qíngtiān háishi yǔtiān, tā yě qù cǎifǎng.} & Qu'il fasse beau ou qu'il pleuve, il va interviewer. \\
\hline
不管 + mot interrogatif，都 + V & \zh{不管谁来，都欢迎。}\emph{Bùguǎn shéi lái, dōu huānyíng.} & Peu importe qui vient, on l'accueille. \\
\hline
\end{tabularx}
\end{center}

\begin{center}
\begin{tabularx}{\linewidth}{|X|X|}
\hline
\rowcolor{popblueL} \textbf{Français} & \textbf{Chinois} \\
\hline
que ce soit A ou B & 不管是A，或者是B，都…… \\
\hline
qu'il fasse beau ou qu'il pleuve & 不管是晴天还是雨天 \\
\hline
peu importe qui / quoi / où & 不管谁 / 什么 / 哪儿 \\
\hline
dans tous les cas (marque obligatoire) & 都 / 也 \\
\hline
\end{tabularx}
\end{center}

\courssub{Exercice résolu et pièges des francophones}

\begin{exoresolu}[Compléter avec 都 ou 也]
Énoncé. Complétez les phrases avec \zh{都} ou \zh{也}.
\begin{enumerate}
\item \zh{不管是星期一，或者是星期五，他\_\_\_\_听广播。}
\item \zh{不管是中文课，还是法文课，她\_\_\_\_喜欢。}
\item \zh{不管谁打电话来，我\_\_\_\_接。}
\end{enumerate}
\tcblower
Solution détaillée.
\begin{enumerate}
\item \zh{都} : \zh{不管是星期一，或者是星期五，他都听广播。} « Que ce soit le lundi ou le vendredi, il écoute la radio. » \zh{都} est le choix par défaut pour « dans tous les cas ».
\item \zh{都} (ou \zh{也}) : \zh{不管是中文课，还是法文课，她都喜欢。} « Que ce soit le cours de chinois ou de français, elle les aime. »
\item \zh{都} : \zh{不管谁打电话来，我都接。} « Peu importe qui téléphone, je décroche. »
\end{enumerate}
Méthode : repérer d'abord \zh{不管} ; vérifier ensuite que le prédicat contient \zh{都/也} ; si le mot manque, la phrase est incomplète et doit être corrigée.
\end{exoresolu}

\begin{attention}[Erreurs fréquentes des francophones]
\begin{enumerate}
\item \textbf{Traduire « ou » par \zh{或者} seul, sans \zh{不管}.} Le français « le matin ou le soir, il écoute la radio » ne se traduit pas par \zh{早上或者晚上，他听广播} : cette phrase signifie simplement « le matin ou le soir (une fois), il écoute la radio » et perd toute la valeur concessive. Il faut : \zh{不管是早上，或者是晚上，他都听广播。}
\item \textbf{Oublier \zh{都} ou \zh{也}.} La phrase \zh{不管是老师，或者是学生，喜欢这个节目} est agrammaticale : sans \zh{都}, la conclusion n'est pas marquée comme universelle.
\item \textbf{Placer \zh{都} au mauvais endroit.} \zh{都} se place \emph{avant} le verbe, \emph{après} le sujet : \zh{他都在网上}, jamais \zh{都他在网上}.
\item \textbf{Confondre \zh{或者} et \zh{还是}.} Dans une question directe, on utilise \zh{还是} (\zh{你喝茶还是喝咖啡？}) ; dans la structure concessive, les deux sont possibles, mais \zh{还是} est préféré avec des alternatives verbales ou adjectivales.
\end{enumerate}
\end{attention}

\begin{exercice}[À vous de jouer]
\begin{enumerate}
\item Traduisez en chinois : « Que ce soit la télévision ou la radio, mon grand-père aime les deux. »
\item Traduisez en chinois : « Qu'il pleuve ou qu'il vente, les élèves vont tous à l'école. »
\item Complétez : \zh{不管是谁，\_\_\_\_要遵守规则。} (tout le monde doit respecter les règles)
\end{enumerate}
\end{exercice}

\begin{corrige}[Exercice « À vous de jouer »]
\begin{enumerate}
\item \zh{不管是电视，或者是收音机，我爷爷都喜欢。} \emph{Bùguǎn shì diànshì, huòzhě shì shōuyīnjī, wǒ yéye dōu xǐhuan.}
\item \zh{不管是下雨还是刮风，学生们都去上学。} \emph{Bùguǎn shì xiàyǔ háishi guāfēng, xuéshengmen dōu qù shàngxué.}
\item \zh{不管是谁，都要遵守规则。} \emph{Bùguǎn shì shéi, dōu yào zūnshǒu guīzé.}
\end{enumerate}
\end{corrige}

\begin{aretenir}[L'essentiel de la structure 1]
\begin{itemize}
\item Schéma : 不管是 + A，或者是 + B，(sujet) + 都/也 + prédicat.
\item Sens : « que ce soit A ou B », la conclusion reste vraie dans tous les cas.
\item \zh{都} ou \zh{也} est \textbf{obligatoire} dans la proposition principale.
\item Variante : \zh{还是} pour les alternatives de type question (\zh{晴天还是雨天}).
\item Extension d'examen : \zh{不管} + mot interrogatif (\zh{不管谁来，都……}).
\end{itemize}
\end{aretenir}

\coursec{Structure 2 : 是……的 (insistance sur une action passée)}

Dans la section précédente, nous avons appris à exprimer l'idée de « que ce soit… ou… » avec \zh{不管是……或者是……}. Nous restons dans le domaine de la mise en relief : le chinois possède une tournure très fréquente, la construction \cle{是……的}, qui permet d'insister sur les \textbf{circonstances} d'une action déjà accomplie (quand ? où ? comment ? avec qui ?). C'est l'équivalent fonctionnel du français « c'est… que / c'est… qui », mais avec une restriction essentielle : elle ne s'emploie que pour le \textbf{passé}.

\courssub{Observation guidée}

\begin{textebox}[À la radio de Yaoundé]
记者：这条新闻是昨天播出的。\\
Jìzhě : Zhè tiáo xīnwén shì zuótiān bōchū de.\\
« Journaliste : Cette nouvelle, c'est hier qu'elle a été diffusée. »\\
听众：我是在网上看到这个视频的，不是在电视上看到的。\\
Tīngzhòng : Wǒ shì zài wǎngshàng kàndào zhège shìpín de, bú shì zài diànshì shàng kàndào de.\\
« Auditeur : C'est sur Internet que j'ai vu cette vidéo, pas à la télévision. »\\
记者：你是怎么知道这个消息的？\\
Jìzhě : Nǐ shì zěnme zhīdào zhège xiāoxi de ?\\
« Journaliste : Comment as-tu appris cette nouvelle ? »\\
听众：是我朋友告诉我的。\\
Tīngzhòng : Shì wǒ péngyou gàosu wǒ de.\\
« Auditeur : C'est un ami qui me l'a dit. »
\end{textebox}

Observez : dans chaque phrase, l'action est \textbf{déjà terminée} (diffuser, voir, apprendre, dire), et le locuteur ne s'intéresse pas au fait lui-même mais à une \textbf{circonstance} : le moment (\zh{昨天}), le lieu ou le canal (\zh{在网上}), la manière (\zh{怎么}), l'agent (\zh{我朋友}). Remarquez aussi qu'aucune de ces phrases ne contient \zh{了}.

\courssub{Schéma général et fonction}

\begin{definition}[La construction 是……的]
La tournure \cle{是……的} met en relief le \textbf{temps}, le \textbf{lieu}, la \textbf{manière}, le \textbf{moyen} ou l'\textbf{agent} d'une action \textbf{déjà accomplie}. L'information nouvelle, celle que l'on veut souligner, se place entre \zh{是} et \zh{的}.
\end{definition}

\begin{propriete}[Schéma de la structure]
\begin{center}
\textbf{Sujet + 是 + (temps / lieu / manière / agent) + verbe + 的}
\end{center}
\begin{itemize}
\item \zh{是} se place \textbf{avant} la circonstance soulignée (jamais avant le verbe seul).
\item \zh{的} se place \textbf{en fin de phrase}, ou juste après le verbe si le complément d'objet est un nom (voir plus bas).
\item La structure ne s'emploie que pour des actions \textbf{passées} ; pour le présent ou le futur, le français « c'est… qui/que » n'a pas d'équivalent direct en chinois : on insiste autrement (intonation, adverbes comme \zh{就} ou \zh{才}).
\end{itemize}
\end{propriete}

\begin{exemplebox}[Exemples fondamentaux]
\zh{他是坐飞机去北京的。}\\
\emph{Tā shì zuò fēijī qù Běijīng de.}\\
« C'est en avion qu'il est allé à Pékin. » (moyen)\\[4pt]
\zh{我们是去年在杜阿拉认识这个记者的。}\\
\emph{Wǒmen shì qùnián zài Dù'ālā rènshi zhège jìzhě de.}\\
« C'est l'année dernière, à Douala, que nous avons connu ce journaliste. » (temps + lieu)\\[4pt]
\zh{这封信不是他写的。}\\
\emph{Zhè fēng xìn bú shì tā xiě de.}\\
« Ce n'est pas lui qui a écrit cette lettre. » (agent, négation)\\[4pt]
\zh{这条消息是广播里播的。}\\
\emph{Zhè tiáo xiāoxi shì guǎngbō lǐ bō de.}\\
« C'est à la radio que cette nouvelle a été diffusée. » (canal)
\end{exemplebox}

\courssub{La position de 的 : deux cas}

\begin{propriete}[Position de 的]
\begin{enumerate}
\item \textbf{Cas général} : \zh{的} se place à la fin de la phrase :\\
\zh{你是在哪儿买这份报纸的？} \emph{Nǐ shì zài nǎr mǎi zhè fèn bàozhǐ de ?} « Où as-tu acheté ce journal ? »
\item \textbf{Objet = nom} : on peut (et on préfère souvent) placer \zh{的} juste après le verbe, avant le complément d'objet nominal :\\
\zh{我是昨天在杜阿拉买的这本书。} \emph{Wǒ shì zuótiān zài Dù'ālā mǎi de zhè běn shū.} « C'est hier, à Douala, que j'ai acheté ce livre. »
\item \textbf{Objet = pronom} : \zh{的} reste obligatoirement en fin de phrase :\\
\zh{这本书是他送给我的。} \emph{Zhè běn shū shì tā sòng gěi wǒ de.} « C'est lui qui m'a offert ce livre. »
\end{enumerate}
\end{propriete}

\begin{attention}[Erreur fréquente : la place de 的]
Les francophones placent souvent \zh{的} au mauvais endroit ou l'oublient. Les deux phrases suivantes sont \textbf{correctes} (deux positions possibles pour un objet nominal) :\\
\zh{我是昨天买这本书的。} (Cas 1 : \zh{的} en fin de phrase)\\
\zh{我是昨天买的这本书。} (Cas 2 : \zh{的} après le verbe, avant l'objet)\\
\emph{Attention :} on ne doit \textbf{jamais} écrire deux \zh{的} (✗ \zh{我是昨天买的这本书的}) ni omettre \zh{的} (✗ \zh{我是昨天买这本书}).
\end{attention}

\courssub{Négation, question et omission de 是}

\begin{propriete}[Négation et interrogation]
\begin{itemize}
\item \textbf{Négation} : on nie avec \cle{不是……的}, jamais avec \zh{没} dans cette tournure.\\
\zh{我们不是坐火车去的，是坐汽车去的。} \emph{Wǒmen bú shì zuò huǒchē qù de, shì zuò qìchē qù de.} « Nous n'y sommes pas allés en train, mais en bus. »\\
✗ \zh{我没是昨天来的。} est impossible : la négation porte sur \zh{是}, pas sur le verbe.
\item \textbf{Question avec 吗} : \zh{你是昨天来的吗？} \emph{Nǐ shì zuótiān lái de ma ?} « C'est hier que tu es arrivé ? »
\item \textbf{Question avec mot interrogatif} : \zh{你是怎么知道这个消息的？} \emph{Nǐ shì zěnme zhīdào zhège xiāoxi de ?} « Comment as-tu appris cette nouvelle ? »
\item \textbf{Langue parlée} : à l'affirmative et dans les questions, \zh{是} peut être omis : \zh{你什么时候来的？} \emph{Nǐ shénme shíhou lái de ?} « Quand es-tu arrivé ? » À l'écrit et à l'examen, gardez \zh{是}.
\end{itemize}
\end{propriete}

\begin{exemplebox}[Dans le domaine des médias]
\zh{我们是在网上看到这个视频的。}\\
\emph{Wǒmen shì zài wǎngshàng kàndào zhège shìpín de.}\\
« C'est sur Internet que nous avons vu cette vidéo. »\\[4pt]
\zh{这个节目是上个星期录的。}\\
\emph{Zhège jiémù shì shàng ge xīngqī lù de.}\\
« C'est la semaine dernière que cette émission a été enregistrée. »\\[4pt]
\zh{这张照片不是我拍的，是我弟弟拍的。}\\
\emph{Zhè zhāng zhàopiàn bú shì wǒ pāi de, shì wǒ dìdi pāi de.}\\
« Ce n'est pas moi qui ai pris cette photo, c'est mon petit frère. »
\end{exemplebox}

\courssub{是……的 ou 了 ? Ne pas confondre}

\begin{propriete}[Distinction fondamentale]
\begin{itemize}
\item \cle{了} marque l'\textbf{accomplissement} d'une action : le fait lui-même est l'information nouvelle. \zh{他昨天来了。} « Il est venu hier. » (je vous informe qu'il est venu)
\item \cle{是……的} suppose le fait \textbf{déjà connu} et insiste sur ses \textbf{circonstances} : \zh{他是昨天来的。} « C'est hier qu'il est venu. » (on sait qu'il est venu ; je précise quand)
\item \textbf{Exception à retenir} : on ne cumule pas \zh{了} et \zh{是……的} pour le même fait : ✗ \zh{他是昨天来了的。}
\end{itemize}
\end{propriete}

\begin{center}
\begin{tabularx}{\linewidth}{|X|X|X|}
\hline
\rowcolor{popblueL} Structure & Exemple (caractères + pinyin) & Traduction \\
\hline
了 (fait accompli) & 他昨天来了。Tā zuótiān lái le. & Il est venu hier. \\
\hline
是……的 (circonstance) & 他是昨天来的。Tā shì zuótiān lái de. & C'est hier qu'il est venu. \\
\hline
Négation 不是……的 & 他不是昨天来的。Tā bú shì zuótiān lái de. & Ce n'est pas hier qu'il est venu. \\
\hline
Question & 他是怎么来的？Tā shì zěnme lái de ? & Comment est-il venu ? \\
\hline
\end{tabularx}
\end{center}

\begin{popfigure}
\begin{tikzpicture}[
node distance=0.9cm and 1.6cm,
startstop/.style={rectangle, rounded corners, draw=popblue, fill=popblueL, text width=4.6cm, align=center, font=\small},
decision/.style={diamond, draw=poporange, fill=poporangeL, text width=3.4cm, align=center, font=\small, inner sep=1pt},
result/.style={rectangle, rounded corners, draw=popgreen, fill=popgreenL, text width=4.2cm, align=center, font=\small},
arrow/.style={-{Stealth[length=2.5mm]}, thick, popdark}]
\node[startstop] (start) {Je parle d'une action passée};
\node[decision, below=of start] (q) {Je veux souligner une circonstance ? (quand, où, comment, qui)};
\node[result, below left=1.1cm and -0.4cm of q, align=center] (oui){Oui : Sujet + 是 + circonstance + verbe + 的\\ 他是坐飞机去的。};
\node[result, below right=1.1cm and -0.4cm of q, align=center] (non){Non : phrase simple avec 了\\ 他坐飞机去了。};
\draw[arrow] (start) -- (q);
\draw[arrow] (q) -| node[pos=0.25, above]{oui} (oui);
\draw[arrow] (q) -| node[pos=0.25, above]{non} (non);
\end{tikzpicture}
\legende{Organigramme de choix : 是……的 ou phrase simple avec 了}
\end{popfigure}

\courssub{Comparaison avec le français}

Le français dit « c'est \emph{en avion} qu'il est parti », « c'est \emph{hier} qu'elle est arrivée », « ce n'est pas \emph{lui} qui a écrit ». Le chinois rend exactement cette mise en relief, mais avec deux différences majeures : d'une part la tournure est limitée au passé ; d'autre part la particule \zh{的}, sans équivalent français, vient « refermer » la phrase comme une seconde parenthèse. Pensez à \zh{是……的} comme à un cadre : tout ce qui est souligné entre dans le cadre, et \zh{的} ferme le cadre à la fin.

\begin{methode}[Transformer une phrase simple en 是……的]
\begin{enumerate}
\item Partez de la phrase simple au passé : \zh{他昨天坐飞机去北京了。}
\item Identifiez la circonstance à souligner (ici le moyen : \zh{坐飞机}).
\item Retirez \zh{了} et placez \zh{是} juste avant cette circonstance : \zh{他是坐飞机……}
\item Refermez avec \zh{的} en fin de phrase : \zh{他是坐飞机去北京的。}
\item Vérifiez : négation éventuelle avec \zh{不是}, et \zh{的} avant un objet nominal si besoin.
\end{enumerate}
\end{methode}

\begin{exoresolu}[Transformation guidée]
Énoncé : transformez la phrase \zh{我去年在雅温得学了汉语。} (Wǒ qùnián zài Yǎwēndé xué le Hànyǔ.) en mettant en relief le lieu « à Yaoundé », puis mettez-la à la forme négative.
\tcblower
Solution détaillée :
\begin{enumerate}
\item Circonstance à souligner : le lieu \zh{在雅温得}.
\item On supprime \zh{了} et on place \zh{是} devant le lieu : \zh{我是在雅温得……}
\item On referme avec \zh{的} : \zh{我是在雅温得学的汉语。} \emph{Wǒ shì zài Yǎwēndé xué de Hànyǔ.} « C'est à Yaoundé que j'ai appris le chinois. » (\zh{的} placé après le verbe car l'objet \zh{汉语} est un nom.)
\item Négation : \zh{我不是在雅温得学的汉语，是在杜阿拉学的。} \emph{Wǒ bú shì zài Yǎwēndé xué de Hànyǔ, shì zài Dù'ālā xué de.} « Ce n'est pas à Yaoundé que j'ai appris le chinois, c'est à Douala. »
\end{enumerate}
\end{exoresolu}

\begin{attention}[Pièges typiques des francophones]
\begin{itemize}
\item Oublier \zh{是} à la négation : ✗ \zh{我没是昨天来的。} → \zh{我不是昨天来的。}
\item Cumuler \zh{了} et \zh{是……的} : ✗ \zh{他是昨天来了的。} → \zh{他是昨天来的。}
\item Employer la tournure pour le futur : ✗ \zh{我是明天去的。} (pour insister sur le futur, dites simplement \zh{我明天才去。} « Ce n'est que demain que j'y vais. »)
\item Placer \zh{是} directement devant le verbe sans circonstance : ✗ \zh{我是去的。} ne veut rien dire seul.
\end{itemize}
\end{attention}

\begin{exercice}[À vous]
Mettez les phrases suivantes à la forme \zh{是……的} en soulignant l'élément indiqué, puis traduisez : 1) \zh{他昨天晚上看了这个电视节目。} (quand) ; 2) \zh{我们在网上找到了这条新闻。} (où) ; 3) \zh{她用手机拍了这张照片。} (comment). Puis mettez la phrase 3 à la forme négative.
\end{exercice}

\begin{corrige}[Exercice « À vous »]
1) \zh{他是昨天晚上看的这个电视节目。} \emph{Tā shì zuótiān wǎnshang kàn de zhège diànshì jiémù.} « C'est hier soir qu'il a regardé cette émission. »\\
2) \zh{我们是在网上找到的这条新闻。} \emph{Wǒmen shì zài wǎngshàng zhǎodào de zhè tiáo xīnwén.} « C'est sur Internet que nous avons trouvé cette nouvelle. »\\
3) \zh{她是用手机拍的这张照片。} \emph{Tā shì yòng shǒujī pāi de zhè zhāng zhàopiàn.} « C'est avec son téléphone portable qu'elle a pris cette photo. »\\
Négation : \zh{她不是用手机拍的这张照片。} « Ce n'est pas avec son téléphone qu'elle a pris cette photo. »
\end{corrige}

\begin{aretenir}[L'essentiel de la structure 是……的]
\begin{itemize}
\item Schéma : \textbf{Sujet + 是 + circonstance + verbe + 的} ; actions passées uniquement.
\item \zh{的} en fin de phrase, ou après le verbe devant un objet nominal.
\item Négation : \zh{不是……的} ; jamais \zh{没}.
\item \zh{了} constate le fait ; \zh{是……的} souligne les circonstances ; on ne les cumule pas.
\item À l'oral, \zh{是} peut être omis (\zh{你什么时候来的？}), mais gardez-le à l'écrit.
\end{itemize}
\end{aretenir}

\coursec{Structure 3 : la phrase comparative en 比 (比字句)}

Dans la section précédente, nous avons vu comment la construction \zh{是……的} permet d'insister sur les circonstances d'une action passée (quand, où, comment). Nous quittons maintenant le passé pour une question que les journalistes se posent chaque jour : \emph{comparer}. Quel média est plus rapide ? Lequel est plus populaire ? En chinois, la comparaison repose sur une structure figée et très régulière : la \cle{phrase en 比} (\zh{比字句}, \emph{bǐ zì jù}). C'est l'une des structures les plus fréquentes de l'examen et de la langue médiatique.

\courssub{Observation guidée : comparer les médias}

\begin{textebox}[Les médias au Cameroun : qui gagne ?]
在中国，网络比报纸快得多。\\
\emph{Zài Zhōngguó, wǎngluò bǐ bàozhǐ kuài de duō.}\\
« En Chine, Internet est beaucoup plus rapide que les journaux. »\\[4pt]
在喀麦隆，广播没有电视受欢迎。\\
\emph{Zài Kāmàilóng, guǎngbō méiyǒu diànshì shòu huānyíng.}\\
« Au Cameroun, la radio n'est pas aussi populaire que la télévision. »\\[4pt]
手机比电脑方便一点儿。\\
\emph{Shǒujī bǐ diànnǎo fāngbiàn yìdiǎnr.}\\
« Le téléphone portable est un peu plus pratique que l'ordinateur. »\\[4pt]
我弟弟比我小三岁，但是他比我更喜欢看新闻。\\
\emph{Wǒ dìdi bǐ wǒ xiǎo sān suì, dànshì tā bǐ wǒ gèng xǐhuan kàn xīnwén.}\\
« Mon petit frère a trois ans de moins que moi, mais il aime regarder les informations plus que moi. »
\end{textebox}

Observez bien ces quatre phrases : dans chaque cas, on trouve l'ordre \textbf{A + 比 + B + adjectif}. Remarquez aussi deux choses essentielles : l'adjectif n'est \emph{jamais} précédé de \zh{很}, et le degré de la différence s'exprime \emph{après} l'adjectif (\zh{得多}, \zh{一点儿}, \zh{三岁}).

\courssub{Le schéma général : A + 比 + B + adjectif}

\begin{propriete}[La phrase en 比 : forme de base]
\textbf{Schéma :} \cle{A} + \zh{比} + \cle{B} + \textbf{adjectif}\\
A = l'élément comparé (le sujet) ; B = la référence ; l'adjectif exprime la qualité comparée.\\
\textbf{Emploi :} affirmer que A dépasse B dans une qualité.\\
\textbf{Point capital :} l'ordre est \textbf{figé} et \textbf{inverse du français}. En français on dit « A est plus grand \emph{que} B » ; en chinois, la référence B vient \emph{avant} l'adjectif : \zh{A比B大} (littéralement « A — comparé à B — grand »).
\end{propriete}

\begin{exemplebox}[Exemples de base]
\zh{手机比电脑方便。} \emph{Shǒujī bǐ diànnǎo fāngbiàn.} « Le portable est plus pratique que l'ordinateur. »\\
\zh{今天的作业比昨天的多。} \emph{Jīntiān de zuòyè bǐ zuótiān de duō.} « Les devoirs d'aujourd'hui sont plus nombreux que ceux d'hier. »\\
\zh{雅温得比杜阿拉凉快。} \emph{Yǎwēndé bǐ Dù'ālā liángkuai.} « Yaoundé est plus fraîche que Douala. »\\
\zh{他比我小三岁。} \emph{Tā bǐ wǒ xiǎo sān suì.} « Il a trois ans de moins que moi. »
\end{exemplebox}

\begin{attention}[Erreur n° 1 des francophones : calquer l'ordre français]
Le français place l'adjectif \emph{avant} « que » : « plus grand \textbf{que} lui ». Le calque donne des phrases impossibles :\\
✗ \zh{大比他} (calque de « plus grand que lui »)\\
✗ \zh{他大比我}\\
✓ \zh{他比我大。} « Il est plus grand (âgé) que moi. »\\
Retenez le réflexe : dès que vous entendez « que » dans une comparaison, placez la référence \emph{juste après} \zh{比} et \emph{avant} l'adjectif.
\end{attention}

\begin{popfigure}
\begin{tikzpicture}[scale=0.95, every node/.style={font=\small}]
% Ordre correct
\node[draw=popgreen, fill=popgreenL, rounded corners, minimum width=1.6cm, minimum height=0.9cm] (A) at (0,0) {A};
\node[draw=popblue, fill=popblueL, rounded corners, minimum width=1.1cm, minimum height=0.9cm] (bi) at (2.2,0) {比};
\node[draw=popgreen, fill=popgreenL, rounded corners, minimum width=1.6cm, minimum height=0.9cm] (B) at (4.2,0) {B};
\node[draw=poporange, fill=poporangeL, rounded corners, minimum width=2.2cm, minimum height=0.9cm] (adj) at (6.8,0) {adjectif};
\draw[-{Stealth}, very thick, popblue] (A) -- (bi);
\draw[-{Stealth}, very thick, popblue] (bi) -- (B);
\draw[-{Stealth}, very thick, popblue] (B) -- (adj);
\node[above=0.15cm of bi, text=popdark] {ordre figé};
\node[below=0.3cm of B, text=popgreen, align=center] {手机 比 电脑 方便};
% Ordre français interdit
\node[draw=popink, fill=poppinkL, rounded corners, minimum width=2.6cm, minimum height=0.9cm] (fr) at (10.8,0) {A + adj + que B};
\draw[popink, line width=2pt] (9.3,0.65) -- (12.3,-0.65);
\draw[popink, line width=2pt] (9.3,-0.65) -- (12.3,0.65);
\node[below=0.3cm of fr, text=popink, align=center] {ordre français\\interdit en chinois};
\end{tikzpicture}
\legende{La phrase en 比 : A précède 比, la référence B suit 比, l'adjectif ferme la phrase.}
\end{popfigure}

\courssub{Exprimer le degré de la différence}

\begin{propriete}[Dans la 比字句, l'adjectif ne prend jamais 很]
Après \zh{比}, l'adjectif est déjà « comparatif » par nature : ajouter \zh{很} ou \zh{非常} est une faute.\\
✗ \zh{他比我}很\zh{高。} \quad ✗ \zh{网络比报纸}非常\zh{快。}\\
Le \textbf{degré} s'exprime autrement :
\begin{itemize}
\item \zh{比……得多} / \zh{比……多了} : « beaucoup plus… que » ;
\item \zh{比……一点儿} / \zh{比……一些} : « un peu plus… que » ;
\item un \textbf{complément de quantité} après l'adjectif : \zh{高两厘米} (plus grand de deux centimètres), \zh{早十分钟} (dix minutes plus tôt), \zh{小三岁} (trois ans de moins) ;
\item \zh{更} ou \zh{还} \emph{avant} l'adjectif pour insister : \zh{他比我\cle{更}高。} « Il est encore plus grand que moi. »
\end{itemize}
\end{propriete}

\begin{exemplebox}[Le degré en contexte médias]
\zh{网络比报纸快得多。} \emph{Wǎngluò bǐ bàozhǐ kuài de duō.} « Internet est beaucoup plus rapide que les journaux. »\\
\zh{手机比电脑方便一点儿。} \emph{Shǒujī bǐ diànnǎo fāngbiàn yìdiǎnr.} « Le portable est un peu plus pratique que l'ordinateur. »\\
\zh{我比他早来十分钟。} \emph{Wǒ bǐ tā zǎo lái shí fēnzhōng.} « Je suis arrivé dix minutes plus tôt que lui. »\\
\zh{他比我更喜欢看新闻。} \emph{Tā bǐ wǒ gèng xǐhuan kàn xīnwén.} « Il aime regarder les informations plus que moi. »
\end{exemplebox}

Notez le dernier exemple : la comparaison fonctionne aussi avec des \textbf{verbes} (aimer, préférer, comprendre). Le schéma devient A + \zh{比} + B + (\zh{更}) + verbe psychologique ; et avec un verbe d'action, la quantité se place \emph{après} le verbe : \zh{我比他早来十分钟。}

\courssub{La négation : 不比 ou 没有 ?}

\begin{definition}[\zh{没有} dans la comparaison]
\cle{A 没有 B + adjectif} signifie « A n'atteint pas B en… », c'est-à-dire « A est \textbf{moins}… que B » (infériorité réelle).\\
\zh{广播没有电视受欢迎。} \emph{Guǎngbō méiyǒu diànshì shòu huānyíng.} « La radio est moins populaire que la télévision. »\\
C'est différent de \cle{A 不比 B + adjectif}, qui est un \textbf{simple constat} : « A n'est pas plus… que B » (les deux peuvent être égaux, ou A légèrement inférieur, sans idée d'infériorité marquée).
\end{definition}

\begin{exemplebox}[Nuance 不比 / 没有]
\zh{我不比他忙。} \emph{Wǒ bù bǐ tā máng.} « Je ne suis pas plus occupé que lui. » (nous sommes à peu près pareils)\\
\zh{我没有他忙。} \emph{Wǒ méiyǒu tā máng.} « Je suis moins occupé que lui. » (il est clairement plus occupé)\\
\zh{这份报纸不比那份贵。} \emph{Zhè fèn bàozhǐ bù bǐ nà fèn guì.} « Ce journal n'est pas plus cher que celui-là. »\\
\zh{我的手机没有他的新。} \emph{Wǒ de shǒujī méiyǒu tā de xīn.} « Mon portable est moins récent que le sien. »
\end{exemplebox}

\courssub{Compléments utiles à l'examen : 越来越 et 一年比一年}

Deux tournures proches de la comparaison tombent souvent à l'examen :

\begin{propriete}[越来越 + adjectif et 一年比一年]
\begin{itemize}
\item \cle{越来越 + adjectif/verbe} : « de plus en plus… ». \zh{学中文的人越来越多。} \emph{Xué Zhōngwén de rén yuè lái yuè duō.} « De plus en plus de gens apprennent le chinois. » (Jamais \zh{很} après \zh{越来越} !)
\item \cle{一年比一年 + adjectif} : « d'année en année, de plus en plus… ». \zh{喀麦隆的互联网用户一年比一年多。} \emph{Kāmàilóng de hùliánwǎng yònghù yì nián bǐ yì nián duō.} « Au Cameroun, les internautes sont de plus en plus nombreux d'année en année. »
\end{itemize}
\end{propriete}

\courssub{Tableau récapitulatif}

\begin{center}
\begin{tabularx}{\linewidth}{|X|X|X|}
\hline
\rowcolor{popblueL} Structure & Exemple (caractères + pinyin) & Traduction \\
\hline
A 比 B + adj & 手机比电脑方便。Shǒujī bǐ diànnǎo fāngbiàn. & Le portable est plus pratique que l'ordinateur. \\
\hline
A 比 B + adj + 得多/多了 & 网络比报纸快得多。Wǎngluò bǐ bàozhǐ kuài de duō. & Internet est beaucoup plus rapide que les journaux. \\
\hline
A 比 B + adj + 一点儿 & 今天比昨天冷一点儿。Jīntiān bǐ zuótiān lěng yìdiǎnr. & Il fait un peu plus froid qu'hier. \\
\hline
A 比 B + adj + quantité & 他比我小三岁。Tā bǐ wǒ xiǎo sān suì. & Il a trois ans de moins que moi. \\
\hline
A 比 B + 更 + verbe & 他比我更喜欢看新闻。Tā bǐ wǒ gèng xǐhuan kàn xīnwén. & Il aime les infos plus que moi. \\
\hline
A 没有 B + adj & 广播没有电视受欢迎。Guǎngbō méiyǒu diànshì shòu huānyíng. & La radio est moins populaire que la télé. \\
\hline
A 不比 B + adj & 我不比他忙。Wǒ bù bǐ tā máng. & Je ne suis pas plus occupé que lui. \\
\hline
越来越 + adj & 新闻越来越有意思。Xīnwén yuè lái yuè yǒu yìsi. & Les infos sont de plus en plus intéressantes. \\
\hline
一年比一年 + adj & 用户一年比一年多。Yònghù yì nián bǐ yì nián duō. & Les usagers sont de plus en plus nombreux chaque année. \\
\hline
\end{tabularx}
\end{center}

\begin{center}
\begin{tabularx}{\linewidth}{|X|X|X|X|}
\hline
\rowcolor{popblueL} Caractères & Pinyin & Nature & Traduction \\
\hline
比较 & bǐjiào & verbe / adverbe & comparer ; assez, relativement \\
\hline
方便 & fāngbiàn & adjectif & pratique, commode \\
\hline
受欢迎 & shòu huānyíng & expression verbale & populaire, apprécié \\
\hline
报纸 & bàozhǐ & nom & journal (presse écrite) \\
\hline
广播 & guǎngbō & nom & radio (diffusion) \\
\hline
用户 & yònghù & nom & usager, utilisateur \\
\hline
差 & chà & adjectif & inférieur, mauvais \\
\hline
\end{tabularx}
\end{center}

\courssub{Exercice résolu et pièges}

\begin{exoresolu}[Traduisez en chinois]
1. « Douala est beaucoup plus chaude que Buea. »\\
2. « Mon frère est arrivé cinq minutes plus tôt que moi. »\\
3. « Les journaux sont moins rapides qu'Internet. »
\tcblower
\textbf{Solution détaillée :}\\
1. La différence est grande → \zh{得多} : \zh{杜阿拉比布埃亚热得多。} \emph{Dù'ālā bǐ Bù'āiyà rè de duō.} (Pas de \zh{很} !)\\
2. Complément de quantité après le verbe : \zh{我哥哥比我早来五分钟。} \emph{Wǒ gēge bǐ wǒ zǎo lái wǔ fēnzhōng.}\\
3. Infériorité → \zh{没有} : \zh{报纸没有网络快。} \emph{Bàozhǐ méiyǒu wǎngluò kuài.}
\end{exoresolu}

\begin{attention}[Les trois pièges de la 比字句]
\begin{enumerate}
\item \textbf{Ordre français} : ✗ \zh{快比报纸} → ✓ \zh{比报纸快}.
\item \textbf{很 interdit} : ✗ \zh{他比我}很\zh{高} → ✓ \zh{他比我高} ou \zh{他比我高得多}.
\item \textbf{Place du complément de quantité} : il vient \emph{après} l'adjectif ou le verbe : ✓ \zh{比我小三岁}, ✓ \zh{比我早来十分钟} ; ne le placez jamais entre \zh{比} et B.
\end{enumerate}
\end{attention}

\begin{exercice}[À vous de comparer]
1. Traduisez : « La télévision est un peu plus chère que la radio. » ; « Je ne suis pas plus grand que lui. » ; « Les élèves de Terminale sont de plus en plus nombreux. »\\
2. Corrigez : \zh{他比我很大。} / \zh{网络快比报纸。}\\
3. Rédigez trois phrases en \zh{比} comparant deux médias que vous utilisez à Yaoundé ou à Douala.
\end{exercice}

\begin{corrige}[Exercice « À vous de comparer »]
1. \zh{电视比广播贵一点儿。} \emph{Diànshì bǐ guǎngbō guì yìdiǎnr.} ; \zh{我没有他高。} / \zh{我不比他高。} \emph{Wǒ méiyǒu tā gāo. / Wǒ bù bǐ tā gāo.} ; \zh{毕业班的学生越来越多。} \emph{Bìyèbān de xuésheng yuè lái yuè duō.}\\
2. \zh{他比我很大} → \zh{他比我大。} (supprimer \zh{很}) ; \zh{网络快比报纸} → \zh{网络比报纸快。} (remettre \zh{比} + référence avant l'adjectif).\\
3. Exemple : \zh{手机比电视方便，但是电视比手机有意思。} \emph{Shǒujī bǐ diànshì fāngbiàn, dànshì diànshì bǐ shǒujī yǒu yìsi.} « Le portable est plus pratique que la télé, mais la télé est plus intéressante que le portable. »
\end{corrige}

\coursec{Comparaison avec le français et pièges des francophones}

Après avoir analysé en détail la structure comparative en \cle{比} (section précédente), nous devons maintenant prendre du recul pour confronter nos trois structures du jour — \cle{不管是……或者是……}, \cle{是……的} et \cle{比字句} — au fonctionnement de la langue française. C'est à cette frontière que se nichent les erreurs les plus tenaces, celles qui coûtent des points au Baccalauréat même quand le vocabulaire est connu. Observons, comparons, et entraînons-nous à déjouer les pièges.

\courssub{« Que ce soit… ou… » : l'obligation du corréla \texorpdfstring{\cle{都/也}}{dans le chinois}}

En français, la concession alternative (« que ce soit X ou Y ») se suffit à elle-même : la principale suit directement, sans mot-outil obligatoire.
\begin{exemplebox}[Français : pas de mot-outil final]
Que ce soit \textbf{la pluie} ou \textbf{le vent}, le match \textbf{aura lieu}. \\
Si c'est \textbf{Paul} ou \textbf{Marie} qui vient, \textbf{on mangera} bien.
\end{exemplebox}

En chinois, la structure \cle{不管是……或者是……} (ou \cle{无论是……还是……}) \textbf{exige} un adverbe corréla dans la principale : \cle{都} (tous, tous les deux) ou \cle{也} (aussi). L'omission de \cle{都/也} est la faute n°1 des francophones.

\begin{propriete}[Règle d'or : le corréla 都/也 est obligatoire]
\textbf{Schéma :} 不管是 A 或者是 B，主语 \textbf{都/也} V / Adj. \\
\textbf{Emploi :} Exprime une concession universelle : quelle que soit l'alternative, le résultat est le même. \\
\textbf{Nuance :} \cle{都} insiste sur la totalité (tous les cas) ; \cle{也} insiste sur l'addition (aussi bien dans un cas que dans l'autre).
\end{propriete}

\begin{center}
\begin{tabularx}{\linewidth}{|X|X|X|X|}
\hline
\rowcolor{popblueL} \textbf{Français} & \textbf{Chinois (Caractères)} & \textbf{Pinyin} & \textbf{Structure} \\
\hline
Que ce soit le vent ou la pluie, je vais y aller. & 不管是刮风，或者是下雨，我\textbf{都}去。 & Bùguǎn shì guāfēng, huòzhě shì xiàyǔ, wǒ \textbf{dōu} qù. & \cle{都} présent \\
\hline
Que ce soit toi ou moi, c'est pareil. & 不管是你，或者是我，\textbf{都}一样。 & Bùguǎn shì nǐ, huòzhě shì wǒ, \textbf{dōu} yīyàng. & \cle{都} présent \\
\hline
Qu'il soit grand ou petit, il est mignon. & 不管是大，或者是小，它\textbf{都}可爱。 & Bùguǎn shì dà, huòzhě shì xiǎo, tā \textbf{dōu} kě'ài. & \cle{都} présent \\
\hline
\end{tabularx}
\end{center}

\begin{attention}[Piège n°1 : L'oubli de \texorpdfstring{\cle{都}}{dou} / \texorpdfstring{\cle{也}}{ye}]
\begin{itemize}
    \item \textbf{Faux :} \zh{不管是下雨或者是刮风，他去学校。} (Bùguǎn shì xiàyǔ huòzhě shì guāfēng, tā qù xuéxiào.) — \emph{Il manque le corréla.}
    \item \textbf{Vrai :} \zh{不管是下雨或者是刮风，他\textbf{都}去学校。} (Tā \textbf{dōu} qù xuéxiào.) — « Qu'il pleuve ou qu'il vente, il va à l'école (dans les deux cas). »
    \item \textbf{Astuce :} À l'oral, \cle{都} se place \textbf{avant} le verbe (ou l'adjectif). Ne le mettez jamais en fin de phrase.
\end{itemize}
\end{attention}

\begin{savaistu}[Au Bac : le test du corréla]
Dans les items de grammaire (QCM ou transformation), si vous voyez \cle{不管是} ou \cle{无论是} sans \cle{都/也} dans la proposition, c'est \textbf{systématiquement faux}. Le correcteur attend la paire indissociable : \cle{不管是……或者是……，都/也……}.
\end{savaistu}

\courssub{« C'est… que/qui » (insistance) : le piège du \texorpdfstring{\cle{是……的}}{shi...de} et l'interdiction du \texorpdfstring{\cle{了}}{le}}

Le français utilise la cleft sentence (« C'est hier \textbf{qu}'il est arrivé ») pour mettre l'accent sur un élément (temps, lieu, manière, but). Le chinois utilise la construction \cle{是……的} \textbf{uniquement pour des actions passées, déjà réalisées}. C'est une « loupe » sur les circonstances d'un fait accompli.

\begin{propriete}[Structure \texorpdfstring{\cle{是……的}}{shi...de} : insistance sur le passé accompli]
\textbf{Schéma :} Sujet + \textbf{是} + \emph{Circonstance (Temps / Lieu / Manière / But / Compagnon)} + Verbe + \textbf{的}. \\
\textbf{Condition sine qua non :} L'action est \textbf{terminée} (aspect parfait). \\
\textbf{Interdiction formelle :} Ne jamais mettre \cle{了} après le verbe dans cette structure. \cle{的} absorbe la valeur d'accompli.
\end{propriete}

\begin{attention}[Piège n°2 : Le cumul interdit \texorpdfstring{\cle{是……了}}{shi...le} / \texorpdfstring{\cle{是……的……了}}{shi...de...le}]
\begin{itemize}
    \item \textbf{Faux :} \zh{我是昨天来\textbf{了}。} (Wǒ shì zuótiān lái \textbf{le}.) — \emph{Double marquage du passé.}
    \item \textbf{Faux :} \zh{我是昨天来\textbf{的了}。} — \emph{Inconcevable.}
    \item \textbf{Vrai :} \zh{我是昨天来\textbf{的}。} (Wǒ shì zuótiān lái \textbf{de}.) — « C'est \textbf{hier} que je suis venu. »
    \item \textbf{Vrai (sans insistance) :} \zh{我昨天来\textbf{了}。} (Wǒ zuótiān lái \textbf{le}.) — Simple narration : « Je suis venu hier. »
\end{itemize}
\end{attention}

\begin{exemplebox}[Comparaison systématique : Français « C'est... que » vs Chinois \texorpdfstring{\cle{是……的}}{shi...de}]
\begin{center}
\begin{tabularx}{\linewidth}{|X|X|X|}
\hline
\rowcolor{popblueL} \textbf{Français (Insistance)} & \textbf{Chinois (\cle{是……的})} & \textbf{Pinyin} \\
\hline
C'est \textbf{hier} qu'il est arrivé. & 他\textbf{是昨天}到\textbf{的}。 & Tā \textbf{shì zuótiān} dào \textbf{de}. \\
\hline
C'est \textbf{à Yaoundé} qu'on s'est vus. & 我们\textbf{是在雅温得}见面\textbf{的}。 & Wǒmen \textbf{shì zài Yǎwēndé} jiànmiàn \textbf{de}. \\
\hline
C'est \textbf{en taxi} qu'il est rentré. & 他\textbf{是坐出租车}回\textbf{的}。 & Tā \textbf{shì zuò chūzūchē} huí \textbf{de}. \\
\hline
C'est \textbf{pour acheter un livre} qu'il est allé en ville. & 他\textbf{是为买书}去城里\textbf{的}。 & Tā \textbf{shì wèi mǎi shū} qù chénglǐ \textbf{de}. \\
\hline
\end{tabularx}
\end{center}
\end{exemplebox}

\begin{methode}[Comment transformer une phrase française « C'est... que » en chinois]
\begin{enumerate}
    \item \textbf{Vérifier} : L'action est-elle passée et terminée ? (Si non $\rightarrow$ pas de \cle{是……的}, utiliser \cle{在} pour le progressif ou structure simple pour le futur/habitude).
    \item \textbf{Identifier} l'élément mis en relief (temps, lieu, manière, but...).
    \item \textbf{Placer} \cle{是} \textbf{avant} cet élément.
    \item \textbf{Placer} \cle{的} \textbf{à la toute fin} de la phrase (après le verbe, après le complément d'objet s'il y a lieu).
    \item \textbf{Supprimer} tout \cle{了}, \cle{过}, \cle{着}.
\end{enumerate}
\end{methode}

\courssub{Le comparatif « plus... que » : l'inversion totale de l'ordre des termes}

C'est le choc frontal le plus violent entre les deux langues. En français : \textbf{A + plus + Adj + que + B}. En chinois : \textbf{A + 比 + B + Adj (+ degré)}. Le terme de comparaison (B) glisse \textbf{avant} l'adjectif.

\begin{propriete}[Structure \texorpdfstring{\cle{比}}{bi} : A compare à B]
\textbf{Schéma :} A + \textbf{比} + B + \textbf{Adjectif} (+ \cle{得多} / \cle{一点儿} / \cle{多了} / quantité). \\
\textbf{Règle d'or :} \textbf{JAMAIS} de \cle{很} (hěn) devant l'adjectif dans une phrase en \cle{比}. L'adjectif est nu (ou suivi d'un complément de degré).
\end{propriete}

\begin{center}
\begin{tabularx}{\linewidth}{|X|X|X|}
\hline
\rowcolor{popblueL} \textbf{Français (Ordre : A > Adj > B)} & \textbf{Chinois (Ordre : A > 比 > B > Adj)} & \textbf{Pinyin} \\
\hline
Internet est \textbf{plus rapide que} le journal. & 网络 \textbf{比} 报纸 \textbf{快得多}。 & Wǎngluò \textbf{bǐ} bàozhǐ \textbf{kuài dé duō}. \\
\hline
Yaoundé est \textbf{plus grand que} Buea. & 雅温得 \textbf{比} 布埃亚 \textbf{大}。 & Yǎwēndé \textbf{bǐ} Bù'āiyà \textbf{dà}. \\
\hline
Il est \textbf{trois ans plus âgé que} moi. & 他 \textbf{比} 我 \textbf{大三岁}。 & Tā \textbf{bǐ} wǒ \textbf{dà sān suì}. \\
\hline
Ce téléphone est \textbf{un peu plus cher}. & 这手机 \textbf{比} 那个 \textbf{贵一点儿}。 & Zhè shǒujī \textbf{bǐ} nàge \textbf{guì yīdiǎnr}. \\
\hline
\end{tabularx}
\end{center}

\begin{attention}[Piège n°3 : L'ordre français \texorpdfstring{\cle{Adj + 比 + B}}{Adj bi B} (Ex: *网络很快比报纸)]
\begin{itemize}
    \item \textbf{Faux :} \zh{网络很快比报纸。} (Wǎngluò hěn kuài bǐ bàozhǐ.) — \emph{Ordre français, présence de \cle{很}.}
    \item \textbf{Vrai :} \zh{网络\textbf{比}报纸\textbf{快得多}。} (Wǎngluò \textbf{bǐ} bàozhǐ \textbf{kuài dé duō}.)
\end{itemize}
\end{attention}

\begin{attention}[Piège n°4 : L'intrus \texorpdfstring{\cle{很}}{hen} (Ex: *他比我\cle{很}高)]
\begin{itemize}
    \item \textbf{Faux :} \zh{他比我\textbf{很}高。} (Tā bǐ wǒ \textbf{hěn} gāo.) — \cle{很} est incompatible avec \cle{比}.
    \item \textbf{Vrai :} \zh{他比我\textbf{高}。} (Tā bǐ wǒ \textbf{gāo}.) — Nuance neutre.
    \item \textbf{Vrai (précisé) :} \zh{他比我\textbf{高五厘米}。} / \zh{他比我\textbf{高得多}。} / \zh{他比我\textbf{高一点儿}。}
\end{itemize}
\end{attention}

\begin{methode}[Check-list « Détective » avant de valider une phrase en \texorpdfstring{\cle{比}}{bi}]
\begin{enumerate}
    \item \textbf{Sujet A} est-il bien \textbf{devant} \cle{比} ? (Celui qui « gagne » la comparaison).
    \item \textbf{B} (le terme de référence) est-il bien \textbf{entre} \cle{比} et l'adjectif ?
    \item \textbf{Adjectif} est-il \textbf{nu} (sans \cle{很}, \cle{非常}, \cle{真}) ?
    \item \textbf{Degré} est-il exprimé ? (\cle{得多}, \cle{一点儿}, \cle{多了}, quantité : \cle{两岁}, \cle{五公斤}...). Si pas de degré $\rightarrow$ nuance neutre acceptée, mais degré préféré à l'écrit.
\end{enumerate}
\end{methode}

\courssub{Synthèse visuelle : Tableau de transposition des trois structures}

\begin{popfigure}
\begin{tikzpicture}[
    node distance=1.2cm and 2cm,
    frbox/.style={rectangle, draw=popblue, thick, fill=popblueL, rounded corners, minimum width=3.5cm, align=center, font=\small},
    zhbox/.style={rectangle, draw=poporange, thick, fill=poporangeL, rounded corners, minimum width=4cm, align=center, font=\small},
    arrow/.style={-Stealth, thick, popdark, shorten >=2pt, shorten <=2pt},
    labelstyle/.style={font=\tiny\itshape, text=popmuted}
]
% Ligne 1 : 不管是
\node[frbox, align=center] (fr1){Français : « Que ce soit X ou Y,\\ le résultat est le même. »\\ \textit{(Pas de mot final obligatoire)}};
\node[zhbox, right=of fr1, align=center] (zh1){Chinois : \textbf{不管是 X 或者是 Y，}\\ \textbf{主语 都/也 V/Adj.}\\ \textit{(Corréla 都/也 OBLIGATOIRE)}};
\draw[arrow] (fr1) -- node[above, labelstyle] {Ajouter 都/也} (zh1);

% Ligne 2 : 是...的
\node[frbox, below=of fr1, align=center] (fr2){Français : « C'est [circonstance]\\ que/qui... » (Insistance)\\ \textit{(Action passée ou présente)}};
\node[zhbox, right=of fr2, align=center] (zh2){Chinois : \textbf{是 [Circonstance] V 的}\\ \textit{(Action PASSÉE uniquement,}\\\textit{PAS de 了, 的 en FIN)}};
\draw[arrow] (fr2) -- node[above, labelstyle] {Vérifier passé + déplacer 的} (zh2);

% Ligne 3 : 比
\node[frbox, below=of fr2, align=center] (fr3){Français : « A plus Adj que B »\\ \textit{(Ordre : A - Adj - B)}};
\node[zhbox, right=of fr3, align=center] (zh3){Chinois : \textbf{A 比 B Adj (Degré)}\\ \textit{(Ordre : A - 比 - B - Adj,}\\\textit{PAS de 很, Degré requis)}};
\draw[arrow] (fr3) -- node[above, labelstyle, align=center]{Inverser B et Adj,\\ supprimer 很} (zh3);
\end{tikzpicture}
\legende{Transposition Français $\rightarrow$ Chinois : les trois pièges majeurs de la leçon 35.}
\end{popfigure}

\courssub{Atelier « Détective » : Corriger les 5 erreurs types}

Appliquez la check-list ci-dessus. Chaque phrase contient \textbf{une seule erreur majeure} (parfois deux liées). Corrigez-la et justifiez en une phrase.

\begin{exoresolu}[Atelier Détective — 5 phrases à corriger]
\textbf{Énoncé :} Corrigez les phrases suivantes et expliquez l'erreur (ordre, mot-outil, aspect, degré).

\begin{enumerate}
    \item \zh{不管是刮风或者是下雨，比赛照常进行。}
    \item \zh{他是2020年来喀麦隆了的。}
    \item \zh{这个手机比那个手机很贵。}
    \item \zh{不管你去不去，我都去。} \textit{(Variante : \cle{不管} + verbe, structure proche)}
    \item \zh{网络比传媒快。} \textit{(Contexte : comparer la vitesse de diffusion)}
\end{enumerate}

\tcblower

\textbf{Solution détaillée :}

\begin{enumerate}
    \item \textbf{Erreur :} Manque du corréla \cle{都/也}. \\
    \textbf{Correction :} \zh{不管是刮风，或者是下雨，比赛\textbf{都}照常进行。} (Bùguǎn shì guāfēng, huòzhě shì xiàyǔ, bǐsài \textbf{dōu} zhàocháng jìnxíng.) \\
    \textbf{Justification :} Structure \cle{不管是……或者是……} impose \cle{都} (ou \cle{也}) avant le verbe de la principale.
    
    \item \textbf{Erreur :} Double marquage du passé (\cle{了} + \cle{的}) + ordre incorrect (\cle{的} avant \cle{了} impossible). \\
    \textbf{Correction :} \zh{他\textbf{是}2020年来喀麦隆\textbf{的}。} (Tā \textbf{shì} 2020 nián lái Kāmàilóng \textbf{de}.) \\
    \textbf{Justification :} \cle{是……的} insiste sur le temps passé ; \cle{了} est interdit, \cle{的} se place en position finale.
    
    \item \textbf{Erreur :} Présence interdite de \cle{很} dans une phrase comparative \cle{比}. \\
    \textbf{Correction :} \zh{这个手机比那个手机\textbf{贵得多} / \textbf{贵一点儿} / \textbf{贵五百元}。} (Zhège shǒujī bǐ nàge shǒujī \textbf{guì dé duō / guì yīdiǎnr / guì wǔbǎi yuán}.) \\
    \textbf{Justification :} \cle{比} exclut \cle{很} ; l'adjectif doit être suivi d'un complément de degré (\cle{得多}, \cle{一点儿}, quantité).
    
    \item \textbf{Erreur :} Structure \cle{不管} + Verbe (conditionnelle) $\neq$ \cle{不管是} + Nom (concession alternative). Ici, « Que tu ailles ou non » $\rightarrow$ \cle{无论你去不去} ou \cle{不管你去不去}. La phrase \zh{不管你去不去，我都去} est \textbf{correcte} ! C'était un \textbf{piège} : pas d'erreur. \cle{都} est bien présent.
    
    \item \textbf{Erreur :} Comparatif incomplet (nuance neutre acceptée à l'oral, mais insuffisant à l'écrit/Bac). Manque de degré. \\
    \textbf{Correction :} \zh{网络比传媒\textbf{快得多}。} (Wǎngluò bǐ chuánméi \textbf{kuài dé duō}.) \\
    \textbf{Justification :} À l'écrit formel, un comparatif \cle{比} sans degré (\cle{得多}, \cle{一点儿}, chiffre) sonne « inachevé ». On précise l'écart.
\end{enumerate}
\end{exoresolu}

\courssub{Exercices d'application et de transformation}

\begin{exercice}[Entraînement écrit — Transposition et correction]
\textbf{A. Traduisez en chinois (caractères + pinyin) en utilisant la structure imposée.}
\begin{enumerate}
    \item [\cle{不管是……或者是……都……}] Que ce soit le journal papier ou la télévision, les jeunes préfèrent Internet.
    \item [\cle{是……的}] C'est à Douala que j'ai acheté ce téléphone l'année dernière. (Insistance sur le lieu).
    \item [\cle{比字句 + degré}] Le réseau 5G est beaucoup plus rapide que la 4G.
    \item [\cle{是……的}] C'est pour regarder la CAN (Coupe d'Afrique des Nations) qu'il a acheté un grand écran.
    \item [\cle{不管是……或者是……也……}] Qu'il fasse chaud ou qu'il fasse froid, il court tous les matins.
\end{enumerate}

\textbf{B. Corrigez l'erreur unique dans chaque phrase (caractères + pinyin).}
\begin{enumerate}
    \item \zh{不管他是老师或者是学生，都要学习。}
    \item \zh{我是坐飞机来了北京的。}
    \item \zh{中国比喀麦隆很大。}
    \item \zh{不管下雨或者不下雨，我去。}
    \item \zh{这本书比那本书贵十块钱了。}
\end{enumerate}
\end{exercice}

\begin{corrige}[Exercice Entraînement écrit]
\textbf{A. Traduction proposée :}
\begin{enumerate}
    \item \zh{不管是报纸，或者是电视，年轻人都喜欢上网。} (Bùguǎn shì bàozhǐ, huòzhě shì diànshì, niánqīngrén dōu xǐhuan shàngwǎng.)
    \item \zh{我是去年在杜阿拉买这手机的。} (Wǒ shì qùnián zài Dù'ālā mǎi zhè shǒujī de.) \textit{Note : \cle{在杜阿拉} placé après \cle{是}.}
    \item \zh{5G网络比4G网络快得多。} (5G wǎngluò bǐ 4G wǎngluò kuài dé duō.)
    \item \zh{他是为看非洲杯买大电视的。} (Tā shì wèi kàn Fēizhōu Bēi mǎi dà diànshì de.) \textit{Note : \cle{为/为了} introduit le but.}
    \item \zh{不管是天热，或者是天冷，他也每天早上跑步。} (Bùguǎn shì tiān rè, huòzhě shì tiān lěng, tā yě měitiān zǎoshang pǎobù.)
\end{enumerate}

\textbf{B. Correction des erreurs :}
\begin{enumerate}
    \item \textbf{Erreur :} \cle{是} placé devant le sujet \cle{他} au lieu de devant les noms \cle{老师/学生}. \cle{不管是} introduit des noms/pronoms, pas une proposition entière avec verbe. \\
    \textbf{Correction :} \zh{不管是\textbf{老师}或者是\textbf{学生}，都要学习。} (Bùguǎn shì \textbf{lǎoshī} huòzhě shì \textbf{xuéshēng}, dōu yào xuéxí.)
    
    \item \textbf{Erreur :} \cle{了} intercalé après \cle{来} (verbe de mouvement) + \cle{的} final. Double marquage. \\
    \textbf{Correction :} \zh{我是坐飞机来北京\textbf{的}。} (Wǒ shì zuò fēijī lái Běijīng \textbf{de}.)
    
    \item \textbf{Erreur :} \cle{很} interdit dans \cle{比} phrase. \\
    \textbf{Correction :} \zh{中国比喀麦隆\textbf{大得多}。} (Zhōngguó bǐ Kāmàilóng \textbf{dà dé duō}.) \textit{Ou : \zh{大很多}。}
    
    \item \textbf{Erreur :} \cle{不管} + Verbe (\cle{下雨}) $\rightarrow$ structure conditionnelle \cle{不管……都……}, mais \cle{或者是} n'est pas utilisé avec des verbes (on utilise \cle{还是} ou \cle{不} + V). De plus, \cle{都} manquant. \\
    \textbf{Correction :} \zh{不管\textbf{下不下雨}，我\textbf{都}去。} (Bùguǎn \textbf{xià bu xiàyǔ}, wǒ \textbf{dōu} qù.) \textit{Structure standard : \cle{不管 V 不 V, 都...}}
    
    \item \textbf{Erreur :} \cle{了} final interdit dans \cle{比} phrase (pas d'aspect perfectif sur la comparaison elle-même). \\
    \textbf{Correction :} \zh{这本书比那本书贵十块钱。} (Zhè běn shū bǐ nà běn shū guì shí kuài qián.)
\end{enumerate}
\end{corrige}

\coursec{Exercices d'application et de transformation}

Dans la section précédente, nous avons comparé le chinois et le français et repéré les pièges qui guettent les apprenants francophones : l'oubli de \zh{都} après \zh{不管是}, la confusion entre \zh{是……的} et le passé composé, le renversement de l'ordre dans la comparaison en \zh{比}. Il est maintenant temps de passer de la théorie à la pratique. Les six exercices qui suivent sont conçus selon une progression pédagogique rigoureuse : de la simple reconnaissance des structures jusqu'à la production libre, en passant par la complétion et la transformation. C'est exactement le parcours exigé au baccalauréat.

\begin{popfigure}
\begin{tikzpicture}[node distance=0.35cm]
\node[draw=popblue, fill=popblueL, rounded corners, minimum width=2.6cm, minimum height=1cm, align=center] (e1) {\textbf{Ex. 1}\\Reconnaissance};
\node[draw=popgreen, fill=popgreenL, rounded corners, minimum width=2.6cm, minimum height=1cm, align=center, right=of e1] (e2) {\textbf{Ex. 2}\\Complétion};
\node[draw=poporange, fill=poporangeL, rounded corners, minimum width=2.6cm, minimum height=1cm, align=center, right=of e2] (e3) {\textbf{Ex. 3--4}\\Transformation};
\node[draw=poppurple, fill=poppurpleL, rounded corners, minimum width=2.6cm, minimum height=1cm, align=center, right=of e3] (e4) {\textbf{Ex. 5--6}\\Production};
\draw[-{Stealth}, thick, popdark] (e1) -- (e2);
\draw[-{Stealth}, thick, popdark] (e2) -- (e3);
\draw[-{Stealth}, thick, popdark] (e3) -- (e4);
\draw[-{Stealth}, very thick, poppink, dashed] ([yshift=-0.5cm]e1.south west) -- ([yshift=-0.5cm]e4.south east) node[midway, below, poppink, align=center]{\textbf{autonomie croissante}};
\end{tikzpicture}
\legende{Progression des exercices : de la reconnaissance guidée à la production autonome.}
\end{popfigure}

\courssub{Rappel synthétique des trois structures}

\begin{propriete}[Structure 1 : Concession indifférente \zh{不管是……或者是……都}]
\emph{Schéma} : \zh{不管是 + Option A + 或者是 + Option B, Sujet + 都 + Verbe/Adjectif}.\\
\emph{Emploi} : Exprime que quelle que soit l'alternative choisie, la conséquence reste la même (« que ce soit… ou…, toujours… »).\\
\emph{Points clés} : \zh{不管} en tête ; \zh{是} devant chaque option (nom, pronom, préposition) ; \zh{或者是} (ou \zh{还是}) relie les options ; \zh{都} (ou \zh{也}) \textbf{obligatoire} avant le prédicat de la principale.
\end{propriete}

\begin{propriete}[Structure 2 : Mise en relief \zh{是……的} (action passée)]
\emph{Schéma} : \zh{Sujet + 是 + Circonstance (temps/lieu/moyen/but) + Verbe + 的 (+ Objet)}.\\
\emph{Emploi} : Met l'accent sur \textbf{une circonstance} d'une action \textbf{déjà accomplie} (passé). Répond à une question « quand ? où ? comment ? pourquoi ? ».\\
\emph{Règle d'or} : \textbf{Pas de \zh{了}} dans la structure \zh{是……的}. L'accompli est porté par la structure elle-même. Position de l'objet : \zh{是……V 的 O} (focus sur la circonstance) ou \zh{Objet 是……V 的} (focus sur l'objet + circonstance).
\end{propriete}

\begin{propriete}[Structure 3 : Comparaison \zh{比} (比字句)]
\emph{Schéma} : \zh{A + 比 + B + Adjectif (+ Degré : 得多 / 一点儿 / Chiffre + Classificateur)}.\\
\emph{Emploi} : Compare A et B ; A possède la qualité à un degré supérieur.\\
\emph{Ordre} : L'élément supérieur (\textbf{A}) vient \textbf{avant} \zh{比}. \zh{没有} pour l'infériorité/égalité. \zh{得多} = écart grand ; \zh{一点儿} = écart petit.
\end{propriete}

\courssub{Exercice 1 : reconnaissance des structures}

\begin{methode}[Comment reconnaître une structure ?]
\begin{enumerate}
\item Lisez la phrase entière et repérez les \cle{mots-outils} : \zh{不管}, \zh{或者}, \zh{都}, \zh{是}, \zh{的}, \zh{比}.
\item Identifiez la structure complète : un mot-outil isolé ne suffit pas (\zh{是} seul n'est pas forcément une structure \zh{是……的} ; il faut vérifier la présence du \zh{的} final et l'absence de \zh{了}).
\item Nommez la \cle{fonction} : concession indifférente (« que ce soit… ou… »), mise en relief d'une circonstance d'une action passée (« c'est… que »), ou comparaison (« plus… que »).
\end{enumerate}
\end{methode}

\begin{exercice}[Exercice 1 — Souligner et nommer]
Dans chaque phrase, soulignez la structure étudiée et nommez sa fonction.\\
1. \zh{不管是晴天，或者是雨天，他都听广播。} (\emph{Bùguǎn shì qíngtiān, huòzhě shì yǔtiān, tā dōu tīng guǎngbō.})\\
2. \zh{这条新闻是我昨天晚上在网上看到的。} (\emph{Zhè tiáo xīnwén shì wǒ zuótiān wǎnshang zài wǎngshang kàndào de.})\\
3. \zh{手机比报纸快得多。} (\emph{Shǒujī bǐ bàozhǐ kuài de duō.})\\
4. \zh{不管是老师，或者是学生，都喜欢这个节目。} (\emph{Bùguǎn shì lǎoshī, huòzhě shì xuésheng, dōu xǐhuan zhège jiémù.})\\
5. \zh{电视的观众比广播的听众多。} (\emph{Diànshì de guānzhòng bǐ guǎngbō de tīngzhòng duō.})\\
6. \zh{他们是在杜阿拉买的这部手机。} (\emph{Tāmen shì zài Dùālā mǎi de zhè bù shǒujī.})
\end{exercice}

\begin{corrige}[Exercice 1]
\begin{enumerate}
\item \zh{不管是晴天，或者是雨天，他\textbf{都}听广播。} → structure \zh{不管是……或者是……都} ; fonction : concession indifférente (« qu'il fasse beau ou qu'il pleuve, il écoute la radio »).
\item \zh{这条新闻\textbf{是}我昨天晚上在网上看到\textbf{的}。} → structure \zh{是……的} ; fonction : mise en relief du moment et du lieu de l'action passée (« c'est hier soir sur Internet que j'ai vu cette nouvelle »).
\item \zh{手机\textbf{比}报纸快\textbf{得多}。} → structure \zh{比} ; fonction : comparaison avec écart marqué (« le téléphone est beaucoup plus rapide que le journal »).
\item \zh{\textbf{不管是}老师，\textbf{或者是}学生，\textbf{都}喜欢这个节目。} → concession indifférente (« professeurs comme élèves, tous aiment cette émission »).
\item \zh{电视的观众\textbf{比}广播的听众多。} → comparaison simple (« la télévision a plus de spectateurs que la radio n'a d'auditeurs »).
\item \zh{他们\textbf{是}在杜阿拉买\textbf{的}这部手机。} → \zh{是……的} ; mise en relief du lieu (« c'est à Douala qu'ils ont acheté ce téléphone »). \emph{Note : l'objet \zh{这部手机} placé après \zh{的} est correct (structure narrative).}
\end{enumerate}
\end{corrige}

\courssub{Exercice 2 : complétion}

\begin{exoresolu}[Exercice 2 — Insérer le bon mot-outil]
Énoncé : complétez les phrases suivantes avec \zh{不管}、\zh{是}、\zh{的}、\zh{比}、\zh{都}、\zh{得多} (chaque mot peut être utilisé plusieurs fois).\\
1. \zh{＿＿＿在城市，或者在乡下，年轻人＿＿＿用手机。}\\
2. \zh{这条消息＿＿＿今天早上广播＿＿＿。}\\
3. \zh{网络＿＿＿报纸快＿＿＿。}
\tcblower
Solution détaillée :\\
1. \zh{\textbf{不管}\textbf{是}在城市，\textbf{或者}\textbf{是}在乡下，年轻人\textbf{都}用手机。} (\emph{Bùguǎn shì zài chéngshì, huòzhě shì zài xiāngxià, niánqīngrén dōu yòng shǒujī.}) — « Qu'en ville ou à la campagne, les jeunes utilisent tous le téléphone. » Justification : la structure \zh{不管是……或者是……} exige \zh{是} devant chaque alternative (même prépositionnelle) et \zh{都} avant le verbe.\\
2. \zh{这条消息\textbf{是}今天早上广播\textbf{的}。} (\emph{Zhè tiáo xiāoxi shì jīntiān zǎoshang guǎngbō de.}) — « C'est ce matin que cette information a été diffusée. » Justification : mise en relief du moment d'une action passée → \zh{是……的}, sans \zh{了}. \zh{广播} fonctionne comme verbe « diffuser » à la voix passive.\\
3. \zh{网络\textbf{比}报纸快\textbf{得多}。} (\emph{Wǎngluò bǐ bàozhǐ kuài de duō.}) — « Internet est beaucoup plus rapide que le journal. » Justification : A + \zh{比} + B + adjectif ; \zh{得多} marque un grand écart.
\end{exoresolu}

\courssub{Exercice 3 : transformation vers \zh{是……的}}

\begin{methode}[Transformer une phrase simple en \zh{是……的}]
\begin{enumerate}
\item Repérez la \cle{circonstance} à mettre en relief (moment, lieu, manière, moyen).
\item Placez \zh{是} juste \textbf{avant} cette circonstance.
\item Placez \zh{的} juste \textbf{après} le verbe (si l'objet suit) ou en \textbf{fin de phrase} (si l'objet est antéposé).
\item \textbf{Supprimez le \zh{了}} de la phrase d'origine : \zh{了} et \zh{是……的} ne cohabitent pas pour la même action.
\end{enumerate}
\end{methode}

\begin{attention}[Le piège du \zh{了}]
Dans la transformation vers \zh{是……的}, il faut \cle{supprimer le \zh{了}} de la phrase d'origine. On ne dit jamais \zh{我是昨天在网上看到了这条新闻的}. La structure \zh{是……的} suffit à elle seule à indiquer que l'action est accomplie ; ajouter \zh{了} est une erreur classique des francophones, qui veulent « traduire » le passé composé deux fois.
\end{attention}

\begin{exercice}[Exercice 3 — Transformation]
Phrase d'origine : \zh{我昨天在网上看到这条新闻。} (\emph{Wǒ zuótiān zài wǎngshang kàndào zhè tiáo xīnwén.}) — « J'ai vu cette nouvelle hier sur Internet. »\\
Transformez-la en \zh{是……的} pour insister sur le \textbf{lieu} (sur Internet).
\end{exercice}

\begin{corrige}[Exercice 3]
\zh{我\textbf{是}昨天\textbf{在网上}看到\textbf{的}这条新闻。} (\emph{Wǒ shì zuótiān zài wǎngshang kàndào de zhè tiáo xīnwén.}) — « C'est \textbf{sur Internet} que j'ai vu cette nouvelle hier. »\\
Vérification : \zh{是} précède le complément de lieu \zh{在网上} ; \zh{的} suit le verbe \zh{看到} ; aucun \zh{了}. Variante acceptée (objet antéposé) : \zh{这条新闻是我昨天在网上看到的。}
\end{corrige}

\courssub{Exercice 4 : construire des comparaisons en \zh{比}}

Rappel du schéma : \cle{A + 比 + B + adjectif (+ 得多 / 一点儿 / 数量)}. L'élément jugé supérieur vient \textbf{avant} \zh{比}.

\begin{center}
\begin{tabularx}{\linewidth}{|X|X|X|}
\hline
\rowcolor{popblueL} Données & Structure & Exemple attendu \\
\hline
Téléphone 80 000 F / téléphone 30 000 F & A + 比 + B + 贵 & \zh{这部手机比那部手机贵。} \\
\hline
Internet / courrier (vitesse) & A + 比 + B + 快得多 & \zh{网络比信件快得多。} \\
\hline
Radio / télévision (audience au village) & A + 比 + B + 多 & \zh{在农村，广播的听众比电视的观众多。} \\
\hline
Télévision / presse écrite (jeunes) & A + 比 + B + 受欢迎 & \zh{对年轻人来说，电视比报纸受欢迎。} \\
\hline
\end{tabularx}
\end{center}

\begin{exercice}[Exercice 4 — Production de comparaisons]
À partir des données du tableau, construisez quatre phrases complètes en \zh{比}. Ajoutez \zh{得多} ou \zh{一点儿} quand l'écart le justifie.
\end{exercice}

\begin{corrige}[Exercice 4]
1. \zh{这部手机比那部手机贵得多。} (\emph{Zhè bù shǒujī bǐ nà bù shǒujī guì de duō.}) — « Ce téléphone est beaucoup plus cher que celui-là. »\\
2. \zh{网络比信件快得多。} (\emph{Wǎngluò bǐ xìnjiàn kuài de duō.}) — « Internet est beaucoup plus rapide que le courrier. »\\
3. \zh{在农村，广播的听众比电视的观众多。} (\emph{Zài nóngcūn, guǎngbō de tīngzhòng bǐ diànshì de guānzhòng duō.}) — « À la campagne, la radio a plus d'auditeurs que la télévision n'a de téléspectateurs. »\\
4. \zh{对年轻人来说，电视比报纸受欢迎一点儿。} (\emph{Duì niánqīngrén láishuō, diànshì bǐ bàozhǐ shòu huānyíng yìdiǎnr.}) — « Pour les jeunes, la télévision est un peu plus appréciée que la presse. »
\end{corrige}

\courssub{Exercice 5 : thème (français → chinois)}

\begin{methode}[Réussir le thème en trois étapes]
\begin{enumerate}
\item \textbf{Repérer} la structure française : « que ce soit… ou… » → \zh{不管是……或者是……都} ; « c'est… que » (action passée) → \zh{是……的} ; « plus… que » → \zh{比}.
\item \textbf{Construire} la phrase chinoise selon le schéma, dans l'ordre chinois (circonstances avant le verbe).
\item \textbf{Vérifier} les mots-outils obligatoires : \zh{都} après \zh{不管是} ; \zh{的} final et absence de \zh{了} dans \zh{是……的} ; A avant \zh{比} dans la comparaison.
\end{enumerate}
\end{methode}

\begin{exemplebox}[Modèle de thème réussi]
Français : « Qu'en ville ou à la campagne, les jeunes ne peuvent se passer du portable. »\\
Chinois : \zh{不管是城市，或者是农村，年轻人都离不开手机。}\\
Pinyin : \emph{Bùguǎn shì chéngshì, huòzhě shì nóngcūn, niánqīngrén dōu líbukāi shǒujī.}\\
Vérification : \zh{不管} en tête, \zh{或者} devant la seconde alternative, \zh{都} avant le verbe \zh{离不开}.
\end{exemplebox}

\begin{exercice}[Exercice 5 — Traduisez en chinois]
1. « Que ce soit en ville ou au village, les jeunes camerounais utilisent tous le téléphone portable. »\\
2. « C'est par la radio que ma grand-mère a appris la nouvelle. »\\
3. « Internet est beaucoup plus rapide que le courrier. »
\end{exercice}

\begin{corrige}[Exercice 5]
1. \zh{不管是在城市，或者是在农村，喀麦隆的年轻人都用手机。} (\emph{Bùguǎn shì zài chéngshì, huòzhě shì zài nóngcūn, Kāmàilóng de niánqīngrén dōu yòng shǒujī.}) — \zh{是} ajouté devant les prépositions \zh{在} pour respecter la structure \zh{不管是……或者是……} ; \zh{都} obligatoire avant \zh{用}.\\
2. \zh{我奶奶是通过广播知道这个消息的。} (\emph{Wǒ nǎinai shì tōngguò guǎngbō zhīdào zhège xiāoxi de.}) — \zh{是} devant \zh{通过广播} (le moyen), \zh{的} final, pas de \zh{了}. \zh{通过} (par le biais de) traduit « par ».\\
3. \zh{网络比信件快得多。} (\emph{Wǎngluò bǐ xìnjiàn kuài de duō.}) — Internet (A) avant \zh{比}, courrier (B) après ; \zh{得多} traduit « beaucoup plus ».
\end{corrige}

\courssub{Exercice 6 : version (chinois → français)}

\begin{textebox}[Mini-texte de version : les médias au village]
\zh{这条消息是今天早上广播的，不管是老人，或者是孩子，都知道了。网络比报纸快得多。}\\
\emph{Zhè tiáo xiāoxi shì jīntiān zǎoshang guǎngbō de, bùguǎn shì lǎorén, huòzhě shì háizi, dōu zhīdào le. Wǎngluò bǐ bàozhǐ kuài de duō.}\\
Traduction : « C'est ce matin que cette information a été diffusée à la radio ; que ce soit les personnes âgées ou les enfants, tous la connaissent déjà. Internet est beaucoup plus rapide que la presse écrite. »
\end{textebox}

\begin{center}
\begin{tabularx}{\linewidth}{|X|X|X|X|}
\hline
\rowcolor{popblueL} Caractères & Pinyin & Nature & Traduction \\
\hline
消息 & xiāoxi & nom & information, nouvelle \\
\hline
广播 & guǎngbō & verbe / nom & diffuser / la radio \\
\hline
老人 & lǎorén & nom & personne âgée \\
\hline
知道 & zhīdào & verbe & savoir, connaître \\
\hline
报纸 & bàozhǐ & nom & journal, presse écrite \\
\hline
\end{tabularx}
\end{center}

\begin{exercice}[Exercice 6 — Version]
Traduisez le mini-texte ci-dessus en français, puis soulignez dans le texte chinois les trois structures étudiées et justifiez chaque choix de traduction.
\end{exercice}

\begin{corrige}[Exercice 6]
Traduction : voir la boîte du texte. Justifications :\\
1. \zh{是今天早上广播的} → structure \zh{是……的} : le français rend la mise en relief par « c'est ce matin que… » ; le passif français « a été diffusée » traduit naturellement le verbe \zh{广播} sans sujet agent.\\
2. \zh{不管是老人，或者是孩子，都} → concession indifférente : « que ce soit… ou…, tous… » ; le \zh{都} se traduit par « tous ».\\
3. \zh{网络比报纸快得多} → comparaison : « beaucoup plus rapide que » ; \zh{得多} impose l'intensif « beaucoup ».
\end{corrige}

\courssub{Correction collective et bilan}

\begin{experience}[Correction collective en classe]
L'enseignant écrit les six exercices au tableau. Pour chaque réponse proposée par un élève, la classe applique la démarche en trois questions : \textbf{Quelle structure ? Quel schéma ? Quels mots-outils obligatoires ?} Toute réponse n'est validée que si l'élève cite la règle (par exemple : « \zh{都} est obligatoire après \zh{不管是} » ou « on supprime \zh{了} dans \zh{是……的} »). Les erreurs sont notées au tableau dans une colonne « pièges des francophones » et corrigées collectivement.
\end{experience}

\begin{aretenir}[Bilan de la leçon]
\begin{itemize}
\item \zh{不管是……或者是……都} : « que ce soit… ou…, tous… » — ne jamais oublier \zh{都} ; répéter \zh{是} devant chaque alternative.
\item \zh{是……的} : mise en relief d'une circonstance d'une action passée — \zh{是} avant la circonstance, \zh{的} après le verbe (ou fin de phrase), \textbf{sans} \zh{了}.
\item \zh{A 比 B + adjectif} : comparaison — l'élément supérieur avant \zh{比} ; \zh{得多} (grand écart), \zh{一点儿} (petit écart).
\item En thème comme en version, toujours repérer d'abord la structure, puis construire, puis vérifier les mots-outils.
\end{itemize}
\end{aretenir}

\newpage

\coursec{Activité d'intégration}

\coursec{Leçon 35 — Grammaire : 不管是…….或者是……. ; 是。。。。。。的 ; 比字句}

\courssub{Objectifs de la leçon}
\begin{itemize}
    \item Maîtriser la construction et l'emploi des trois structures grammaticales majeures : \cle{不管是……或者是……} (universalité / concession), \cle{是……的} (insistance sur les circonstances d'une action passée) et \cle{比字句} (comparaison avec degrés d'intensité).
    \item Mobiliser ces structures dans une production écrite et orale cohérente sur le thème des médias (Module 5).
    \item Identifier et corriger les erreurs fréquentes des francophones (ordre des mots, place des compléments, mots-outils \zh{了}/\zh{很}, classificateurs).
    \item Enrichir le lexique des médias (HSK 3-4) et développer l'esprit critique sur l'accès à l'information au Cameroun.
\end{itemize}

\courssub{Situation déclenchante : Observation guidée}
\begin{textebox}[Contexte : Le journal du lycée bilingue de Yaoundé]
Le club de presse du \zh{雅温得双语高中} (Yǎwēndé Shuāngyǔ Gāozhōng) prépare un numéro spécial sur le thème \zh{« 青年与媒体 »} (Qīngnián yǔ Méitǐ — « Les jeunes et les médias »). En tant que membres de la rédaction, vous devez proposer un article court intitulé \zh{« 在喀麦隆，通过什么媒体获取资讯最好？ »} (Zài Kāmàilóng, tōngguò shénme Méitǐ huòqǔ Zīxùn zuì hǎo ? — « Au Cameroun, quel est le meilleur média pour s'informer ? »).

L'article doit refléter la réalité camerounaise : l'écart d'accès à Internet entre \zh{雅温得} (Yaoundé) et \zh{杜阿拉} (Douala) d'un côté, et les zones rurales (\zh{乡村} xiāngcūn) de l'autre ; le rôle historique de la \zh{广播} (guǎngbō - radio) et de la \zh{电视} (diànshì - télévision) publique (\zh{CRTV}) ; l'essor des \zh{社交媒体} (shèjiāo méitǐ - réseaux sociaux) comme WhatsApp et Facebook chez les lycéens. Vous devez convaincre vos camarades que votre analyse est la plus juste.
\end{textebox}

\courssub{Lexique essentiel : Les médias (Module 5 — HSK 3/4)}
\begin{center}
\begin{tabularx}{\linewidth}{|X|X|X|X|}
\hline
\rowcolor{popblueL} \textbf{Caractères} & \textbf{Pinyin} & \textbf{Nature} & \textbf{Traduction} \\
\hline
媒体 & méitǐ & nom & média, support d'information \\
\hline
电视 & diànshì & nom & télévision \\
\hline
广播 & guǎngbō & nom & radio, radiodiffusion \\
\hline
报纸 & bàozhǐ & nom & journal, presse écrite \\
\hline
互联网 & hùliánwǎng & nom & Internet \\
\hline
手机 & shǒujī & nom & téléphone portable, smartphone \\
\hline
社交媒体 & shèjiāo méitǐ & nom & réseaux sociaux \\
\hline
微信 & Wēixìn & nom propre & WeChat \\
\hline
资讯 & zīxùn & nom & information, actualité (terme formel) \\
\hline
新闻 & xīnwén & nom & nouvelles, journal télévisé \\
\hline
消息 & xiāoxi & nom & nouvelle, info, message \\
\hline
获取 & huòqǔ & verbe & obtenir, se procurer (formel) \\
\hline
浏览 & liúlǎn & verbe & parcourir, consulter (écran, journal) \\
\hline
核实 & héshí & verbe & vérifier, confirmer (exactitude) \\
\hline
可靠 & kěkào & adj. & fiable, sûr \\
\hline
方便 & fāngbiàn & adj. & pratique, commode \\
\hline
偏远 & piānyuǎn & adj. & reculé, isolé (village) \\
\hline
识字 & shízì & verbe & savoir lire/écrire (litt. connaître les caractères) \\
\hline
权威 & quánwēi & adj. & autoritaire, officiel, faisant autorité \\
\hline
组合使用 & zǔhé shǐyòng & v. + adv. & utiliser de façon combinée \\
\hline
\end{tabularx}
\end{center}

\courssub{Structure 1 : 不管是……或者是…… (Qu'il s'agisse de… ou de…)}
\begin{propriete}[Règle : 不管是……或者是…… / 还是……]
\textbf{Schéma :} \textbf{Sujet + 不管是 + Option A + 或者是 / 还是 + Option B + ，都 / 也 + Verbe / Adjectif} \\
\textbf{Emploi :} Énumère deux alternatives contrastées (extrêmes, opposées ou simplement différentes) pour affirmer que le prédicat s'applique aux deux sans distinction. \zh{不管是} introduit la concession. \zh{或者是} (écrit/formel) ou \zh{还是} (oral/courant) lie les options. \zh{都} (tous) ou \zh{也} (aussi) est \textbf{obligatoire} dans la principale pour marquer l'universalité.
\begin{center}
\begin{tabularx}{\linewidth}{|X|X|X|}
\hline
\rowcolor{popblueL} Structure & Exemple (Caractères + Pinyin) & Traduction \\ \hline
Suj + 不管是 A 或者是 B, 都... & \zh{不管是在雅温得，或者是在偏远的乡村，获取资讯都是必需品。} (Bùguǎn shì zài Yǎwēndé, huòzhě shì zài piānyuǎn de xiāngcūn, huòqǔ zīxùn dōu shì bìxūpǐn.) & Que ce soit à Yaoundé ou dans des villages reculés, s'informer est un besoin vital. \\ \hline
Suj + 不管是 A 还是 B, 也... & \zh{不管是看电视，还是刷手机，都要辨别真假。} (Bùguǎn shì kàn diànshì, háishì shuā shǒujī, dōu yào biébié zhēnjiǎ.) & Qu'on regarde la télé ou qu'on scrolle sur son téléphone, il faut distinguer le vrai du faux. \\ \hline
Suj + 不管是 A 或者是 B, 都不... (Négation) & \zh{不管是报纸，或者是广播，年轻人都不太用了。} (Bùguǎn shì bàozhǐ, huòzhě shì guǎngbō, niánqīngrén dōu bù tài yòng le.) & Ni le journal ni la radio, les jeunes ne les utilisent plus beaucoup. \\ \hline
\end{tabularx}
\end{center}
\end{propriete}

\begin{attention}[Différence avec \zh{无论……还是……}]
\zh{无论……还是……} est plus formel, plus littéraire, souvent à l'écrit (éditoriaux, discours). \zh{不管是……或者是/还是……} est standard à l'oral et à l'écrit courant (presse, essais scolaires). Au Bac, les deux sont acceptés, mais \zh{不管是} est plus naturel pour un article de presse lycéenne.
\end{attention}

\courssub{Structure 2 : 是……的 (Insistance sur les circonstances d'une action passée)}
\begin{propriete}[Règle : 是……的 — Focus sur le passé accompli]
\textbf{Schéma :} \textbf{Sujet + 是 + [Complément Temps / Lieu / Manière / Auteur / But / Moyen] + Verbe + 的} \\
\textbf{Règle d'or :} La structure \zh{是……的} ne s'utilise que pour des \textbf{actions déjà accomplies} (passé). Elle ne décrit \textbf{PAS} l'action elle-même (déjà connue), mais met l'accent sur \textbf{QUAND, OÙ, COMMENT, PAR QUI, POURQUOI, AVEC QUOI} elle s'est produite. \textbf{Interdiction :} \zh{了} ne s'emploie généralement \textbf{PAS} après le verbe dans cette structure (sauf complément de résultat/direction : \zh{看到的}, \zh{拿来的}).
\begin{center}
\begin{tabularx}{\linewidth}{|X|X|X|}
\hline
\rowcolor{poppurpleL} Complément mis en relief & Exemple (Caractères + Pinyin) & Traduction \\ \hline
\textbf{Temps} (\zh{时候}, \zh{昨天}, \zh{上周}) & \zh{这条消息是我们昨天在手机上看到的。} (Zhè tiáo xiāoxi shì wǒmen zuótiān zài shǒujī shàng kàn dào de.) & C'est \textbf{hier} sur le portable qu'on a vu cette info. \\ \hline
\textbf{Lieu} (\zh{地方}, \zh{学校}, \zh{网上}) & \zh{这个视频是在学校图书馆下载的。} (Zhège shìpín shì zài xuéxiào túshūguǎn xiàzǎi de.) & C'est \textbf{à la biblio de l'école} que cette vidéo a été téléchargée. \\ \hline
\textbf{Moyen / Instrument} (\zh{通过}, \zh{用}, \zh{坐}) & \zh{这个消息是老师通过微信群发给我们的。} (Zhège xiāoxi shì lǎoshī tōngguò Wēixìn qún fā gěi wǒmen de.) & C'est \textbf{via le groupe WeChat} que le prof nous l'a envoyée. \\ \hline
\textbf{Auteur / Agent} & \zh{这篇报道是记者写的。} (Zhè piān bàodǎo shì jìzhě xiě de.) & C'est \textbf{le journaliste} qui a rédigé ce reportage. \\ \hline
\end{tabularx}
\end{center}
\end{propriete}

\begin{exemplebox}[Transformation SVO + 了 $\rightarrow$ 是……的]
\begin{itemize}
    \item Fait brut (nouvelle info) : \zh{我们昨天在微信群里看到了这条消息。} (Nous avons vu cette info hier sur WeChat.) $\rightarrow$ Focus sur les circonstances : \zh{这条消息是我们昨天在微信群里看到的。} (C'est hier sur WeChat qu'on a vu cette info.)
    \item \textbf{Erreur type :} *\zh{我们是昨天看了新闻的} (Redondance \zh{了}+\zh{的}, ordre incorrect). $\rightarrow$ \zh{我们是昨天看新闻的} ou \zh{新闻是我们昨天看的}.
\end{itemize}
\end{exemplebox}

\courssub{Structure 3 : La phrase comparative en 比 (比字句) — Degrés d'intensité}
\begin{propriete}[Règle : 比字句 — Structure de base et degrés]
\textbf{Schéma de base :} \textbf{A + 比 + B + Adjectif (+ Degré)} \\
\textbf{Règles absolues :}
\begin{enumerate}
    \item \zh{很}, \zh{非常}, \zh{特别}, \zh{真} sont \textbf{INTERDITS} devant l'adjectif dans une phrase \zh{比}.
    \item Le degré (complément de degré) se place \textbf{APRÈS} l'adjectif.
    \item Pour exprimer « A n'est pas aussi Adj que B » : \textbf{A + 没有 + B + 这么/那么 + Adj}.
\end{enumerate}
\begin{center}
\begin{tabularx}{\linewidth}{|X|X|X|}
\hline
\rowcolor{popgreenL} Degré d'intensité & Structure / Exemple (Caractères + Pinyin) & Traduction / Nuance \\ \hline
\textbf{Beaucoup plus} (marqué) & \zh{得多} / \zh{多了} & \zh{手机比报纸快得多。} (Shǒujī bǐ bàozhǐ kuài dé duō.) Le portable est \textbf{beaucoup plus} rapide que le journal. \\ \hline
\textbf{Un peu plus} (nuancé) & \zh{一点儿} / \zh{一些} / \zh{有点儿} & \zh{买报纸比上网贵一点儿。} (Mǎi bàozhǐ bǐ shàngwǎng guì yìdiǎnr.) Acheter le journal est \textbf{un peu plus} cher que le net. \\ \hline
\textbf{Équivalence (négatif)} & \zh{A 没有 B 这么/那么 Adj} & \zh{广播没有电视那么直观。} (Guǎngbō méiyǒu diànshì nàme zhíguān.) La radio n'est pas aussi visuelle que la télé. \\ \hline
\textbf{Superlatif (dans un groupe)} & \zh{最} / \zh{最...的} & \zh{在这四种媒体里，手机最方便。} (Zài zhè sì zhǒng méitǐ lǐ, shǒujī zuì fāngbiàn.) Parmi ces 4 médias, le portable est \textbf{le plus} pratique. \\ \hline
\end{tabularx}
\end{center}
\end{propriete}

\begin{exemplebox}[Comparatif avec verbes d'action (比 + V + 得 + Adj/Degré)]
Lorsque l'on compare la \textbf{manière} dont une action est faite, on utilise \zh{比} + Verbe + \zh{得} + Adjectif/Degré.
\begin{itemize}
    \item \zh{他跑得比我快得多。} (Tā pǎo de bǐ wǒ kuài dé duō.) Il court beaucoup plus vite que moi.
    \item \zh{这个软件用起来比那个方便多了。} (Zhège ruǎnjiàn yòng qǐlái bǐ nàge fāngbiàn duō le.) Ce logiciel est bien plus pratique à utiliser que celui-là.
\end{itemize}
\end{exemplebox}

\courssub{Comparaison avec le français et pièges des francophones (Terminale A)}
\begin{attention}[Erreurs fréquentes — Check-list de relecture]
\begin{itemize}
    \item \textbf{Ordre \zh{是……的} :} *\zh{我们是昨天看了新闻的} (Faux : \zh{了} redondant + place de \zh{昨天}). $\rightarrow$ \textbf{Correct :} \zh{我们是昨天看新闻的} (Sujet = nous) OU \zh{新闻是我们昨天看的} (Sujet = nouvelle, passif naturel).
    \item \textbf{Degré dans \zh{比} :} *\zh{手机比电视很方便} / *\zh{手机比电视更方便得多} (Faux : \zh{很} interdit ; \zh{更} et \zh{得多} redondants). $\rightarrow$ \textbf{Correct :} \zh{手机比电视方便得多} / \zh{方便多了} / \zh{方便一点儿}.
    \item \textbf{Sujet commun dans \zh{不管是} :} *\zh{不管是下雨或者是晴天，我去} (Maladroit : sujets météorologiques vs agent humain). $\rightarrow$ \textbf{Correct :} \zh{不管是下雨还是晴天，我都去} (Sujet unique \zh{我} pour les deux options).
    \item \textbf{Classificateurs médias (obligatoires) :} \zh{一条} 消息/新闻 (info), \zh{一期} 节目 (émission), \zh{一份} 报纸 (journal), \zh{一个} 视频 / 网站 / 群 / 手机.
    \item \textbf{Vocabulaire piège :} \zh{上网} (verbe : aller sur le net) vs \zh{网上} (locatif : sur le net / en ligne). Ex: \zh{我在网上看新闻} (Je lis les infos en ligne) vs \zh{我每天上网} (Je vais sur le net chaque jour). \zh{看新闻} (regarder/consulter les infos, générique) vs \zh{读报} (lire le journal papier, spécifique).
    \item \textbf{Ponctuation :} Utiliser la ponctuation chinoise pleine chasse (\zh{，} \zh{。} \zh{、} \zh{？}) dans les phrases chinoises. Pas d'espace entre caractères.
\end{itemize}
\end{attention}

\courssub{Activité d'intégration : Production écrite collaborative (Article de presse)}
\courssub{Consignes détaillées (Durée : 40 min)}
\begin{enumerate}
    \item \textbf{Analyse du sujet (5 min) :} En groupe de 4, identifiez les 4 médias à comparer : \zh{电视}, \zh{广播}, \zh{报纸}, \zh{互联网/手机}. Listez 2 avantages et 2 inconvénients pour chacun en chinois (vocabulaire Module 5 + Leçon 35).
    \item \textbf{Rédaction collaborative (25 min) :} Rédigez l'article (7--8 phrases) en respectant \textbf{scrupuleusement} les contraintes syntaxiques suivantes :
    \begin{itemize}
        \item \textbf{Phrase 1 (Accroche universelle) :} Utiliser \cle{不管是……或者是……} pour poser l'universalité du besoin d'info (villes/villages, jeunes/adultes, riches/pauvres). \zh{都} obligatoire.
        \item \textbf{Phrase 2 (Ancrage factuel passé) :} Utiliser \cle{是……的} pour situer précisément (temps, lieu, moyen) la réception d'une info récente (ex: match de foot, résultat bac, actualité locale). Pas de \zh{了}.
        \item \textbf{Phrases 3--6 (Comparaisons argumentées) :} Utiliser \textbf{au moins 3 phrases en \cle{比字句}} avec des degrés \textbf{variés} : \zh{得多}, \zh{多了}, \zh{一点儿}, \zh{一些}, \zh{没有...那么...}. Comparer : rapidité (\zh{快}), coût (\zh{便宜}/\zh{贵}), accessibilité (\zh{方便}), fiabilité (\zh{可靠}), profondeur (\zh{深入}).
        \item \textbf{Phrases 7--8 (Conclusion / Recommandation) :} Structure libre (\zh{因此}, \zh{建议}, \zh{首选}, \zh{组合使用}). Parallélisme rhétorique apprécié (\zh{手机查..., 电视看...}).
    \end{itemize}
    \item \textbf{Relecture croisée (5 min) :} Échangez avec un autre groupe. Vérifiez : ordre \zh{比} (A 比 B Adj), place du complément dans \zh{是……的} (avant \zh{的}), absence de \zh{很} dans comparatif, classificateurs, ponctuation chinoise.
    \item \textbf{Présentation orale et vote (5 min) :} Un porte-parole lit l'article (prononciation, tons, liaisons). La classe note les structures repérées sur grille. Vote : \zh{最有说服力} (le plus convaincant) / \zh{最地道} (le plus naturel).
\end{enumerate}

\courssub{Production modèle et corrigé commenté}
\begin{exoresolu}[Article modèle : « 在喀麦隆，通过什么媒体获取资讯最好？ »]
\textbf{Texte élève (Modèle cible HSK 4) :}
\zh{不管是在雅温得，或者是在偏远的乡村，获取资讯都是每天的必需品。} \\
\zh{这条关于足球比赛的消息，是我们昨天在微信群里看到的。} \\
\zh{手机上网比看电视快得多，也方便多了。} \\
\zh{但是，电视新闻比网上的消息可靠一些，因为有专业记者核实。} \\
\zh{广播比报纸便宜，而且不需要识字，适合老年人。} \\
\zh{报纸虽然贵一点儿，但是分析比较深入，适合做研究的学生。} \\
\zh{因此，我们建议同学们组合使用：手机查快讯，电视看权威新闻。} \\

\textbf{Pinyin (standard, tons sur voyelles, espaces inter-mots) :}
Bùguǎn shì zài Yǎwēndé, huòzhě shì zài piānyuǎn de xiāngcūn, huòqǔ zīxùn dōu shì měitiān de bìxūpǐn. \\
Zhè tiáo guānyú zúqiú bǐsài de xiāoxi, shì wǒmen zuótiān zài Wēixìn qún lǐ kàn dào de. \\
Shǒujī shàngwǎng bǐ kàn diànshì kuài dé duō, yě fāngbiàn duō le. \\
Dànshì, diànshì xīnwén bǐ wǎngshàng de xiāoxi kěkào yìxiē, yīnwèi yǒu zhuānyè jìzhě héshí. \\
Guǎngbō bǐ bàozhǐ piányi, érqiě bù xūyào shízì, shìhé lǎoniánrén. \\
Bàozhǐ suīrán guì yìdiǎnr, dànshì fēnxī bǐjiào shēnrù, shìhé zuò yánjiū de xuésheng. \\
Yīncǐ, wǒmen jiànyì tóngxuémen zǔhé shǐyòng: shǒujī chá kuàixùn, diànshì kàn quánwēi xīnwén. \\

\textbf{Traduction française littéraire :}
1. Que ce soit à Yaoundé ou dans des villages reculés, s'informer est un besoin quotidien. \\
2. Cette information concernant le match de football, c'est hier dans le groupe WeChat que nous l'avons vue. \\
3. Surfer sur smartphone est bien plus rapide et bien plus pratique que de regarder la télévision. \\
4. Cependant, les informations télévisées sont un peu plus fiables que celles du net, car des journalistes professionnels les vérifient. \\
5. La radio est moins chère que le journal, et ne nécessite pas de savoir lire, elle convient aux personnes âgées. \\
6. Le journal, bien que légèrement plus cher, offre des analyses relativement approfondies, il convient aux élèves faisant de la recherche. \\
7. C'est pourquoi nous conseillons aux camarades une utilisation combinée : smartphone pour les flashes, télévision pour les informations officielles. \\

\textbf{Commentaire linguistique détaillé (Corrigé type Bac) :}
\begin{enumerate}
    \item \textbf{Phrase 1 (\zh{不管是……或者是……}) :} Structure canonique. Opposition géographique \zh{雅温得} (urbain) vs \zh{偏远的乡村} (rural). \zh{获取资讯} (objet verbalisé) comme sujet commun. \zh{都是} répond à \zh{不管是}. \zh{必需品} (HSK 4) vocabulaire précis pour « besoin vital ».
    \item \textbf{Phrase 2 (\zh{是……的}) :} Topicalisation de l'objet \zh{这条消息} (classif. \zh{条} correct). Focus double : \textbf{Temps} \zh{昨天} + \textbf{Lieu virtuel} \zh{在微信群里}. Verbe \zh{看到} (V + complé. de résultat) compatible avec \zh{的} sans \zh{了}. Structure passive naturelle en chinois (\zh{消息是...看到的}).
    \item \textbf{Phrase 3 (\zh{比} + \zh{得多} / \zh{多了}) :} Sujet verbal \zh{手机上网} vs \zh{看电视}. Double adjectif \zh{快} (monosyll.) + \zh{方便} (bisyll.). Degrés \zh{得多} (écrit) et \zh{多了} (oral/écrit courant) en alternance stylistique. \zh{也} coordonne les deux qualités.
    \item \textbf{Phrase 4 (\zh{比} + \zh{一些}) :} Nuance \zh{一些} (un peu) sur \zh{可靠} (bisyll.). Justification causale \zh{因为...} (omission de \zh{所以} standard). Vocabulaire précis \zh{专业记者核实} (journalistes pros vérifient).
    \item \textbf{Phrase 5 (\zh{比} + Adj nu + \zh{而且}) :} Comparatif simple \zh{便宜} (moins cher). \zh{而且} ajoute argument fort (accessibilité \zh{不需要识字}). \zh{适合} + cible \zh{老年人}.
    \item \textbf{Phrase 6 (Concession \zh{虽然...但是...} + \zh{比} + \zh{一点儿}) :} Structure complexe HSK 4 maîtrisée. \zh{分析...深入} (collocation). \zh{做研究} (activité académique).
    \item \textbf{Phrase 7 (Conclusion \zh{因此} + \zh{建议} + Parallélisme) :} \zh{组合使用} (terme technique). Parallélisme \zh{手机查快讯 / 电视看权威新闻} (V+O // V+O) rhétoriquement très chinois, efficace pour une conclusion.
\end{enumerate}
\tcblower
\textbf{Variantes acceptées / Pistes de différenciation pédagogique :}
\begin{itemize}
    \item \textit{Niveau HSK 3 (simplifié / remédiation) :} Remplacer \zh{偏远的} par \zh{远的}, \zh{核实} par \zh{查} ou \zh{确认}, \zh{组合使用} par \zh{一起用}. Accepter \zh{多了} systématiquement au lieu de \zh{得多}. Tolérer \zh{比...更...} sans degré si l'élève bloque sur les degrés.
    \item \textit{Niveau HSK 4+ (approfondissement / excellence) :} Exiger \zh{不仅...而且...} dans la phrase 5 (\zh{广播不仅比报纸便宜，而且不需要识字}). Ajouter une phrase avec \zh{跟...比} (\zh{跟电视比，手机更灵活}) ou \zh{相比之下} (\zh{相比之下，手机的优势很明显}). Utiliser \zh{相比之下} en ouverture de phrase 7.
    \item \textit{Erreur corrigée type (exemple pour le tableau) :} Élève écrit : *\zh{手机比电视更快得多}. $\rightarrow$ Correction prof : \zh{更} et \zh{得多} redondants (double comparatif). Choix : \zh{手机比电视快得多} OU \zh{手机比电视更快}.
\end{itemize}
\end{exoresolu}

\courssub{Bilan grammatical et communicatif}
\begin{aretenir}[Synthèse des acquis — Leçon 35]
Cette activité d'intégration en situation authentique (presse lycéenne, contexte Cameroun) a permis de valider l'acquisition \textbf{opérationnelle} des trois structures cibles au seuil B1/B2 (HSK 4).
\begin{itemize}
    \item \textbf{不管是……或者是……} : Les élèves l'utilisent en \textbf{macro-structure} d'ouverture (thèse générale). Ils gèrent l'accord du sujet commun (\zh{获取资讯}) et la réponse \zh{都/也}.
    \item \textbf{是……的} : La contrainte « fait précis, passé, circoncis » force la distinction cognitive \textbf{Action nouvelle (SVO+了) vs Circonstance connue (是...的)}. Le piège \zh{了}+\zh{的} est évité par 90\% des groupes grâce à la relecture croisée.
    \item \textbf{比字句} : L'exigence de \textbf{degrés variés} (\zh{得多}, \zh{一点儿}, \zh{一些}, \zh{多了}, \zh{没有...那么...}) sort les élèves du \zh{比...更...} scolaire (HSK 2) vers la nuance HSK 4 (\zh{可靠一些}, \zh{方便多了}). L'usage spontané de \zh{没有} pour inverser la comparaison (radio vs journal) prouve la maîtrise.
    \item \textbf{Cohérence textuelle :} Les connecteurs logiques (\zh{但是}, \zh{虽然...但是...}, \zh{而且}, \zh{因此}) et la structure parallélique de la conclusion valident la compétence \textbf{Expression écrite} (niveau Bac / HSK 4).
    \item \textbf{Socioculturel :} L'ancrage camerounais (CRTV, WhatsApp, zones rurales, illettrisme personnes âgées) donne du sens à la comparaison, évitant l'exercice purement mécanique.
\end{itemize}
\end{aretenir}

\courssub{Prolongements et évaluation (Devoirs / Projet Module 5)}
\begin{methode}[Activités de transfert et d'approfondissement]
\begin{enumerate}
    \item \textbf{Publication numérique (Compétence numérique + Écrite) :} Publier l'article corrigé sur le blog du club de presse (WordPress / Compte WeChat Official Account) avec deux photos légendées en chinois : \zh{雅温得街头卖报纸} (Vendeur de journaux rue Yaoundé) vs \zh{学生看手机刷新闻} (Étudiants lisant news sur téléphone). Légende type : \zh{传统与新媒体的共存。} (Coexistence médias traditionnels/nouveaux).
    \item \textbf{Débat oral formel (Compétence orale - Expression continue/interaction) :} Sujet : \zh{「传统媒体会消失吗？」} (Les médias traditionnels vont-ils disparaître ?). Règles : 3 groupes (Pour, Contre, Jury). Chaque intervention doit contenir \textbf{au minimum} : 1 \zh{不管是}, 1 \zh{是...的} (témoignage personnel : \zh{我小时候是听广播长大的}), 1 \zh{比} avec degré. Évaluation sur grille : justesse structures (4 pts), vocabulaire médias (3 pts), prononciation/tons (2 pts), argumentation (1 pt).
    \item \textbf{Entraînement Bac Blanc — Thème (Français $\rightarrow$ Chinois) :} Traduire en chinois en mobilisant les 3 structures :
    \zh{« Que ce soit à Douala ou à Bafoussam, les jeunes préfèrent Internet pour s'informer. C'est sur TikTok que j'ai vu la vidéo de la finale hier. Internet est bien plus rapide que la télévision, mais la télévision est un peu plus sérieuse. Par conséquent, je conseille de regarder la télévision pour les grandes nouvelles et d'utiliser le portable pour les petites infos. »} \\
    \textbf{Corrigé attendu :} \zh{不管是在杜阿拉，还是在巴福萨姆，年轻人都喜欢用互联网获取资讯。这个决赛的视频，是我昨天在抖音上看到的。上网比看电视快得多，但是电视新闻比网上的消息严肃一些。因此，我建议：大新闻看电视，小消息查手机。}
\end{enumerate}
\end{methode}

\newpage

\coursec{Méthodes et exercices résolus}

\coursec{Leçon 35 — Grammaire : 不管是…….或者是……. ; 是。。。。。。的 ; 比字句}

\courssub{Observation guidée d'exemples authentiques}

\begin{textebox}[Dialogue : Choix du média et vérification de l'info (Contexte : Lycée de Yaoundé)]
\textbf{Personnages :} \zh{马文} (Élève, Terminale A) -- \zh{李娜} (Éléve, Terminale A)
\zh{马文：} 这个周末，\textbf{不管是看电视} \textbf{或者是刷手机}，我都要看喀麦隆队的比赛直播。\\
\zh{李娜：} 我\textbf{是上周在抖音上}看到预告\textbf{的}。电视信号有时候不好。\\
\zh{马文：} 没错，网络直播\textbf{比}电视\textbf{流畅得多}。\\
\zh{李娜：} \textbf{不管是}哪个平台，\textbf{都要}注意辨别假新闻。
\end{textebox}

\begin{experience}[Analyse inductive (Travail en binôme)]
Sur le dialogue ci-dessus, identifiez et soulignez :
\begin{enumerate}
    \item La structure exprimant l'indifférence du choix (« que ce soit... ou... »).
    \item La structure mettant l'accent sur une circonstance d'une action passée (« c'est... que... »).
    \item La structure de comparaison (« plus... que... »).
\end{enumerate}
Discutez : Quelle est la place du mot \zh{都} dans la première phrase ? Pourquoi \zh{了} est-il absent de la deuxième phrase ?
\end{experience}

\courssub{Structure 1 : 不管是……或者是…… (Indifférence du choix)}

\begin{definition}[Définition]
La structure \zh{不管是 A 或者是 B，主语 都/也 VP} exprime que \textbf{quelle que soit l'alternative choisie parmi A et B}, le résultat, l'attitude ou l'action du sujet principal reste inchangé. Elle correspond au français « Que ce soit A ou B, ... » ou « Qu'il s'agisse de A ou de B, ... ».
\end{definition}

\begin{propriete}[Schéma et Règles de construction]
\begin{center}
\begin{tabularx}{\linewidth}{|X|X|X|}
\hline
\rowcolor{popblueL} \textbf{Composant} & \textbf{Position} & \textbf{Nature / Contraintes} \\
\hline
\zh{不管是} & Début de la proposition subordonnée & Marque le concessionnel universel. \\
\hline
A / B & Après \zh{是} / après \zh{或者是} & Noms, pronoms, verbes, adjectifs, prépositionnels (\zh{在}...). \textbf{Répéter la préposition} si A et B sont des compléments de lieu/temps/moyen. \\
\hline
\zh{或者是} & Entre A et B & « Ou bien ». \textbf{Interdit :} \zh{还是} (réservé aux questions alternatives). \\
\hline
\zh{，} & Fin de la subordonnée & Virgule obligatoire avant la principale. \\
\hline
\zh{主语 + 都 / 也} & Début de la principale & \zh{都} (totalité, insistance) ou \zh{也} (addition). \textbf{Obligatoire}. \\
\hline
VP (Prédicat) & Fin & Verbe / Adjectif + Complément. Doit convenir sémantiquement à A et B. \\
\hline
\end{tabularx}
\end{center}
\end{propriete}

\begin{methode}[Construire une phrase avec \zh{不管是……或者是……}]
\begin{enumerate}
    \item \textbf{Identifier les alternatives A et B} (noms, verbes, prépositionnels). Ex: \zh{电视} / \zh{手机} ; \zh{在雅温得} / \zh{在杜阿拉}.
    \item \textbf{Écrire la concession} : \zh{不管是} + A + \zh{或者是} + B + \zh{，}. (Répéter la préposition : \zh{不管是在雅温得或者是在杜阿拉}).
    \item \textbf{Écrire la conséquence} : Sujet + \zh{都/也} + Prédicat.
    \item \textbf{Vérifier} : \zh{都/也} présent ? Prédicat compatible avec les deux alternatives ?
\end{enumerate}
\end{methode}

\begin{exemplebox}[Exemples variés]
\begin{enumerate}
    \item \zh{不管是学生或者是老师，学习中文都不容易。} (Bùguǎn shì xuéshēng huòzhě shì lǎoshī, xuéxí Zhōngwén dōu bù róngyì.) --- « Que l'on soit étudiant ou professeur, apprendre le chinois n'est pas facile. » (Alternatives : noms).
    \item \zh{不管是在家里或者是在学校，我们都要遵守规则。} (Bùguǎn shì zài jiālǐ huòzhě shì zài xuéxiào, wǒmen dōu yào zūnshǒu guīzé.) --- « Que ce soit à la maison ou à l'école, nous devons respecter les règles. » (Alternatives : prépositionnels \zh{在}... \rightarrow répétition de \zh{在}).
    \item \zh{不管是看报纸或者是看视频，都要核实来源。} (Bùguǎn shì kàn bàozhǐ huòzhě shì kàn shìpín, dōu yào héshí lái yuán.) --- « Que ce soit lire le journal ou regarder des vidéos, il faut vérifier la source. » (Alternatives : verbes. Sujet omis/général).
\end{enumerate}
\end{exemplebox}

\begin{exoresolu}[Application : Médias et géographie camerounaise]
\textbf{Énoncé :} Traduisez en chinois en utilisant \zh{不管是……或者是……}.
\begin{enumerate}
    \item Que ce soit WeChat ou les SMS, il a déjà envoyé l'info. (\zh{微信 Wēixìn}, \zh{短信 duǎnxìn}, \zh{发信息 fā xìnxī}, \zh{已经 yǐjīng}).
    \item Que ce soit à Buea ou à Douala, la 4G est rapide. (\zh{布埃亚 Bù'āiyà}, \zh{杜阿拉 Dùālā}, \zh{4G网络 4G wǎngluò}).
    \item Que ce soit lire un journal papier ou écouter la radio, il faut être critique. (\zh{纸质报纸 zhǐzhì bàozhǐ}, \zh{收音机 shōuyīnjī}, \zh{听 tīng}, \zh{保持批判性 bǎochí pīpànxìng}).
\end{enumerate}
\tcblower
\textbf{Solution détaillée :}
\begin{enumerate}
    \item \zh{不管是微信或者是短信，他已经发信息了。} (Alternatives nominales, action accomplie \zh{已经...了}).
    \item \zh{不管是在布埃亚或者是在杜阿拉，4G网络都很快。} (\textbf{Répétition de \zh{在}}).
    \item \zh{不管是看纸质报纸或者是听收音机，都要保持批判性。} (Alternatives verbales, sujet générique omis).
\end{enumerate}
\end{exoresolu}

\courssub{Structure 2 : 是……的 (Focus sur action passée : temps, lieu, manière, but...)}

\begin{definition}[Définition]
La construction \zh{是……的} (souvent appelée « cleft sentence » ou phrase de mise en relief) sert à \textbf{mettre l'accent sur une circonstance} (temps, lieu, moyen, compagnon, but, manière) d'une action \textbf{déjà accomplie}. Le verbe principal perd son \zh{le} (aspect accompli) au profit du \zh{de} final qui clôture la structure.
\end{definition}

\begin{propriete}[Schéma, Ordre des mots et Cas particuliers]
\textbf{Schéma de base :} \zh{主语 + 是 + [Élément mis en relief] + 动词 + 的。} (Objet direct \textbf{après le verbe}).
\begin{center}
\begin{tabularx}{\linewidth}{|X|X|X|}
\hline
\rowcolor{popblueL} \textbf{Focus (Élément mis en relief)} & \textbf{Exemple (Caractères + Pinyin)} & \textbf{Traduction / Note} \\
\hline
Temps & \zh{他是去年来喀麦隆的。} (Tā shì qùnián lái Kāmàilóng de.) & C'est l'an dernier qu'il est venu au Cameroun. \zh{了} disparu. \\
\hline
Lieu / Moyen & \zh{她是在网上买的手机。} (Tā shì zài wǎngshàng mǎi de shǒujī.) & C'est en ligne qu'elle a acheté le téléphone. Objet \zh{手机} après \zh{买}. \\
\hline
Compagnon / Manière & \zh{我们是一起去看球的。} (Wǒmen shì yīqǐ qù kàn qiú de.) & C'est ensemble que nous sommes allés voir le match. \\
\hline
But & \zh{我是为了学中文来中国的。} (Wǒ shì wèile xué Zhōngwén lái Zhōngguó de.) & C'est pour apprendre le chinois que je suis venu en Chine. \\
\hline
\multicolumn{3}{|l|}{\textbf{Cas particulier : Focus sur l'OBJET}} \\
\hline
Objet & \zh{这是我买的书。} (Zhè shì wǒ mǎi de shū.) & « C'est CE LIVRE que j'ai acheté. » Structure : \zh{是 + Objet + Verbe + 的 + Nom}. Différent de la mise en relief de circonstance. \\
\hline
\end{tabularx}
\end{center}
\end{propriete}

\begin{attention}[Pièges majeurs pour francophones (Structure \zh{是...的})]
\begin{enumerate}
    \item \textbf{Garder \zh{了} :} \zh{*我是昨天买了书的。} $\rightarrow$ \textbf{Faux}. \zh{了} (aspect) et \zh{的} (focus passé) sont en distribution complémentaire. \textbf{Correct :} \zh{我是昨天买书的} (focus temps) ou \zh{我是昨天买的书} (focus objet = « le livre que j'ai acheté hier »).
    \item \textbf{Déplacer l'objet dans la structure :} \zh{*我是书买的。} $\rightarrow$ Sens : « C'est LE LIVRE (et pas autre chose) que j'ai acheté » (Focus objet). Si le focus est le temps (\zh{昨天}), l'objet reste après le verbe : \zh{我是昨天买书的}.
    \item \textbf{Négation :} Utiliser \zh{不是...的}. \zh{我不是昨天买的书。} (Ce n'est pas hier que j'ai acheté le livre). \textbf{Jamais} \zh{没是...的} ni \zh{是...没的}.
\end{enumerate}
\end{attention}

\begin{methode}[Transformer une phrase passée en phrase de focus \zh{是...的}]
\begin{enumerate}
    \item Vérifier que l'action est \textbf{accomplie} (contexte passé, \zh{了} présent ou implicite).
    \item Identifier l'élément à souligner (Temps ? Lieu \zh{在}... ? Moyen \zh{坐}... ? But \zh{为了}... ?).
    \item Construire : \zh{Sujet + 是 + [Élément] + Verbe (+ Complément de résultat/direction) + 的}.
    \item \textbf{Placer l'objet :} S'il y a un objet direct et que le focus n'est PAS l'objet $\rightarrow$ Objet \textbf{après le verbe}. Si focus = Objet $\rightarrow$ Structure \zh{是 + Objet + V + 的 (+ Nom)}.
    \item Supprimer \zh{了} s'il était présent.
\end{enumerate}
\end{methode}

\begin{exoresolu}[Transformation : Focus sur circonstances passées]
\textbf{Énoncé :} Transformez pour insister sur l'élément \textbf{souligné}.
\begin{enumerate}
    \item Il a appris la nouvelle \textbf{hier}. (\zh{他昨天知道了这个消息。})
    \item Nous avons regardé le match \textbf{sur le téléphone}. (\zh{我们在手机上看了比赛。})
    \item Elle est allée à Douala \textbf{en avion}. (\zh{她坐飞机去杜阿拉了。})
    \item Je n'ai \textbf{PAS} lu ce journal \textbf{hier}. (\zh{我昨天没看这份报纸。} $\rightarrow$ Focus sur \zh{昨天} en niant).
\end{enumerate}
\tcblower
\textbf{Solution détaillée :}
\begin{enumerate}
    \item Focus Temps \zh{昨天} : \zh{他是昨天知道这个消息的。} (\emph{Tā shì zuótiān zhīdào zhège xiāoxi de.}) \textit{Verbe \zh{知道} (état/résultat), pas de \zh{了}.}
    \item Focus Lieu/Moyen \zh{在手机上} : \zh{我们是在手机上看比赛的。} (\emph{Wǒmen shì zài shǒujī shàng kàn bǐsài de.}) \textit{Objet \zh{比赛} après verbe.}
    \item Focus Moyen \zh{坐飞机} : \zh{她是坐飞机去杜阿拉的。} (\emph{Tā shì zuò fēijī qù Dùālā de.}) \textit{Complément de direction \zh{去杜阿拉} gardé après verbe.}
    \item Négation Focus Temps : \zh{我不是昨天看这份报纸的。} (\emph{Wǒ bù shì zuótiān kàn zhè fèn bàozhǐ de.}) \textit{Sens : C'est un autre jour.}
\end{enumerate}
\end{exoresolu}

\courssub{Structure 3 : La phrase comparative en 比 (比字句)}

\begin{definition}[Définition]
La structure \zh{A 比 B + Adjectif/Verbe d'état} exprime la \textbf{supériorité} de A sur B pour une qualité donnée. C'est l'équivalent principal de « A est plus [adj] que B ». L'ordre \zh{A 比 B} est \textbf{immuable}.
\end{definition}

\begin{propriete}[Tableau récapitulatif des comparatifs (Niveau HSK 3-4)]
\begin{center}
\begin{tabularx}{\linewidth}{|X|X|X|}
\hline
\rowcolor{popblueL} \textbf{Type} & \textbf{Structure} & \textbf{Exemple (Caractères + Pinyin + Traduction)} \\
\hline
Supériorité simple & \zh{A 比 B Adj} & \zh{中国比喀麦隆大。} (Zhōngguó bǐ Kāmàilóng dà.) --- La Chine est plus grande que le Cameroun. \\
\hline
Supériorité quantifiée & \zh{A 比 B Adj + Degré} & \zh{中国比喀麦隆大得多 / 大多了 / 大两万公里。} (Degré \textbf{APRÈS} l'adjectif). \\
\hline
Infériorité (Négation standard) & \zh{A 没有 B 这么/那么 Adj} & \zh{喀麦隆没有中国那么大。} (Kāmàilóng méiyǒu Zhōngguó nàme dà.) --- Le Cameroun n'est pas aussi grand que la Chine. \textbf{Pas de \zh{比}}. \\
\hline
Égalité & \zh{A 跟/和 B 一样 Adj} & \zh{这手机跟那手机一样贵。} (Zhège shǒujī gēn nàge shǒujī yīyàng guì.) \\
\hline
Supériorité (Choix/Préférence) & \zh{跟 B 比(起来)，A 更/更加 Adj} & \zh{跟报纸比(起来)，视频更直观。} (Gēn bàozhǐ bǐqǐlái, shìpín gèng zhíguān.) \\
\hline
\end{tabularx}
\end{center}
\end{propriete}

\begin{tikzpicture}[node distance=1.2cm, every node/.style={font=\small}]
\node[draw=popblue, thick, rounded corners, fill=popblueL, align=center] (A) {A \zh{比} B};
\node[draw=poporange, thick, rounded corners, fill=poporangeL, align=center, right=of A] (Adj) {Adjectif / Verbe d'état};
\node[draw=popgreen, thick, rounded corners, fill=popgreenL, align=center, right=of Adj] (Deg) {Degré (\zh{一点儿}, \zh{得多}, \zh{三岁}...)};
\draw[-Stealth, thick, popdark] (A) -- (Adj);
\draw[-Stealth, thick, popdark] (Adj) -- (Deg);
\node[align=center, text=popred, below=0.5cm of Adj] {\textbf{ORDRE FIXE : A 比 B Adj Degré}};
\end{tikzpicture}

\begin{methode}[Construire et manipuler le comparatif \zh{比}]
\begin{enumerate}
    \item \textbf{Ordre immuable :} \zh{A + 比 + B + Adjectif}. (A > B).
    \item \textbf{Degré (Quantification) :} \textbf{Toujours APRÈS l'adjectif} : \zh{一点儿 / 一些 / 得多 / 多了 / 两岁 / 三厘米}. \textit{Jamais avant l'adjectif.}
    \item \textbf{Négation « Pas aussi... que » :} \zh{A + 没有 + B + 这么/那么 + Adj}. \textit{Interdiction formelle de \zh{比} dans cette structure standard.}
    \item \textbf{Comparatif d'égalité :} \zh{A + 跟/和 + B + 一样 + Adj}.
    \item \textbf{Verbes d'action (coût, durée, effort) :} \zh{做这道题比做那道题费时间。} (Cette exo prend plus de temps que celle-là).
\end{enumerate}
\end{methode}

\begin{exemplebox}[Exemples contextuels Cameroun/Chine]
\begin{enumerate}
    \item \zh{雅温得比北京热得多。} (Yǎwēndé bǐ Běijīng rè de duō.) --- Yaoundé est beaucoup plus chaud que Pékin.
    \item \zh{看视频比看报纸轻松一点儿。} (Kàn shìpín bǐ kàn bàozhǐ qīngsōng yīdiǎnr.) --- Regarder des vidéos est un peu plus facile que lire le journal. (Sujets verbaux).
    \item \zh{这个APP没有那个APP那么难用。} (Zhège APP méiyǒu nàge APP nàme nányòng.) --- Cette appli n'est pas aussi difficile à utiliser que celle-là.
    \item \zh{喀麦隆足球队跟中国队一样强。} (Kāmàilóng zúqiúduì gēn Zhōngguó duì yīyàng qiáng.) --- L'équipe du Cameroun est aussi forte que l'équipe de Chine.
\end{enumerate}
\end{exemplebox}

\begin{exoresolu}[Examen blanc : Comparatifs médias et géographie]
\textbf{Énoncé :}
\begin{enumerate}
    \item Traduisez : « La télévision est plus intéressante que la radio. » (\zh{有意思 yǒu yìsi}).
    \item Traduisez : « Mon téléphone est un peu plus rapide que le tien. » (\zh{快 kuài}, \zh{一点儿 yīdiǎnr}).
    \item Traduisez : « Douala n'est pas aussi froid que Pékin en janvier. » (\zh{冷 lěng}, \zh{一月 yīyuè}).
    \item Transformez en phrase avec \zh{比} : \zh{这个网站有意思。那个网站没意思。}
    \item Corrigez l'erreur : \zh{*我比他高两岁。} (Âge : utiliser \zh{大}).
\end{enumerate}
\tcblower
\textbf{Solution détaillée :}
\begin{enumerate}
    \item \zh{电视比收音机有意思。} (\emph{Diànshì bǐ shōuyīnjī yǒu yìsi.})
    \item \zh{我的手机比你的手机快一点儿。} (\emph{Wǒ de shǒujī bǐ nǐ de shǒujī kuài yīdiǎnr.}) \textit{Degré \zh{一点儿} APRÈS \zh{快}.}
    \item \zh{杜阿拉一月没有北京那么冷。} (\emph{Dùālā yīyuè méiyǒu Běijīng nàme lěng.}) \textit{Structure \zh{没有...那么}.}
    \item \zh{这个网站比那个网站有意思。} (\emph{Zhège wǎngzhàn bǐ nàge wǎngzhàn yǒu yìsi.})
    \item \textbf{Erreur :} Degré \zh{两岁} avant l'adjectif + mauvais adjectif (\zh{高} = taille). \textbf{Correct :} \zh{我比他大两岁。} (\emph{Wǒ bǐ tā dà liǎng suì.})
\end{enumerate}
\end{exoresolu}

\courssub{Comparaison avec le français et pièges des francophones (Synthèse)}

\begin{attention}[Les 5 erreurs qui coûtent des points au Bac (Récapitulatif)]
\begin{enumerate}
    \item \textbf{\zh{比} + Degré : Position.} \zh{*我比他高两岁} $\rightarrow$ \zh{我比他大两岁}. Règle : \zh{A 比 B Adj + Degré}.
    \item \textbf{Négation comparatif : \zh{比} + \zh{不}.} \zh{*我不比他高} (Ambigu). \textbf{Correct :} \zh{我没有他那么高} (Infériorité/Égalité). \zh{不比} existe mais signifie « ne pas comparer » ou « n'être pas meilleur que » (ex: \zh{这不比那个好}).
    \item \textbf{\zh{是...的} : \zh{了} résiduel ou Objet déplacé.} \zh{*我是昨天买了书的} (Double passé). \zh{*我是书买的} (Focus objet involontaire). \textbf{Correct :} \zh{我是昨天买书的} (Focus temps) / \zh{我是昨天买的书} (Focus objet = « le livre que j'ai acheté hier »).
    \item \textbf{\zh{不管是...} : Oubli de \zh{都/也}.} \zh{*不管是手机或者是电脑，我看新闻。} $\rightarrow$ \zh{我\underline{都}看新闻}. Sans \zh{都}, structure « Topic-Comment » brisée, sens incomplet.
    \item \textbf{\zh{或者是} vs \zh{还是}.} \zh{还是} = Choix alternatif dans une \textbf{QUESTION} (\zh{你喝茶还是咖啡?}). \zh{或者是} / \zh{或} = Énumération / Indifférence dans une \textbf{AFFIRMATION} (\zh{不管是茶或者是咖啡...}). \textbf{Piège Bac :} \zh{不管是茶还是咖啡} $\rightarrow$ \textbf{Faute lourde}.
\end{enumerate}
\end{attention}

\begin{propriete}[Tableau de correspondance Français $\leftrightarrow$ Chinois (Thème / Version)]
\begin{center}
\begin{tabularx}{\linewidth}{|X|X|X|}
\hline
\rowcolor{popblueL} \textbf{Français (Source)} & \textbf{Chinois (Cible)} & \textbf{Point de vigilance} \\
\hline
« Que ce soit A ou B, ... » & \zh{不管是 A 或者是 B，都/也...} & \zh{都/也} obligatoire. \zh{或者是} pas \zh{还是}. \\
\hline
« C'est [Hier/À Paris/En avion] que j'ai fait... » (Insistance passée) & \zh{是 [Hier/À Paris/En avion] V 的} & \textbf{Pas de \zh{了}}. Objet après V. Négation : \zh{不是...的}. \\
\hline
« A est plus [Adj] que B » & \zh{A 比 B Adj} & Ordre fixe A > B. \\
\hline
« A est [beaucoup/un peu] plus [Adj] que B » & \zh{A 比 B Adj 得多 / 一点儿} & Degré \textbf{après} Adj. \\
\hline
« A n'est pas aussi [Adj] que B » & \zh{A 没有 B 那么/这么 Adj} & \textbf{Jamais \zh{比}}. \zh{那么} (loin) / \zh{这么} (proche/présent). \\
\hline
« A est aussi [Adj] que B » & \zh{A 跟 B 一样 Adj} & \zh{跟/和...一样}. \\
\hline
\end{tabularx}
\end{center}
\end{propriete}

\courssub{Exercices d'application et de transformation}

\begin{exercice}[Exercice 1 : Choix de structure (QCM / Complétion)]
Complétez avec la structure adaptée (\zh{不管是...或者是...}, \zh{是...的}, \zh{比}, \zh{没有...那么}, \zh{跟...一样}).
\begin{enumerate}
    \item \_\_\_\_\_\_ 电视 \_\_\_\_\_\_ 手机，年轻人 \_\_\_\_\_\_ 看新闻。 (Que ce soit TV ou téléphone, les jeunes regardent les infos.)
    \item 这个视频 \_\_\_\_\_\_ 那个视频 \_\_\_\_\_\_ 有意思。 (Cette vidéo est plus intéressante que celle-là.)
    \item 他 \_\_\_\_\_\_ 2020年 \_\_\_\_\_\_ 买 \_\_\_\_\_\_ 车。 (C'est en 2020 qu'il a acheté la voiture.)
    \item 喀麦隆 \_\_\_\_\_\_ 中国 \_\_\_\_\_\_ 大。 (Le Cameroun n'est pas aussi grand que la Chine.)
    \item 学中文 \_\_\_\_\_\_ 学英文 \_\_\_\_\_\_ 难。 (Apprendre le chinois est aussi difficile qu'apprendre l'anglais.)
\end{enumerate}

\end{exercice}

\begin{exercice}[Exercice 2 : Transformation de phrases]
Transformez les phrases suivantes selon l'indication.
\begin{enumerate}
    \item \zh{我在雅温得见到他了。} (Focus sur le lieu \zh{在雅温得} $\rightarrow$ \zh{是...的})
    \item \zh{这个手机比那个手机贵五百元。} (Négation : « Pas aussi cher » $\rightarrow$ \zh{没有...那么})
    \item \zh{无论你是男生还是女生，都要努力。} (Remplacez \zh{无论...还是...} par \zh{不管是...或者是...}).
    \item \zh{我看了这本书。} (Focus sur l'objet \zh{这本书} $\rightarrow$ \zh{是...的} structure focus objet).
    \item \zh{北京比上海冷。} (Quantifiez : « de 5 degrés » $\rightarrow$ \zh{五度}).
\end{enumerate}

\end{exercice}

\begin{exercice}[Exercice 3 : Traduction Thème (Français $\rightarrow$ Chinois)]
Traduisez en chinois (caractères + pinyin).
\begin{enumerate}
    \item Que ce soit par WeChat ou par email, le professeur a déjà envoyé le devoir. (\zh{微信}, \zh{邮箱 yóuxiāng}, \zh{老师 lǎoshī}, \zh{作业 zuòyè}, \zh{发送 fāsòng}).
    \item C'est l'année dernière que j'ai commencé à utiliser TikTok (Douyin). (\zh{开始 kāishǐ}, \zh{用 yòng}, \zh{抖音 Dǒuyīn}).
    \item Yaoundé est beaucoup plus chaud que Pékin en décembre. (\zh{十二月 shíèryuè}, \zh{热 rè}, \zh{得多 de duō}).
    \item Mon portable n'est pas aussi rapide que le tien. (\zh{手机 shǒujī}, \zh{快 kuài}).
    \item Regarder les infos à la télé ou sur le portable, il faut toujours vérifier la source. (\zh{看新闻 kàn xīnwén}, \zh{电视 diànshì}, \zh{核实 héshí}, \zh{来源 lái yuán}).
\end{enumerate}

\end{exercice}

\begin{exercice}[Exercice 4 : Traduction Version (Chinois $\rightarrow$ Français)]
Traduisez en français naturel.
\begin{enumerate}
    \item \zh{不管是看纸质报纸或者是刷短视频，我们都要独立思考。}
    \item \zh{我是上个月在朋友圈看到这个广告的。}
    \item \zh{现在的网速比十年前快得多。}
    \item \zh{传统媒体没有新媒体那么快。}
    \item \zh{不管你信不信，这个消息是真的。} (Bonus : structure \zh{不管} + Verbe).
\end{enumerate}

\end{exercice}

\begin{exercice}[Exercice 5 : Production écrite guidée (60-80 caractères)]
Rédigez un court paragraphe sur « Mon rapport à l'information » en respectant \textbf{strictement} les 3 contraintes :
\begin{itemize}
    \item 1 phrase avec \zh{比} (comparatif quantifié : \zh{得多} ou \zh{一点儿}).
    \item 1 phrase avec \zh{是...的} (raconter comment/quand vous avez appris une nouvelle importante).
    \item 1 phrase avec \zh{不管是...或者是...} (exprimer votre exigence de vérité).
\end{itemize}
Utilisez les connecteurs : \zh{首先}, \zh{其次}, \zh{总之}.
\end{exercice}

\courssub{Corrigés des exercices}

\begin{corrige}[Exercice 1]
\begin{enumerate}
    \item \zh{不管是} 电视 \zh{或者是} 手机，年轻人 \zh{都} 看新闻。
    \item 这个视频 \zh{比} 那个视频 \zh{有意思}。 (ou \zh{更有意思}).
    \item 他 \zh{是} 2020年 \zh{买} \zh{的} 车。
    \item 喀麦隆 \zh{没有} 中国 \zh{那么} 大。
    \item 学中文 \zh{跟} 学英文 \zh{一样} 难。
\end{enumerate}
\end{corrige}

\begin{corrige}[Exercice 2]
\begin{enumerate}
    \item \zh{我是在雅温得见到他的。} (Ou \zh{见到他}).
    \item \zh{那个手机没有这个手机那么贵。} (Inversion A/B nécessaire pour garder le sens « celui-ci est plus cher »).
    \item \zh{不管是男生或者是女生，都要努力。}
    \item \zh{这是我看的书。} (Structure Focus Objet : \zh{是 + Objet + V + 的}).
    \item \zh{北京比上海冷五度。}
\end{enumerate}
\end{corrige}

\begin{corrige}[Exercice 3 : Thème]
\begin{enumerate}
    \item \zh{不管是微信或者是邮箱，老师已经发送作业了。} (Bùguǎn shì Wēixìn huòzhě shì yóuxiāng, lǎoshī yǐjīng fāsòng zuòyè le.)
    \item \zh{我是去年开始用抖音的。} (Wǒ shì qùnián kāishǐ yòng Dǒuyīn de.)
    \item \zh{雅温得比北京十二月热得多。} (Yǎwēndé bǐ Běijīng shíèryuè rè de duō.)
    \item \zh{我的手机没有你的手机那么快。} (Wǒ de shǒujī méiyǒu nǐ de shǒujī nàme kuài.)
    \item \zh{不管是看电视新闻或者是看手机新闻，都要核实来源。} (Bùguǎn shì kàn diànshì xīnwén huòzhě shì kàn shǒujī xīnwén, dōu yào héshí lái yuán.)
\end{enumerate}
\end{corrige}

\begin{corrige}[Exercice 4 : Version]
\begin{enumerate}
    \item « Que ce soit lire le journal papier ou faire défiler des vidéos courtes, nous devons tous penser de manière indépendante. »
    \item « C'est le mois dernier que j'ai vu cette publicité sur les "Moments" (WeChat). »
    \item « La vitesse du net actuelle est bien plus rapide qu'il y a dix ans. »
    \item « Les médias traditionnels ne sont pas aussi rapides que les nouveaux médias. »
    \item « Que tu y croies ou non, cette nouvelle est vraie. » (\zh{不管你信不信} = « Que tu croies ou pas », structure \zh{不管} + Verbe affirmatif/négatif).
\end{enumerate}
\end{corrige}

\begin{corrige}[Exercice 5 : Production écrite (Modèle de référence)]
\zh{首先，看视频比看文字轻松得多。其次，我是昨天在抖音上看到喀麦隆队赢球的消息。总之，不管是短视频或者是长文章，我都希望内容是真实的。}
\textbf{Analyse grammaticale du modèle :}
\begin{enumerate}
    \item \zh{看视频比看文字轻松得多} : Sujets verbaux (\zh{看视频}, \zh{看文字}), \zh{比}, Adj \zh{轻松}, Degré \zh{得多} (position correcte). \checkmark
    \item \zh{我是昨天在抖音上看到...的} : \zh{是...的} focus double (Temps \zh{昨天} + Lieu/Moyen \zh{在抖音上}). Verbe \zh{看到} (résultat). Pas de \zh{了}. \checkmark
    \item \zh{不管是短视频或者是长文章，我都希望...} : Structure complète, \zh{都} présent, alternatives nominales. \checkmark
\end{enumerate}
\end{corrige}

\courssub{Vocabulaire clé de la leçon (HSK 3-4)}

\begin{center}
\begin{tabularx}{\linewidth}{|X|X|X|X|}
\hline
\rowcolor{popblueL} \textbf{Caractères} & \textbf{Pinyin} & \textbf{Nature} & \textbf{Traduction} \\
\hline
不管是 & bùguǎn shì & Conjonction & Quoi que ce soit / Que ce soit \\
\hline
或者是 & huòzhě shì & Conjonction & Ou bien (affirmation) \\
\hline
还是 & hái shi & Conjonction & Ou bien (question/choix) \\
\hline
都 & dōu & Adv. & Tous / Tous deux (obligatoire avec 不管) \\
\hline
是...的 & shì...de & Structure & C'est... que... (focus passé) \\
\hline
比 & bǐ & Préposition & Comparatif « plus que » \\
\hline
没有 & méiyǒu & Verbe & N'avoir pas / Comparatif « pas aussi... que » \\
\hline
那么 & nà me & Adv. & Tellement / Si (loin) \\
\hline
这么 & zhè me & Adv. & Si / Tellement (proche) \\
\hline
一样 & yī yàng & Adj/Adv & Pareil / Identique \\
\hline
得多 & de duō & Complément & Beaucoup plus (degré fort) \\
\hline
一点儿 & yī diǎnr & Complément & Un peu (degré faible) \\
\hline
核实 & hé shí & Verbe & Vérifier (information) \\
\hline
来源 & lái yuán & Nom & Source (info) \\
\hline
真实 & zhēn shí & Adj & Vrai / Authentique \\
\hline
假 & jiǎ & Adj & Faux \\
\hline
流畅 & liú chàng & Adj & Fluide (connexion, vidéo) \\
\hline
直观 & zhí guān & Adj & Intuitif / Visuel \\
\hline
批判性 & pī pàn xìng & Nom & Esprit critique \\
\hline
\end{tabularx}
\end{center}

\begin{savaistu}[Socioculturel : Paysage médiatique Cameroun vs Chine]
\textbf{Cameroun :} Coexistence médias traditionnels (CRTV - Radio/Télé publique, presse écrite \zh{报纸} comme \zh{Cameroon Tribune}, \zh{Mutations}) et nouveaux médias (Facebook, WhatsApp, chaînes YouTube locales). Problématique : \zh{假新闻} (Fake news), \zh{网络谣言} (Rumeurs en ligne), coût de la data (\zh{流量 liàngliang}).
\textbf{Chine :} Écosystème numérique intégré \zh{微信 Wēixìn} (WeChat : messagerie, paiement, réseaux sociaux, mini-programmes), \zh{抖音 Dǒuyīn} (TikTok chinois), \zh{微博 Wēibó} (Twitter-like), \zh{今日头条 Jīnrì Tóutiáo} (Agrégateur d'infos IA). Contrôle strict (\zh{审核 shēnhé}), grande rapidité (\zh{5G/光纤 guāngxiān}), paiement sans espèces généralisé (\zh{扫码支付 sǎomǎ zhīfù}).
\textbf{Point commun :} Transition vers le \zh{新媒体 xīn méitǐ} (Nouveaux médias), déclin de la \zh{纸媒 zhǐ méi} (Presse papier), nécessité d'\zh{媒介素养 méijiè sùyǎng} (Éducation aux médias).
\end{savaistu}

\newpage

\coursec{Exercices}

\courssub{Je vérifie mes connaissances}

\begin{exercice}[QCM : Identification des structures]
Pour chaque phrase, cocher la seule proposition correcte.
\begin{enumerate}
    \item La structure qui exprime « quel que soit le choix entre A et B, le résultat est le même » est :
    \begin{itemize}
        \item[A)] 是……的
        \item[B)] 不管是……或者是……
        \item[C)] 比字句
        \item[D)] 不但……而且……
    \end{itemize}
    \item Dans une phrase comparative avec \zh{比} (Bǐ), l'adverbe de degré \zh{很} (hěn) :
    \begin{itemize}
        \item[A)] Se place toujours devant l'adjectif.
        \item[B)] Est interdit (on utilise \zh{更} gèng, \zh{得多} de duō, \zh{一点儿} yìdiǎnr).
        \item[C)] Est obligatoire.
        \item[D)] Se place après l'adjectif.
    \end{itemize}
    \item La structure \zh{是……的} (shì… de) sert principalement à :
    \begin{itemize}
        \item[A)] Comparer deux objets.
        \item[B)] Insister sur le temps, le lieu, la manière ou l'auteur d'une action \textbf{déjà accomplie}.
        \item[C)] Exprimer un choix alternatif.
        \item[D)] Nier une action passée.
    \end{itemize}
    \item Dans la phrase \zh{不管你是老师或者是学生，都要遵守规则。} (Bùguǎn nǐ shì lǎoshī huòzhě shì xuéshēng, dōu yào zūnshǒu guīzé.), le mot \zh{都} (dōu) :
    \begin{itemize}
        \item[A)] Est facultatif.
        \item[B)] Se place avant \zh{不管}.
        \item[C)] Est nécessaire dans la clause principale pour marquer l'universalité.
        \item[D)] Remplace \zh{也} (yě).
    \end{itemize}
\end{enumerate}
\end{exercice}

\begin{exercice}[Vrai ou Faux ? Justifiez en français]
Dire si l'affirmation est vraie (V) ou fausse (F) et corriger l'erreur grammaticale le cas échéant.
\begin{enumerate}
    \item \zh{手机比电脑便宜很。} (Shǒujī bǐ diànnǎo piányi hěn.)
    \item \zh{是昨天我看了这部电影的。} (Shì zuótiān wǒ kànle zhè bù diànyǐng de.)
    \item \zh{不管是下雨或者是刮风，我都去跑步。} (Bùguǎn shì xiàyǔ huòzhě shì guāfēng, wǒ dōu qù pǎobù.)
    \item \zh{这本书比那本书有意思得多了。} (Zhè běn shū bǐ nà běn shū yǒuyìsi de duō le.)
\end{enumerate}
\end{exercice}

\begin{exercice}[Vocabulaire : Les médias et les télécommunications]
Associer chaque terme chinois à sa définition française. Écrire la lettre correspondante.
\begin{center}
\begin{tabularx}{\linewidth}{|X|X|X|X|}
\hline
\rowcolor{popblueL} \textbf{Caractères} & \textbf{Pinyin} & \textbf{Nature} & \textbf{Définition / Traduction} \\
\hline
1. 电视台 & diànshìtái & n. & A. Réseau social / Media social \\
\hline
2. 社交媒体 & shèjiāo méitǐ & n. & B. Journal (quotidien) \\
\hline
3. 报纸 & bàozhǐ & n. & C. Chaîne de télévision / Station de TV \\
\hline
4. 直播 & zhíbō & v. & D. Diffuser en direct (live) \\
\hline
5. 网络 & wǎngluò & n. & E. Information / Nouvelles \\
\hline
6. 资讯 & zīxùn & n. & F. Réseau / Internet \\
\hline
7. 收音机 & shōuyīnjī & n. & G. Radio (poste de radio) \\
\hline
8. 记者 & jìzhě & n. & H. Journaliste / Reporter \\
\hline
\end{tabularx}
\end{center}
\end{exercice}

\begin{exercice}[Pinyin et Caractères : Reconnaissance]
Relier le pinyin (avec tons) au caractère correspondant.
\begin{center}
\begin{tabularx}{\linewidth}{|X|X|X|X|}
\hline
\rowcolor{poporangeL} \textbf{Pinyin} & \textbf{Caractère cible} & \textbf{Sens} & \textbf{Propositions (mélangées)} \\
\hline
bùguǎn & \zh{不管} & Quoi qu'il en soit / Peu importe & \zh{不管、比较、必须、媒体} \\
\hline
bǐjiào & \zh{比较} & Comparer / Relativement & \zh{比较、比赛、笔记、被子} \\
\hline
bìxū & \zh{必须} & Devoir / Falloir (obligation forte) & \zh{必须、比较、媒体、报纸} \\
\hline
méitǐ & \zh{媒体} & Média (support d'information) & \zh{媒体、每天、美丽、妹妹} \\
\hline
\end{tabularx}
\end{center}
\end{exercice}

\courssub{Je m'entraîne}

\begin{exercice}[Traduction : Français $\rightarrow$ Chinois]
Traduire les phrases suivantes en chinois (caractères + pinyin). Utiliser la structure imposée entre parenthèses.
\begin{enumerate}
    \item Que tu regardes la télévision ou que tu surfes sur Internet, les publicités sont partout. (\zh{不管是……或者是……})
    \item C'est par téléphone portable qu'il a payé l'addition hier. (\zh{是……的})
    \item La 5G est beaucoup plus rapide que la 4G. (\zh{比字句} + \zh{得多})
    \item Peu importe que ce soit une vraie nouvelle ou une fausse nouvelle, il faut vérifier. (\zh{不管是……或者是……} + \zh{都})
    \item Ce journaliste a écrit cet article à Yaoundé le mois dernier. (\zh{是……的} : insister sur le lieu et le temps)
\end{enumerate}
\end{exercice}

\begin{exercice}[Transformation de phrases]
Réécrire chaque phrase en utilisant la structure indiquée pour mettre l'accent sur l'élément souligné.
\begin{enumerate}
    \item \zh{我昨天在喀麦隆看了足球比赛。} (Insister sur le \textbf{temps} et le \textbf{lieu} $\rightarrow$ \zh{是……的})
    \item \zh{收音机不如电视直观。} (Transformer en phrase comparative avec \zh{比} : \textbf{La TV est plus intuitive que la radio})
    \item \zh{无论新闻是真还是假，我们要小心。} (Réécrire avec \zh{不管是……或者是……})
    \item \zh{他用中文写了这封信。} (Insister sur la \textbf{langue/instrument} $\rightarrow$ \zh{是……的})
    \item \zh{看报纸比刷手机慢。} (Ajouter un complément de degré \textbf{「 beaucoup plus »} $\rightarrow$ \zh{得多})
\end{enumerate}
\end{exercice}

\begin{exercice}[Texte à trous : Mots-outils]
Compléter le texte suivant sur les télécommunications au Cameroun avec les mots manquants : \zh{不管}、\zh{是}、\zh{的}、\zh{比}、\zh{都}、\zh{得多}、\zh{或者}、\zh{更}。
\begin{textebox}[L'évolution des médias au Cameroun]
\zh{现在} \zh{\_\_\_\_\_\_} \zh{是} \zh{城市} \zh{\_\_\_\_\_\_} \zh{是} \zh{农村}，\zh{手机} \zh{\_\_\_\_\_\_} \zh{是} \zh{智能机} \zh{\_\_\_\_\_\_} \zh{是} \zh{功能机}，\zh{人人} \zh{\_\_\_\_\_\_} \zh{有}。\zh{网络} \zh{信号} \zh{\_\_\_\_\_\_} \zh{以前} \zh{\_\_\_\_\_\_} \zh{好} \zh{\_\_\_\_\_\_} \zh{了}。\zh{获取} \zh{资讯} \zh{\_\_\_\_\_\_} \zh{过去} \zh{\_\_\_\_\_\_} \zh{方便} \zh{\_\_\_\_\_\_}。\zh{是} \zh{通过} \zh{微信} \zh{\_\_\_\_\_\_} \zh{朋友圈} \zh{看} \zh{新闻} \zh{\_\_\_\_\_\_} \zh{很多} \zh{年轻人} \zh{的} \zh{习惯}。
\end{textebox}
\end{exercice}

\begin{exercice}[Remise en ordre / Appariement]
Remettre les mots dans l'ordre pour former des phrases correctes. Écrire la phrase en caractères et en pinyin.
\begin{enumerate}
    \item \zh{的 / 这条 / 是 / 新闻 / 真 / 的 / 是} \quad (Structure \zh{是……的} : insister sur l'authenticité)
    \item \zh{比 / 电视 / 手机 / 方便 / 得多 / 是} \quad (Structure \zh{比} : le téléphone plus pratique que la TV)
    \item \zh{或者 / 是 / 不管 / 老人 / 年轻人 / 是 / 都 / 喜欢 / 看 / 短视频} \quad (Structure \zh{不管是……或者是……})
    \item \zh{在 / 是 / 我 / 家 / 的 / 看 / 电视 / 的} \quad (Structure \zh{是……的} : insister sur le lieu)
    \item \zh{报纸 / 环保 / 电子版 / 比 / 更 / 传统 / 是} \quad (Structure \zh{比} + \zh{更} : version électronique plus écologique)
\end{enumerate}
\end{exercice}

\courssub{Je me prépare au Bac}

\begin{exercice}[Type Bac — Compréhension de texte et questions de langue]
Lire le texte suivant et répondre aux questions en français.
\begin{textebox}[Le choix de l'information : TV vs Internet]
\zh{在喀麦隆，不管是住在雅温得的公务员，或者是住在杜阿拉的商人，现在都离不开手机。是十年前，大家主要通过电视台和报纸了解新闻的。可是现在，短视频平台比传统电视受欢迎得多。因为手机上网比买电视机便宜，而且随时随地都能看直播。是通过微信群，我昨天才知道了那条关于足球比赛的假新闻。所以，不管是真消息或者是假消息，我们都要学会辨别。}
\par \emph{(Zài Kāmàilóng, bùguǎn shì zhù zài Yǎwēndé de gōngwùyuán, huòzhě shì zhù zài Dūālā de shāngrén, xiànzài dōu líbukāi shǒujī. Shì shí nián qián, dàjiā zhǔyào tōngguò diànshìtái hé bàozhǐ liǎojiě xīnwén de. Kěshì xiànzài, duǎn shìpín píngtái bǐ chuántǒng diànshì shòuhuān de duō. Yīnwèi shǒujī shàngwǎng bǐ mǎi diànshījī piányi, érqiě suíshí suídì dōu néng kàn zhíbō. Shì tōngguò Wēixìn qún, wǒ zuótiān cái zhīdàole nà tiáo guānyú zúqiú bǐsài de jiǎ xīnwén. Suǒyǐ, bùguǎn shì zhēn xiāoxi huòzhě shì jiǎ xiāoxi, wǒmen dōu xuéhuì biànbié.)}
\end{textebox}
\textbf{Questions de compréhension :}
\begin{enumerate}
    \item Selon le texte, quels sont les deux profils de personnes mentionnés et dans quelles villes habitent-ils ?
    \item Quels étaient les médias principaux il y a dix ans ?
    \item Pourquoi les plateformes de vidéos courtes sont-elles plus populaires que la télévision traditionnelle aujourd'hui ? (Deux raisons)
    \item Comment le narrateur a-t-il appris la fausse nouvelle sur le match de football ?
\end{enumerate}
\textbf{Questions de langue :}
\begin{enumerate}
    \setcounter{enumi}{4}
    \item Identifier dans le texte \textbf{deux} occurrences de la structure \zh{不管是……或者是……} et les souligner.
    \item Identifier \textbf{trois} occurrences de la structure \zh{是……的} et préciser pour chacune l'élément mis en relief (temps, lieu, moyen, etc.).
    \item Relever \textbf{deux} phrases comparatives avec \zh{比} et indiquer le complément de degré utilisé (\zh{得多}, \zh{更}, etc.).
\end{enumerate}
\end{exercice}

\begin{exercice}[Type Bac — Thème et Version]
\textbf{A. Thème (Français $\rightarrow$ Chinois) :} Traduire en chinois (caractères + pinyin).
\begin{enumerate}
    \item Que je regarde la télévision ou que j'écoute la radio, j'aime suivre les informations. (Utiliser \zh{不管是……或者是……})
    \item C'est hier soir que j'ai appris cette nouvelle par Internet. (Utiliser \zh{是……的})
    \item Internet est beaucoup plus rapide que les journaux. (Utiliser \zh{比} + \zh{得多})
\end{enumerate}

\textbf{B. Version (Chinois $\rightarrow$ Français) :} Traduire en français correct.
\begin{enumerate}
    \setcounter{enumi}{3}
    \item \zh{不管是传统媒体还是新媒体，都有优点和缺点。}
    \item \zh{这篇报道是记者上周在现场采访写的。}
    \item \zh{看直播比看录播有意思多了。}
\end{enumerate}
\end{exercice}

\begin{exercice}[Type Bac — Expression écrite guidée]
\textbf{Sujet :} \emph{« Les réseaux sociaux ont remplacé la télévision comme source principale d'information pour les jeunes. »}
Rédiger un court paragraphe d'environ \textbf{80 à 100 caractères chinois} pour exprimer votre avis.
\textbf{Contraintes obligatoires :}
\begin{itemize}
    \item Utiliser \textbf{au moins deux} des trois structures étudiées : \zh{不管是……或者是……}、\zh{是……的}、\zh{比字句}.
    \item Utiliser au moins 5 mots du vocabulaire de la leçon (ex: \zh{社交媒体}、\zh{直播}、\zh{资讯}、\zh{真假}、\zh{辨别}、\zh{方便}、\zh{习惯}).
    \item Respecter l'ordre des mots (sujet + temps + lieu + verbe) et la place des mots-outils (\zh{都}、\zh{得多}、\zh{更}).
\end{itemize}
\begin{center}
\begin{tabularx}{\linewidth}{|X|}
\hline
\rowcolor{popgreenL} \textbf{Amorce suggérée :} \zh{现在的年轻人，\_\_\_\_\_\_……} \\
\hline
\end{tabularx}
\end{center}
\end{exercice}

\newpage

\coursec{Corrigés des exercices}

\courssub{Corrigés des exercices — Leçon 35}

\begin{corrige}[Exercice 1 — QCM : Identification des structures]
\begin{enumerate}
    \item \textbf{Réponse B}. \zh{不管是……或者是……} (bùguǎn shì… huòzhě shì…) exprime « quel que soit le choix entre A et B, le résultat est le même ». \zh{是……的} insiste sur une action passée ; \zh{比字句} compare ; \zh{不但……而且……} exprime l'addition (« non seulement… mais aussi »).
    \item \textbf{Réponse B}. Dans une phrase comparative en \zh{比}, l'adjectif ne peut pas être précédé de \zh{很} (hěn). On exprime le degré avec \zh{更} (gèng, « encore plus ») devant l'adjectif, ou \zh{得多} (de duō, « beaucoup plus ») / \zh{一点儿} (yìdiǎnr, « un peu plus ») après l'adjectif.
    \item \textbf{Réponse B}. La structure \zh{是……的} met en relief le temps, le lieu, la manière, le moyen ou l'auteur d'une action \textbf{déjà accomplie}. Elle ne sert ni à comparer, ni à exprimer un choix, ni à nier (la négation d'une action passée se fait avec \zh{没（有）}).
    \item \textbf{Réponse C}. Dans la clause principale qui suit \zh{不管是……或者是……}, l'adverbe \zh{都} (dōu) — ou \zh{也} (yě) — est \textbf{nécessaire} : il marque l'universalité du résultat (« dans tous les cas »). Il se place après le sujet de la clause principale, jamais avant \zh{不管}.
\end{enumerate}
\end{corrige}

\begin{corrige}[Exercice 2 — Vrai ou Faux ?]
\begin{enumerate}
    \item \textbf{Faux.} Dans une phrase comparative en \zh{比}, on ne peut pas utiliser \zh{很} devant l'adjectif. Correction : \zh{手机比电脑便宜得多。} (Shǒujī bǐ diànnǎo piányi de duō.) « Le téléphone portable est beaucoup moins cher que l'ordinateur. » On peut aussi dire \zh{手机比电脑便宜一点儿。}
    \item \textbf{Faux.} Dans la structure \zh{是……的}, le sujet se place normalement avant \zh{是}, et le \zh{的} final remplace le \zh{了} de l'accompli (on ne cumule pas \zh{了} après le verbe avec \zh{的} final). Correction : \zh{我是昨天看这部电影的。} (Wǒ shì zuótiān kàn zhè bù diànyǐng de.) « C'est hier que j'ai vu ce film. »
    \item \textbf{Vrai.} Structure correcte : \zh{不管是} + A + \zh{或者是} + B, sujet + \zh{都} + verbe. Exemple : \zh{不管是下雨或者是刮风，我都去跑步。} (Bùguǎn shì xiàyǔ huòzhě shì guāfēng, wǒ dōu qù pǎobù.) « Qu'il pleuve ou qu'il vente, je vais courir. »
    \item \textbf{Vrai.} \zh{比} + adjectif + \zh{得多} est correct ; le \zh{了} final marque ici le constat du changement. « Ce livre est bien plus intéressant que celui-là. » (On peut aussi écrire sans \zh{了} : \zh{有意思得多。})
\end{enumerate}
\end{corrige}

\begin{corrige}[Exercice 3 — Vocabulaire : médias et télécommunications]
\begin{itemize}
    \item 1 — C (\zh{电视台} diànshìtái : chaîne de télévision)
    \item 2 — A (\zh{社交媒体} shèjiāo méitǐ : réseaux sociaux)
    \item 3 — B (\zh{报纸} bàozhǐ : journal)
    \item 4 — D (\zh{直播} zhíbō : diffuser en direct)
    \item 5 — F (\zh{网络} wǎngluò : réseau, Internet)
    \item 6 — E (\zh{资讯} zīxùn : informations)
    \item 7 — G (\zh{收音机} shōuyīnjī : radio)
    \item 8 — H (\zh{记者} jìzhě : journaliste)
\end{itemize}
\textbf{Point de vigilance :} ne pas confondre \zh{报纸} (bàozhǐ, le journal papier) et \zh{报道} (bàodào, le reportage) ; \zh{记者} contient la clé \zh{讠} (parole), comme beaucoup de mots liés à la communication.
\end{corrige}

\begin{corrige}[Exercice 4 — Pinyin et caractères]
\begin{itemize}
    \item bùguǎn $\rightarrow$ \zh{不管} (« peu importe »). Piège : ne pas choisir \zh{必须} (bìxū) ni \zh{媒体} (méitǐ) : seul \zh{不管} porte les tons 4-3.
    \item bǐjiào $\rightarrow$ \zh{比较} (« comparer »). Pièges : \zh{比赛} (bǐsài, « compétition »), \zh{笔记} (bǐjì, « notes »), \zh{被子} (bèizi, « couette ») — mêmes premières syllabes proches, tons différents.
    \item bìxū $\rightarrow$ \zh{必须} (« devoir, falloir »). Attention au ton 4 sur bì, contrairement à bǐ (3e ton) de \zh{比较}.
    \item méitǐ $\rightarrow$ \zh{媒体} (« média »). Pièges homophones proches : \zh{每天} (měitiān, « chaque jour »), \zh{美丽} (měilì, « beau »), \zh{妹妹} (mèimei, « petite sœur »).
\end{itemize}
\textbf{Méthode :} vérifier d'abord les tons, puis les clés (radicaux) : \zh{媒} contient la clé \zh{女} (femme) à gauche ; \zh{较} contient la clé \zh{车} (véhicule).
\end{corrige}

\begin{corrige}[Exercice 5 — Traduction : Français $\rightarrow$ Chinois]
\begin{enumerate}
    \item \zh{不管你是看电视，或者是上网，广告都无处不在。} (Bùguǎn nǐ shì kàn diànshì, huòzhě shì shàngwǎng, guǎnggào dōu wúchùbùzài.) « Que tu regardes la télévision ou que tu surfes sur Internet, les publicités sont partout. » Noter le \zh{都} obligatoire dans la clause principale.
    \item \zh{他是昨天用手机付的账。} (Tā shì zuótiān yòng shǒujī fù de zhàng.) « C'est par téléphone portable qu'il a payé l'addition hier. » Ici \zh{的} se place entre le verbe et l'objet (\zh{付的账}) ; on peut aussi dire \zh{他是昨天用手机付账的。}
    \item \zh{5G比4G快得多。} (5G bǐ 4G kuài de duō.) « La 5G est beaucoup plus rapide que la 4G. » Pas de \zh{很} devant \zh{快}.
    \item \zh{不管是真新闻，或者是假新闻，我们都要核实。} (Bùguǎn shì zhēn xīnwén, huòzhě shì jiǎ xīnwén, wǒmen dōu yào héshí.) « Peu importe que ce soit une vraie ou une fausse nouvelle, il faut vérifier. »
    \item \zh{这位记者是上个月在雅温得写这篇文章的。} (Zhè wèi jìzhě shì shàng gè yuè zài Yǎwēndé xiě zhè piān wénzhāng de.) « C'est à Yaoundé, le mois dernier, que ce journaliste a écrit cet article. » L'élément mis en relief (\zh{上个月在雅温得}) se place entre \zh{是} et le verbe (temps avant lieu).
\end{enumerate}
\end{corrige}

\begin{corrige}[Exercice 6 — Transformation de phrases]
\begin{enumerate}
    \item \zh{我是昨天在喀麦隆看足球比赛的。} (Wǒ shì zuótiān zài Kāmàilóng kàn zúqiú bǐsài de.) « C'est hier, au Cameroun, que j'ai regardé le match de football. » Le \zh{了} disparaît ; temps avant lieu entre \zh{是} et le verbe.
    \item \zh{电视比收音机直观。} (Diànshì bǐ shōuyīnjī zhíguān.) « La télévision est plus intuitive que la radio. » On inverse l'ordre des termes : A + \zh{比} + B + adjectif.
    \item \zh{不管是真新闻，或者是假新闻，我们都要小心。} (Bùguǎn shì zhēn xīnwén, huòzhě shì jiǎ xīnwén, wǒmen dōu yào xiǎoxīn.) « Que la nouvelle soit vraie ou fausse, nous devons être prudents. »
    \item \zh{他是用中文写这封信的。} (Tā shì yòng Zhōngwén xiě zhè fēng xìn de.) « C'est en chinois qu'il a écrit cette lettre. » L'instrument (\zh{用中文}) est placé entre \zh{是} et le verbe.
    \item \zh{看报纸比刷手机慢得多。} (Kàn bàozhǐ bǐ shuā shǒujī màn de duō.) « Lire le journal est beaucoup plus lent que consulter son téléphone. » \zh{得多} se place après l'adjectif \zh{慢}.
\end{enumerate}
\end{corrige}

\begin{corrige}[Exercice 7 — Texte à trous]
Texte complété :
\par \zh{现在不管是城市或者是农村，手机不管是智能机或者是功能机，人人都有。网络信号比以前好得多了。获取资讯比过去更方便。很多年轻人的习惯是通过微信朋友圈看新闻的。}
\par \emph{(Xiànzài bùguǎn shì chéngshì huòzhě shì nóngcūn, shǒujī bùguǎn shì zhìnéngjī huòzhě shì gōngnéngjī, rénrén dōu yǒu. Wǎngluò xìnhào bǐ yǐqián hǎo de duō le. Huòqǔ zīxùn bǐ guòqù gèng fāngbiàn. Hěn duō niánqīngrén de xíguàn shì tōngguò Wēixìn péngyouquān kàn xīnwén de.)}
\par Traduction : « Aujourd'hui, que ce soit en ville ou à la campagne, que ce soit un smartphone ou un téléphone basique, tout le monde a un téléphone. Le signal du réseau est bien meilleur qu'avant. Obtenir des informations est plus pratique qu'autrefois. C'est par le cercle d'amis WeChat que beaucoup de jeunes ont l'habitude de lire les nouvelles. »
\par \textbf{Ordre des réponses :} \zh{不管}、\zh{或者}、\zh{不管}、\zh{或者}、\zh{都}、\zh{比}、\zh{得}（\zh{得多}）、\zh{比}、\zh{更}、\zh{的}. (Le dernier trou correspond au \zh{的} final de la structure \zh{是……的} : \zh{很多年轻人的习惯是通过微信朋友圈看新闻的}.)
\end{corrige}

\begin{corrige}[Exercice 8 — Remise en ordre]
\begin{enumerate}
    \item \zh{这条新闻是真的。} (Zhè tiáo xīnwén shì zhēn de.) « Cette nouvelle est vraie. » Ici \zh{是……的} renforce l'affirmation (valeur d'insistance sur l'authenticité). Sans insistance, on dirait simplement \zh{这条新闻是真。} (peu usité isolément).
    \item \zh{手机比电视方便得多。} (Shǒujī bǐ diànshì fāngbiàn de duō.) « Le téléphone est beaucoup plus pratique que la télévision. » (\zh{是} était un distracteur : pas de \zh{是} dans une phrase comparative simple.)
    \item \zh{不管是老人或者是年轻人，都喜欢看短视频。} (Bùguǎn shì lǎorén huòzhě shì niánqīngrén, dōu xǐhuan kàn duǎnshìpín.) « Que ce soit les personnes âgées ou les jeunes, tous aiment regarder les vidéos courtes. » \zh{都} se place juste après la virgule, devant le verbe.
    \item \zh{我是在家看的电视。} (Wǒ shì zài jiā kàn de diànshì.) « C'est à la maison que j'ai regardé la télévision. » Variante acceptée : \zh{我是在家看电视的。} (Wǒ shì zài jiā kàn diànshì de.)
    \item \zh{电子版比传统报纸更环保。} (Diànzǐ bǎn bǐ chuántǒng bàozhǐ gèng huánbǎo.) « La version électronique est plus écologique que le journal traditionnel. » (\zh{是} était un distracteur.)
\end{enumerate}
\end{corrige}

\begin{corrige}[Exercice 9 — Type Bac : Compréhension et questions de langue]
\textbf{Questions de compréhension (4 pts, 1 pt par réponse) :}
\begin{enumerate}
    \item Un fonctionnaire habitant à Yaoundé (\zh{住在雅温得的公务员}) et un commerçant habitant à Douala (\zh{住在杜阿拉的商人}).
    \item Il y a dix ans, on s'informait principalement par les chaînes de télévision et les journaux (\zh{电视台和报纸}).
    \item Deux raisons : (a) surfer sur Internet avec un téléphone coûte moins cher qu'acheter un téléviseur (\zh{手机上网比买电视机便宜}) ; (b) on peut regarder des directs partout et à tout moment (\zh{随时随地都能看直播}).
    \item Il l'a apprise par un groupe WeChat (\zh{是通过微信群}), hier.
\end{enumerate}
\textbf{Questions de langue (6 pts) :}
\begin{enumerate}
    \setcounter{enumi}{4}
    \item (2 pts) Occurrences de \zh{不管是……或者是……} : (a) \zh{不管是住在雅温得的公务员，或者是住在杜阿拉的商人} ; (b) \zh{不管是真消息或者是假消息}.
    \item (2 pts) Trois occurrences de \zh{是……的} : (a) \zh{是十年前，大家主要通过电视台和报纸了解新闻的} — relief sur le \textbf{temps} (il y a dix ans) ; (b) \zh{是通过微信群，我昨天才知道了那条……假消息} — relief sur le \textbf{moyen} (par le groupe WeChat) ; (c) la tournure \zh{是……的} du point (a) porte aussi sur le moyen (\zh{通过电视台和报纸}).
    \item (2 pts) Phrases comparatives : (a) \zh{短视频平台比传统电视受欢迎得多} — complément de degré \zh{得多} ; (b) \zh{手机上网比买电视机便宜} — comparaison simple sans complément de degré.
\end{enumerate}
\textbf{Points de vigilance :} ne pas confondre le \zh{的} de \zh{是……的} avec le \zh{的} possessif ; dans \zh{受欢迎得多}, \zh{得多} suit l'adjectif et ne peut pas être remplacé par \zh{很}.
\end{corrige}

\begin{corrige}[Exercice 10 — Type Bac : Thème et Version]
\textbf{A. Thème (3 pts, 1 pt par phrase) :}
\begin{enumerate}
    \item \zh{不管是看电视，或者是听收音机，我都喜欢关注新闻。} (Bùguǎn shì kàn diànshì, huòzhě shì tīng shōuyīnjī, wǒ dōu xǐhuan guānzhù xīnwén.)
    \item \zh{我是昨天晚上通过互联网知道这个消息的。} (Wǒ shì zuótiān wǎnshang tōngguò hùliánwǎng zhīdào zhè ge xiāoxi de.)
    \item \zh{互联网比报纸快得多。} (Hùliánwǎng bǐ bàozhǐ kuài de duō.)
\end{enumerate}
\textbf{B. Version (3 pts, 1 pt par phrase) :}
\begin{enumerate}
    \setcounter{enumi}{3}
    \item \zh{不管是传统媒体或者是新媒体，都有优点和缺点。} (Bùguǎn shì chuántǒng méitǐ huòzhě shì xīn méitǐ, dōu yǒu yōudiǎn hé quēdiǎn.) « Que ce soient les médias traditionnels ou les nouveaux médias, tous ont des avantages et des inconvénients. »
    \item \zh{这篇报道是记者上周在现场采访写的。} (Zhè piān bàodào shì jìzhě shàng zhōu zài xiànchǎng cǎifǎng xiě de.) « C'est la semaine dernière, en interviewant sur le terrain, que le journaliste a écrit ce reportage. » (Relief sur le temps et le lieu/la manière.)
    \item \zh{看直播比看录播有意思多了。} (Kàn zhíbō bǐ kàn lùbō yǒuyìsi duō le.) « Regarder en direct est bien plus intéressant que regarder en différé. » (\zh{多了} = variante courante de \zh{得多}.)
\end{enumerate}
\textbf{Points de vigilance :} en thème, ne jamais mettre \zh{很} dans la comparaison (phrase 3) ; ne pas garder \zh{了} après le verbe dans \zh{是……的} (phrase 2) ; en version, bien rendre la valeur d'insistance par « c'est… que » et respecter \zh{或者是} (leçon) plutôt que \zh{还是}.
\end{corrige}

\begin{corrige}[Exercice 11 — Type Bac : Expression écrite guidée]
\textbf{Proposition de rédaction modèle (environ 95 caractères) :}
\par \zh{现在的年轻人，不管是上学或者是休息，都喜欢刷社交媒体。短视频比传统电视方便得多，随时随地都能看直播、获取资讯。我是上个月通过朋友圈知道那条假新闻的。所以，不管是真消息或者是假消息，我们都要学会辨别，养成好的习惯。}
\par \emph{(Xiànzài de niánqīngrén, bùguǎn shì shàngxué huòzhě shì xiūxi, dōu xǐhuan shuā shèjiāo méitǐ. Duǎnshìpín bǐ chuántǒng diànshì fāngbiàn de duō, suíshí suídì dōu néng kàn zhíbō, huòqǔ zīxùn. Wǒ shì shàng gè yuè tōngguò péngyouquān zhīdào nà tiáo jiǎ xīnwén de. Suǒyǐ, bùguǎn shì zhēn xiāoxi huòzhě shì jiǎ xiāoxi, wǒmen dōu yào xuéhuì biànbié, yǎngchéng hǎo de xíguàn.)}
\par Traduction : « Les jeunes d'aujourd'hui, qu'ils soient en cours ou en repos, aiment tous consulter les réseaux sociaux. Les vidéos courtes sont beaucoup plus pratiques que la télévision traditionnelle : on peut regarder des directs et s'informer partout, à tout moment. C'est le mois dernier, par le cercle d'amis, que j'ai appris cette fausse nouvelle. Donc, que les informations soient vraies ou fausses, nous devons apprendre à les distinguer et prendre de bonnes habitudes. »
\par \textbf{Barème indicatif (10 pts) :} respect de la longueur et du sujet (2 pts) ; deux structures de la leçon correctement employées (4 pts, 2 pts chacune) ; au moins 5 mots de vocabulaire du module (2 pts) ; exactitude des caractères et de l'ordre des mots (2 pts).
\par \textbf{Points de vigilance :} \zh{都} obligatoire après la clause en \zh{不管是……或者是……} ; pas de \zh{很} dans la phrase comparative ; pas de \zh{了} après le verbe dans \zh{是……的} ; le sujet précède \zh{是} dans la structure d'insistance.
\end{corrige}

\newpage

\coursec{Fiche bilan}

\begin{popfigure}
\begin{center}
\begin{tikzpicture}[mindmap, grow cyclic, every node/.style={concept, align=center, font=\small},
root concept/.append style={concept color=popblue, font=\bfseries},
level 1/.append style={level distance=3.4cm, sibling angle=60},
level 2/.append style={level distance=2.2cm, font=\scriptsize}]
\node[root concept, align=center]{Leçon 35\\Grammaire\\Tle A}
child[concept color=poporange]{ node[align=center]{不管是……\\或者是……}
  child { node[align=center]{quoi que…\\ou…} }
  child { node[align=center]{都 / 也\\obligatoire} }
}
child[concept color=popgreen]{ node[align=center]{是……的\\(passé)}
  child { node[align=center]{quand ?\\où ? comment ?} }
  child { node[align=center]{是 peut\\tomber} }
}
child[concept color=poppurple]{ node {比字句}
  child { node {A 比 B + Adj} }
  child { node {更 / 还 / 得多} }
}
child[concept color=poppink]{ node[align=center]{Pièges\\francophones}
  child { node[align=center]{pas de\\« très » avec 比} }
  child { node[align=center]{ordre\\des mots} }
}
child[concept color=popgold]{ node[align=center]{Lexique\\médias}
  child { node[align=center]{报纸 电视\\网络 手机} }
}
child[concept color=popblue!60]{ node[align=center]{Méthodes\\express}
  child { node[align=center]{identifier\\le schéma} }
  child { node[align=center]{traduire\\par étapes} }
};
\end{tikzpicture}
\end{center}
\legende{Carte mentale de la leçon : trois structures, lexique, pièges et méthodes.}
\end{popfigure}

\begin{bilan}
\end{bilan}

\courssub{Lexique clé (mass médias et télécommunication)}
\begin{center}
\begin{tabularx}{\linewidth}{|X|X|X|X|}
\hline
\rowcolor{popblueL} Caractères & Pinyin & Nature & Traduction \\
\hline
报纸 & bàozhǐ & nom & journal (papier) \\
\hline
电视 & diànshì & nom & télévision \\
\hline
网络 & wǎngluò & nom & Internet, réseau \\
\hline
手机 & shǒujī & nom & téléphone portable \\
\hline
消息 & xiāoxi & nom & nouvelle, information \\
\hline
节目 & jiémù & nom & émission, programme \\
\hline
记者 & jìzhě & nom & journaliste \\
\hline
发 & fā & verbe & envoyer (message) \\
\hline
不管 & bùguǎn & conj. & peu importe, quoi que \\
\hline
或者 & huòzhě & conj. & ou (dans une phrase) \\
\hline
比 & bǐ & prép. & comparé à, que (comparaison) \\
\hline
更 & gèng & adv. & encore plus \\
\hline
\end{tabularx}
\end{center}

\courssub{Règles de grammaire à connaître par cœur}
\begin{center}
\begin{tabularx}{\linewidth}{|X|X|X|}
\hline
\rowcolor{popblueL} Structure & Exemple (caractères + pinyin) & Traduction \\
\hline
不管是 A，或者是 B，都/也 + V & 不管是报纸，或者是网络，我都看。 \emph{Bùguǎn shì bàozhǐ, huòzhě shì wǎngluò, wǒ dōu kàn.} & Que ce soit le journal ou Internet, je lis tout. \\
\hline
Sujet + 是 + (temps/lieu/manière) + V + 的 & 我是昨天在网上看到这个消息的。 \emph{Wǒ shì zuótiān zài wǎngshang kàndào zhège xiāoxi de.} & C'est hier sur Internet que j'ai vu cette nouvelle. \\
\hline
A + 比 + B + Adj & 手机比报纸快。 \emph{Shǒujī bǐ bàozhǐ kuài.} & Le portable est plus rapide que le journal. \\
\hline
A + 比 + B + Adj + 得多/一点儿 & 网络比电视快得多。 \emph{Wǎngluò bǐ diànshì kuài de duō.} & Internet est beaucoup plus rapide que la télé. \\
\hline
A + 比 + B + 更/还 + Adj & 这个节目比那个节目更有意思。 \emph{Zhège jiémù bǐ nàge jiémù gèng yǒu yìsi.} & Cette émission est encore plus intéressante que celle-là. \\
\hline
\end{tabularx}
\end{center}

\textbf{Points critiques :}
\begin{itemize}
\item \cle{不管是……或者是……} : la réponse principale contient presque toujours \zh{都} ou \zh{也} ; on ne peut pas les omettre.
\item \cle{是……的} : sert à insister sur les circonstances (temps, lieu, manière) d'une action \textbf{déjà accomplie} ; \zh{是} peut être omis, jamais \zh{的} final ; la négation est \zh{不是……的}.
\item \cle{比字句} : l'adjectif seul suffit, sans \zh{很} ; pour nuancer : \zh{得多} (beaucoup plus), \zh{一点儿} (un peu plus), \zh{更/还} (encore plus).
\end{itemize}

\courssub{Méthodes express}
\begin{enumerate}
\item Repérer le mot-clé (\zh{不管}, \zh{是…的}, \zh{比}) avant de traduire.
\item Pour \zh{是……的} : identifier d'abord l'action, puis la circonstance mise en relief, traduire par « c'est … que ».
\item Pour \zh{比} : poser A, puis B, puis l'adjectif ; vérifier qu'il n'y a pas \zh{很} avant l'adjectif.
\item Relire la phrase chinoise à voix haute pour vérifier l'ordre : sujet → circonstance → verbe.
\end{enumerate}


\begin{aretenir}[Checklist avant le Bac]
\begin{itemize}
\item[$\square$] Je connais le schéma \zh{不管是……或者是……都/也……} et je place \zh{都} ou \zh{也} dans la proposition principale.
\item[$\square$] Je sais que \zh{是……的} insiste sur le moment, le lieu ou la manière d'une action passée.
\item[$\square$] Je sais que dans \zh{是……的}, le \zh{的} final est obligatoire et la négation est \zh{不是……的}.
\item[$\square$] Je construis la comparaison avec \zh{A 比 B + adjectif} sans mettre \zh{很} devant l'adjectif.
\item[$\square$] Je sais nuancer la comparaison avec \zh{得多}, \zh{一点儿}, \zh{更} et \zh{还}.
\item[$\square$] Je connais la négation de la comparaison : \zh{A 没有 B + Adj} ou \zh{A 不比 B + Adj}.
\item[$\square$] Je maîtrise le lexique des médias : \zh{报纸}, \zh{电视}, \zh{网络}, \zh{手机}, \zh{消息}, \zh{节目}, \zh{记者}.
\item[$\square$] Je place correctement les compléments de temps et de lieu avant le verbe.
\item[$\square$] Je traduis \zh{是……的} par « c'est … que » et \zh{比} par « plus … que ».
\item[$\square$] Je vérifie les tons et le pinyin de chaque mot nouveau que j'écris.
\item[$\square$] Je sais transformer une phrase simple en phrase d'insistance avec \zh{是……的}.
\item[$\square$] Je peux rédiger un court texte de 4 à 5 phrases utilisant au moins deux des trois structures.
\end{itemize}
\end{aretenir}

\findecours

\end{document}
